% ============================================================================
% ALEXANDRIA EN LA PRÁCTICA
% El lenguaje de perfiles, alexandria-server y alexandria-client,
% con un espacio de datos de medida energética como caso de estudio
% ============================================================================
\documentclass[11pt,a4paper,openany,oneside]{book}

\usepackage{iftex}
\ifPDFTeX
  \usepackage[utf8]{inputenc}
  \usepackage[T1]{fontenc}
  \usepackage{lmodern}
\else
  \usepackage{fontspec}
\fi

\usepackage[spanish,es-tabla,es-nodecimaldot,es-noshorthands]{babel}
\usepackage[top=2.2cm,bottom=2.2cm,left=2.25cm,right=2.25cm,headheight=15pt]{geometry}
\usepackage{setspace}
\setstretch{1.07}
\usepackage{microtype}
\usepackage{xcolor}
\usepackage{graphicx}
\usepackage{booktabs}
\usepackage{tabularx}
\usepackage{longtable}
\usepackage{array}
\usepackage{enumitem}
\usepackage{amsmath,amssymb}
\usepackage{upquote}
\usepackage{fancyhdr}
\usepackage{titlesec}
\usepackage[most]{tcolorbox}
\usepackage{listings}
\usepackage{tikz}
\usetikzlibrary{arrows.meta,positioning,shapes.geometric,fit,backgrounds,calc,decorations.pathreplacing,babel}
\usepackage{xurl}
\usepackage[hypertexnames=false,colorlinks=true,linkcolor=alxNavy,citecolor=alxTeal,urlcolor=alxBlue]{hyperref}

% --- Evitar desbordamientos horizontales feos ---
\setlength{\emergencystretch}{3em}
\hfuzz=2pt
\setlength{\parskip}{0.35em}

% --- Paleta ---
\definecolor{alxNavy}{HTML}{0F2942}
\definecolor{alxBlue}{HTML}{1D5B9B}
\definecolor{alxTeal}{HTML}{0D7A6F}
\definecolor{alxAmber}{HTML}{B45309}
\definecolor{alxRed}{HTML}{991B1B}
\definecolor{alxPurple}{HTML}{5B21B6}
\definecolor{alxSlate}{HTML}{F4F6F9}
\definecolor{alxLightBlue}{HTML}{EBF3FA}
\definecolor{alxLightTeal}{HTML}{EAF6F4}
\definecolor{alxLightAmber}{HTML}{FEF6E7}
\definecolor{alxLightRed}{HTML}{FDF2F2}
\definecolor{alxLightPurple}{HTML}{F3EFFB}
\definecolor{alxBorder}{HTML}{CBD5E1}
\definecolor{codeComment}{HTML}{57606A}
\definecolor{codeString}{HTML}{0A3069}
\definecolor{codeKeyword}{HTML}{8250DF}

% --- Columnas útiles ---
\newcolumntype{L}[1]{>{\raggedright\arraybackslash}p{#1}}
\newcolumntype{C}[1]{>{\centering\arraybackslash}p{#1}}
\newcolumntype{Y}{>{\raggedright\arraybackslash}X}

% --- Código en línea: se parte en cualquier carácter y no interpreta guiones bajos ---
\Urlmuskip=0mu\relax
\newcommand{\cod}[1]{\nolinkurl{#1}}
\newcommand{\ok}{\textcolor{alxTeal}{\textbf{ok}}}
\newcommand{\warn}{\textcolor{alxAmber}{\textbf{aviso}}}
\newcommand{\viol}{\textcolor{alxRed}{\textbf{violación}}}
\newcommand{\si}{\textcolor{alxTeal}{\textbf{sí}}}
\newcommand{\no}{\textcolor{alxRed}{\textbf{no}}}
\newcommand{\pasa}{\textcolor{alxTeal}{pasa}}
\newcommand{\falla}{\textcolor{alxRed}{falla}}

% --- Encabezados y pies ---
\pagestyle{fancy}
\fancyhf{}
\fancyhead[L]{\small\textcolor{alxNavy!80}{\nouppercase{\leftmark}}}
\fancyhead[R]{\small\textcolor{alxBlue}{\textbf{\thepage}}}
\fancyfoot[C]{\footnotesize\textcolor{gray}{Alexandria en la práctica $\cdot$ lenguaje de perfiles, hub y cliente}}
\renewcommand{\headrulewidth}{0.4pt}
\renewcommand{\footrulewidth}{0.2pt}
\fancypagestyle{plain}{%
  \fancyhf{}%
  \fancyfoot[C]{\small\textcolor{alxBlue}{\textbf{\thepage}}}%
  \renewcommand{\headrulewidth}{0pt}%
  \renewcommand{\footrulewidth}{0pt}%
}

% --- Partes, capítulos y secciones ---
\titleformat{\part}[display]
  {\centering\Huge\bfseries\color{alxNavy}}
  {\normalfont\Large\scshape\color{alxTeal} Parte \thepart}
  {1ex}
  {\titlerule[1.5pt]\vspace{1.5ex}}
  [\vspace{1.5ex}\titlerule]
\titleformat{\chapter}[display]
  {\normalfont\bfseries\color{alxNavy}}
  {\filright\normalfont\large\scshape\color{alxTeal} \chaptertitlename\ \thechapter}
  {1ex}
  {\titlerule[1pt]\vspace{1ex}\huge\raggedright}
  [\vspace{1ex}\titlerule]
\titlespacing*{\chapter}{0pt}{-15pt}{25pt}
\titleformat{\section}
  {\normalfont\Large\bfseries\color{alxNavy}}
  {\thesection}{0.8em}{}
\titleformat{\subsection}
  {\normalfont\large\bfseries\color{alxBlue}}
  {\thesubsection}{0.7em}{}
\titleformat{\subsubsection}
  {\normalfont\normalsize\bfseries\color{alxTeal}}
  {\thesubsubsection}{0.6em}{}

% --- Índice: hueco suficiente para números de capítulo de dos cifras ---
\usepackage{etoolbox}
\makeatletter
\patchcmd{\l@chapter}{1.5em}{2.1em}{}{}
\renewcommand*\l@section{\@dottedtocline{1}{2.1em}{2.7em}}
\makeatother

% --- Cajas pedagógicas ---
\newtcolorbox{ideaClave}[1][Idea clave]{
  enhanced, breakable,
  colback=alxLightBlue, colframe=alxBlue,
  coltitle=white, fonttitle=\bfseries,
  title={#1},
  boxrule=0.8pt, arc=3pt, left=8pt, right=8pt, top=6pt, bottom=6pt
}
\newtcolorbox{notaPractica}[1][Nota práctica]{
  enhanced, breakable,
  colback=alxLightTeal, colframe=alxTeal,
  coltitle=white, fonttitle=\bfseries,
  title={#1},
  boxrule=0.8pt, arc=3pt, left=8pt, right=8pt, top=6pt, bottom=6pt
}
\newtcolorbox{trampaHabitual}[1][Trampa habitual]{
  enhanced, breakable,
  colback=alxLightAmber, colframe=alxAmber,
  coltitle=white, fonttitle=\bfseries,
  title={#1},
  boxrule=0.8pt, arc=3pt, left=8pt, right=8pt, top=6pt, bottom=6pt
}
\newtcolorbox{bajoElCapo}[1][Bajo el capó]{
  enhanced, breakable,
  colback=alxSlate, colframe=alxNavy,
  coltitle=white, fonttitle=\bfseries,
  title={#1},
  boxrule=0.8pt, arc=3pt, left=8pt, right=8pt, top=6pt, bottom=6pt
}
\newtcolorbox{hallazgo}[1]{
  enhanced, breakable,
  colback=alxLightRed, colframe=alxRed,
  coltitle=white, fonttitle=\bfseries,
  title={#1},
  boxrule=0.8pt, arc=3pt, left=8pt, right=8pt, top=6pt, bottom=6pt
}
\newtcolorbox{retos}[1][Para practicar]{
  enhanced, breakable,
  colback=white, colframe=alxPurple,
  coltitle=white, fonttitle=\bfseries,
  title={#1},
  boxrule=0.9pt, arc=3pt, left=8pt, right=8pt, top=6pt, bottom=6pt
}

% --- Listados ---
\lstset{
  basicstyle=\footnotesize\ttfamily,
  backgroundcolor=\color{alxSlate},
  frame=single,
  rulecolor=\color{alxBorder},
  framesep=5pt,
  xleftmargin=6pt,
  xrightmargin=6pt,
  breaklines=true,
  breakatwhitespace=false,
  columns=fullflexible,
  keepspaces=true,
  showstringspaces=false,
  upquote=true,
  captionpos=b,
  aboveskip=0.9em,
  belowskip=0.7em,
  commentstyle=\color{codeComment}\itshape,
  keywordstyle=\color{codeKeyword}\bfseries,
  stringstyle=\color{codeString},
  literate=
    {á}{{\'a}}1 {é}{{\'e}}1 {í}{{\'\i}}1 {ó}{{\'o}}1 {ú}{{\'u}}1
    {Á}{{\'A}}1 {É}{{\'E}}1 {Í}{{\'I}}1 {Ó}{{\'O}}1 {Ú}{{\'U}}1
    {ñ}{{\~n}}1 {Ñ}{{\~N}}1 {ü}{{\"u}}1 {Ü}{{\"U}}1
    {¿}{{?`}}1 {¡}{{!`}}1
    {→}{{$\rightarrow$}}1 {…}{{\ldots}}1 {·}{{$\cdot$}}1 {§}{{\S}}1
}
\lstdefinelanguage{yamlx}{
  morecomment=[l]{\#},
  morestring=[b]",
  sensitive=true
}
\lstdefinelanguage{jsonx}{
  morestring=[b]",
  morekeywords={true,false,null},
  sensitive=true
}
\lstdefinelanguage{turtlex}{
  morekeywords={@prefix,a,true,false},
  morestring=[b]",
  sensitive=true
}
\lstdefinelanguage{shx}{
  morekeywords={curl,cd,source,echo,cat,export,make,go,grep,sed,kill,sleep,python3,docker},
  morecomment=[l]{\#},
  morestring=[b]",
  morestring=[b]',
  sensitive=true
}
\lstdefinestyle{yaml}{language=yamlx}
\lstdefinestyle{json}{language=jsonx}
\lstdefinestyle{ttl}{language=turtlex}
\lstdefinestyle{sh}{language=shx, backgroundcolor=\color{alxLightBlue!55}, rulecolor=\color{alxBlue!35}}
\lstdefinestyle{out}{language={}, backgroundcolor=\color{white}, rulecolor=\color{alxBorder}, basicstyle=\scriptsize\ttfamily}

% --- Metadatos ---
\hypersetup{
  pdftitle={Alexandria en la práctica},
  pdfsubject={Tutorial de alexandriac, alexandria-server y alexandria-client},
  pdfkeywords={Alexandria, vocabulary hub, ODRL, DSP, SHACL, espacios de datos}
}

\begin{document}

% ============================== PORTADA ==============================
\begin{titlepage}
\begin{tikzpicture}[remember picture,overlay]
  \fill[alxNavy] (current page.north west) rectangle ([yshift=-9.4cm]current page.north east);
  \fill[alxTeal] ([yshift=-9.4cm]current page.north west) rectangle ([yshift=-9.75cm]current page.north east);
\end{tikzpicture}
\vspace*{0.6cm}
\begin{flushleft}
{\color{white}\large\scshape Tutorial de uso\par}
\vspace{0.7cm}
{\color{white}\fontsize{34}{40}\selectfont\bfseries Alexandria en la práctica\par}
\vspace{0.5cm}
{\color{white!85}\Large El lenguaje de perfiles \texttt{alexandria.profiles/v1},\\[2pt]
el hub \texttt{alexandria-server} y el borde \texttt{alexandria-client}\par}
\end{flushleft}
\vspace{3.4cm}
\begin{flushleft}
{\Large\color{alxNavy}\bfseries Caso de estudio: un espacio de datos de medida energética\par}
\vspace{0.4cm}
{\large\color{alxNavy!80} De la regla de negocio al veredicto en la red: perfil ODRL, publicación gobernada,
validación de tráfico DSP y una revisión de la parte de SHACL que queda por hacer.\par}
\end{flushleft}
\vfill
\begin{tcolorbox}[enhanced, colback=alxSlate, colframe=alxBorder, boxrule=0.6pt, arc=3pt, left=10pt, right=10pt, top=8pt, bottom=8pt]
\small
\textbf{Qué está comprobado.} Todos los comandos, salidas, cabeceras y códigos de error de este documento se han
ejecutado contra los binarios compilados desde la rama \cod{feat/client-validator-mode} (commit \cod{e824d3a})
con las correcciones del capítulo~\ref{ch:shacl} aplicadas en el árbol de trabajo, el 13 y el 14 de noviembre de
2026. Lo que no se ha ejecutado se marca expresamente como tal.

\medskip
\textbf{Ficheros.} Los ejemplos están en \cod{tutorial/manual/examples/}; el laboratorio usa puertos alternativos
para convivir con el de \cod{config/}.
\end{tcolorbox}
\vspace{0.4cm}
{\small\color{gray} Repositorio \cod{ging/alexandria} $\cdot$ especificación \cod{alexandria-vhub} 0.3.0 $\cdot$ noviembre de 2026}
\end{titlepage}

\frontmatter
\setcounter{tocdepth}{1}
\tableofcontents

\mainmatter

\chapter*{Cómo leer este tutorial}
\addcontentsline{toc}{chapter}{Cómo leer este tutorial}
\markboth{Cómo leer este tutorial}{}

Este documento recorre Alexandria de punta a punta con un único ejemplo que se parece a un problema de
verdad: una distribuidora eléctrica que comparte curvas de carga cuartohorarias con comercializadoras,
agregadores de flexibilidad y un grupo de investigación. No hay colores ni acciones de juguete: hay finalidades
de uso, agregación mínima, plazos de retención y jurisdicciones, que es de lo que tratan las políticas reales.

El recorrido tiene cinco partes y cada una se apoya en la anterior:

\begin{enumerate}[leftmargin=*, itemsep=2pt]
  \item \textbf{El modelo mental.} Qué problema resuelve Alexandria, qué hace cada binario y cuál es el caso de estudio.
  \item \textbf{El lenguaje de perfiles.} Cómo se escribe un \cod{OdrlProfile} y un \cod{SpaceProfile}, qué genera
        \cod{alexandriac}, cómo valida y cómo se versiona.
  \item \textbf{El hub.} Cómo \cod{alexandria-server} convierte ese YAML en una versión firmada, con fechas y auditoría.
  \item \textbf{El borde.} Cómo \cod{alexandria-client} sincroniza esas versiones y decide sobre cada mensaje DSP.
  \item \textbf{Revisión de SHACL.} Qué está hecho, qué no, qué falla cuando se ejecutan de verdad las \emph{shapes}
        generadas y qué haría falta para cerrarlo.
\end{enumerate}

\section*{Convenciones}

\noindent\begin{tabularx}{\linewidth}{@{}L{3.6cm}Y@{}}
\toprule
\textbf{Aspecto} & \textbf{Significado} \\
\midrule
Fondo azul claro & Comando que se escribe en la terminal. \\
Fondo blanco, letra pequeña & Salida real del comando, copiada tal cual (recortada con \texttt{…} cuando es larga). \\
Fondo gris & Contenido de un fichero: YAML, JSON o Turtle. \\
Caja azul, «Idea clave» & El concepto que conviene retener de la sección. \\
Caja ámbar, «Trampa habitual» & Un comportamiento que sorprende la primera vez. \\
Caja gris, «Bajo el capó» & Dónde está en el código lo que se acaba de ver. \\
Caja roja, «Hallazgo» & Un defecto o una carencia detectada durante la revisión. \\
\bottomrule
\end{tabularx}

\section*{Qué necesitas}

\begin{itemize}[leftmargin=*, itemsep=2pt]
  \item Go 1.22 o posterior y \cod{make} para compilar (\texttt{make build}).
  \item \cod{curl} y \cod{python3} (solo biblioteca estándar) para las llamadas y para leer JSON.
  \item Para reproducir la revisión de SHACL del capítulo~\ref{ch:shacl}: un entorno virtual de Python con
        \cod{pyshacl} y \cod{jsonschema}.
\end{itemize}

\section*{Los ficheros del tutorial}

Todo lo que se cita vive en \cod{tutorial/manual/examples/}:

\noindent\begin{tabularx}{\linewidth}{@{}L{5.4cm}Y@{}}
\toprule
\textbf{Ruta} & \textbf{Contenido} \\
\midrule
\cod{dsl/metering-1.0.0.yaml} & El perfil ODRL de medida, solo. \\
\cod{dsl/flexgrid-1.0.0.yaml} & Perfil 1.0.0 más el \cod{SpaceProfile} 1.0.0: el primer borrador que se publica. \\
\cod{dsl/flexgrid-1.1.0.yaml} & Versión compatible: una finalidad nueva y una acción obsoleta. \\
\cod{dsl/flexgrid-2.0.0.yaml} & Versión de ruptura: la acción obsoleta se retira. \\
\cod{dsl/health-0.3.0.yaml} & Segundo perfil, de uso secundario de datos de salud. \\
\cod{vendor/dsp/2025-1/} & Copia local reducida de DSP 2025-1 para el \emph{binding} vendorizado. \\
\cod{msgs/} & Los mensajes DSP que se envían: \cod{m1}–\cod{m6}, el acuerdo \cod{a1} y las transferencias. \\
\cod{rules/} & Las tres reglas del cliente que se adaptan. \\
\cod{lab/} & \cod{server.yaml}, \cod{client.yaml} y \cod{env.sh}, el guion del laboratorio. \\
\cod{shacl/harness.py} & El banco de pruebas de la revisión de SHACL. \\
\bottomrule
\end{tabularx}

\begin{notaPractica}[Sobre los puertos]
El laboratorio de este tutorial escucha en \cod{18443} (hub), \cod{19400} (IdP), \cod{18282} (conector),
\cod{19080}/\cod{19081} (tráfico) y \cod{19090} (administración). Son distintos de los de \cod{config/server.yaml}
y \cod{config/client.yaml} a propósito: así puede convivir con el laboratorio por defecto, que usa otra base de
datos, otro compilador de borradores y otro \emph{namespace}.
\end{notaPractica}


\part{El modelo mental}
\chapter{Qué resuelve Alexandria y con qué piezas}
\label{ch:modelo}

\section{El problema: una pila de especificaciones con huecos}

Un espacio de datos no es una especificación: es una \textbf{pila} de ellas. ODRL para las políticas, DCAT para el
catálogo, DSP para el protocolo, DCP para las credenciales. Cada una deja libertad a propósito. ODRL define qué es
una restricción, pero no qué operandos existen en tu sector. DSP define un \cod{ContractRequestMessage}, pero no
qué política es aceptable dentro.

Cada una de esas libertades es un \textbf{slot}. Un espacio de datos concreto lo cierra con un \textbf{binding}:
qué perfil manda en ese slot, en qué versión, con qué obligación y en qué punto del tráfico se comprueba, que es
el \textbf{gate}.

\begin{ideaClave}
La autoridad del espacio escribe esa pila \emph{una sola vez}, en YAML. Todo lo demás (ontología, contexto JSON-LD,
JSON Schema, \emph{shapes} SHACL, especificación legible) se \textbf{genera}. Nadie edita a mano dos artefactos
que deban decir lo mismo, porque acaban diciendo cosas distintas.
\end{ideaClave}

El problema no termina al generar. Los participantes tienen que enterarse de que hay una versión nueva, fiarse de
que viene de quien dice venir, aplicarla todos a la vez y no romper los contratos ya firmados cuando una regla
cambia. De eso se ocupan el hub y el cliente.

\section{Las tres piezas}

\begin{center}
\begin{tikzpicture}[
  node distance=9mm and 19mm, font=\small,
  box/.style={draw=alxNavy, fill=alxLightBlue, rounded corners=2pt, minimum height=11mm, minimum width=30mm, align=center},
  ext/.style={draw=alxBorder, fill=alxSlate, rounded corners=2pt, minimum height=11mm, minimum width=24mm, align=center},
  art/.style={draw=alxTeal, fill=alxLightTeal, rounded corners=2pt, minimum height=9mm, minimum width=30mm, align=center},
  arr/.style={-{Stealth[length=2.2mm]}, alxNavy, thick}]
  \node[ext] (yaml) {Bundle YAML\\[-1pt]\scriptsize perfil + espacio};
  \node[box, right=of yaml] (hub) {\texttt{alexandria-server}\\[-1pt]\scriptsize hub};
  \node[art, right=of hub] (pub) {Paquete firmado\\[-1pt]\scriptsize \texttt{/\allowbreak{}pub}, listado, registro};
  \node[box, below=14mm of pub] (cli) {\texttt{alexandria-client}\\[-1pt]\scriptsize borde L7};
  \node[ext, left=of cli] (conn) {Conector DSP\\[-1pt]\scriptsize sin modificar};
  \node[ext, right=of cli] (peer) {Contraparte};
  \node[box, below=14mm of yaml] (cc) {\texttt{alexandriac}\\[-1pt]\scriptsize compilador offline};
  \draw[arr] (yaml) -- node[above]{\scriptsize borrador} (hub);
  \draw[arr] (hub) -- node[above]{\scriptsize aprueba} (pub);
  \draw[arr] (pub) -- node[right]{\scriptsize sincroniza y verifica} (cli);
  \draw[arr] (peer) -- node[above]{\scriptsize mensaje} (cli);
  \draw[arr] (cli) -- node[above]{\scriptsize reenvía} (conn);
  \draw[arr, dashed] (yaml) -- node[left]{\scriptsize mismo código} (cc);
\end{tikzpicture}
\end{center}

\noindent\begin{tabularx}{\linewidth}{@{}L{3.5cm}L{2.7cm}Y@{}}
\toprule
\textbf{Binario} & \textbf{Quién lo usa} & \textbf{Qué hace} \\
\midrule
\cod{alexandriac} & Quien escribe el perfil & Valida y compila YAML en local, sin red y sin estado. Es el bucle rápido de edición. \\
\cod{alexandria-server} & La autoridad del espacio & Gobierna el ciclo borrador, candidato, versión y revocación. Firma lo que publica y lo sirve inmutable bajo \cod{/pub}. \\
\cod{alexandria-client} & Cada participante & Se coloca delante del conector. Sincroniza las versiones, verifica las firmas y decide sobre cada mensaje con estado local. \\
\bottomrule
\end{tabularx}

\medskip
El hub usa el mismo paquete Go que \cod{alexandriac} (\cod{internal/dsl}). Lo que compila en tu portátil compila
igual en el hub.

Hay además dos binarios de laboratorio que no son producto: \cod{fake-idp}, un emisor OAuth/OIDC que firma
tokens, y \cod{fake-connector}, un conector DSP de mentira que responde \cod{200} a todo y recuerda lo que recibe.

\section{El caso de estudio: Flexgrid}

Flexgrid es un espacio de datos sectorial ficticio. Los participantes y sus intereses son estos:

\noindent\begin{tabularx}{\linewidth}{@{}L{3.3cm}Y@{}}
\toprule
\textbf{Participante} & \textbf{Qué quiere} \\
\midrule
Distribuidora (DSO) & Publica la curva de carga cuartohoraria de sus puntos de suministro. Es la proveedora y quiere limitar para qué se usa. \\
Agregador & Entrena modelos de previsión de demanda sobre datos agregados por zona. \\
Comercializadora & Liquida desvíos y factura. Hasta ahora también perfilaba consumidores individuales. \\
Universidad & Quiere los datos para investigación, una finalidad que al principio no existe. \\
Autoridad del espacio & Decide el vocabulario común y lo publica. Lucía, de la distribuidora, actúa como \emph{steward}. \\
\bottomrule
\end{tabularx}

\medskip
Las reglas de negocio que la autoridad quiere poder expresar, y con qué pieza de ODRL se corresponden:

\noindent\begin{tabularx}{\linewidth}{@{}Y L{2.6cm} L{4.2cm}@{}}
\toprule
\textbf{Regla de negocio} & \textbf{Pieza ODRL} & \textbf{Término} \\
\midrule
«Puedes usarlos para prever la demanda.» & Acción & \cod{enr:forecastDemand} \\
«Puedes usarlos para liquidar desvíos.» & Acción & \cod{enr:settleImbalance} \\
«Solo para esta finalidad.» & Operando izquierdo & \cod{odrl:purpose}, adoptado y acotado \\
«Cada registro agrega al menos 50 puntos de suministro.» & Operando izquierdo & \cod{enr:minAggregation} \\
«No los guardes más de 90 días.» & Operando izquierdo & \cod{enr:retentionDays} \\
«Solo se tratan en España o en la UE.» & Operando izquierdo & \cod{enr:processingRegion} \\
«El permiso caduca el 1 de enero.» & Operando izquierdo & \cod{odrl:dateTime}, adoptado \\
\bottomrule
\end{tabularx}

\medskip
Sobre este espacio se cuenta además una historia de gobierno, que es la que da sentido al ciclo de vida:

\begin{enumerate}[leftmargin=*, itemsep=2pt]
  \item \textbf{1.0.0.} Primera versión. Incluye la acción \cod{enr:profileConsumer}, perfilar a un consumidor individual.
  \item \textbf{1.1.0.} La universidad entra: se añade la finalidad \cod{enr:research}. A la vez el regulador avisa de
        que el perfilado individual no va a ser admisible, y la acción se marca obsoleta. Cambio \emph{compatible}.
  \item \textbf{2.0.0.} El perfilado se retira. Cambio de \emph{ruptura}, con una ventana de gracia para migrar.
  \item \textbf{Revocación.} El regulador declara ilícito el perfilado. La autoridad revoca la 1.1.0, y los acuerdos
        firmados bajo ella dejan de valer para nuevas transferencias.
\end{enumerate}

\section{Preparación}

Desde la raíz del repositorio:

\begin{lstlisting}[style=sh]
make build
\end{lstlisting}
\begin{lstlisting}[style=out]
go build -o bin/alexandriac ./cmd/alexandriac
go build -o bin/alexandria-server ./cmd/alexandria-server
go build -o bin/alexandria-client ./cmd/alexandria-client
go build -o bin/fake-idp ./deploy/fake-idp
go build -o bin/fake-connector ./deploy/fake-connector
\end{lstlisting}

Los capítulos de la parte II solo necesitan \cod{bin/alexandriac}. Los de las partes III y IV usan los cinco.


\part{El lenguaje de perfiles}
\chapter{Anatomía de un \texttt{OdrlProfile}}
\label{ch:odrlprofile}

Un \cod{OdrlProfile} es un documento YAML que dice tres cosas: qué términos de ODRL \textbf{adopta} el espacio,
qué términos nuevos \textbf{crea} y cómo de \textbf{cerrado} es el vocabulario resultante. Este capítulo recorre
\cod{examples/dsl/metering-1.0.0.yaml} sección por sección.

\section{Cabecera e identidad}

\begin{lstlisting}[style=yaml]
dsl: alexandria.profiles/v1
kind: OdrlProfile

profile:
  id: https://flexgrid.example/odrl/profile/metering
  version: 1.0.0
  title:
    en: Smart Metering Data Usage Profile
    es: Perfil de uso de datos de medida inteligente
  description:
    en: ODRL profile for sharing quarter-hourly smart meter data between DSOs, retailers and aggregators.
  publisher: Flexgrid Data Space Authority
  license: https://creativecommons.org/licenses/by/4.0/
  base_url: https://flexgrid.example/odrl/profile/metering
\end{lstlisting}

\noindent\begin{tabularx}{\linewidth}{@{}L{3.7cm}L{2.2cm}Y@{}}
\toprule
\textbf{Campo} & \textbf{Obligatorio} & \textbf{Qué es} \\
\midrule
\cod{dsl} & sí & Siempre \cod{alexandria.profiles/v1}. \\
\cod{kind} & sí & \cod{OdrlProfile} o \cod{SpaceProfile}. \\
\cod{profile.id} & sí & IRI del perfil. Es lo que una política declara en su campo \cod{profile}. \\
\cod{profile.version} & sí & SemVer completo: \cod{MAYOR.MENOR.PARCHE}. \\
\cod{profile.title} & sí & Mapa de idiomas. La clave \cod{en} es obligatoria en todos los mapas de idiomas. \\
\cod{profile.description} & no & Mapa de idiomas. \\
\cod{profile.publisher} & sí & Texto libre. \\
\cod{profile.license} & sí & IRI de la licencia del vocabulario. \\
\cod{profile.base_url} & sí & De aquí salen las URL de los artefactos: el contexto es \texttt{<base\_url>/\allowbreak{}context.jsonld}. \\
\bottomrule
\end{tabularx}

\begin{trampaHabitual}[El identificador del perfil no es el espacio de nombres]
El \cod{id} del perfil y el \emph{namespace} de sus términos son IRI distintas. Si coinciden, o si una es la otra
más \texttt{\#} o \cod{/}, el compilador lo rechaza con \cod{PROFILE_ID_IS_NAMESPACE}. El perfil es un documento; los
términos viven en otro sitio.
\end{trampaHabitual}

\section{Espacios de nombres}

\begin{lstlisting}[style=yaml, firstnumber=16]
namespaces:
  enr: https://flexgrid.example/ns/metering#
\end{lstlisting}

Hay que declarar al menos uno, y su IRI tiene que acabar en \texttt{\#} o en \cod{/}. Cuando hay varios, el
\textbf{primero que aparece en el fichero} es el espacio de nombres de los términos que crea el perfil. Los demás
solo sirven para citar términos ajenos.

Nueve prefijos están reservados y no se pueden redeclarar: \cod{odrl}, \cod{xsd}, \cod{rdf}, \cod{rdfs}, \cod{owl},
\cod{skos}, \cod{dct}, \cod{sig} y \cod{alx}. Están siempre disponibles.

\section{Términos adoptados}

\begin{lstlisting}[style=yaml, firstnumber=19]
adopted:
  actions:
    - odrl:aggregate
    - odrl:anonymize
  left_operands:
    - term: odrl:purpose
      signature:
        value_kind: iri
        values: [enr:gridBalancing, enr:billing, enr:demandForecast]
        operators: [eq, isAnyOf]
    - term: odrl:dateTime
      signature:
        datatype: xsd:dateTime
        operators: [lt, lteq, gt, gteq]
\end{lstlisting}

Adoptar es decir «este término ya existe y aquí se usa». Las categorías son seis: \cod{actions},
\cod{left_operands}, \cod{operators}, \cod{right_operands}, \cod{party_functions} y \cod{asset_relations}.

Cada entrada admite dos formas. La \textbf{corta} es el CURIE a secas, como \cod{odrl:aggregate}. La
\textbf{larga} es un objeto con \cod{term} y, solo para operandos izquierdos, \cod{signature} y \cod{enforcement}.

Lo interesante está en la forma larga. ODRL define \cod{odrl:purpose}, pero no dice qué finalidades existen. Al
adoptarlo con una firma, el perfil lo acota: en Flexgrid una finalidad es una de tres IRI, y se compara con
\cod{eq} o con \cod{isAnyOf}. El término sigue siendo de ODRL; la restricción es del espacio.

\begin{bajoElCapo}
Los términos \cod{odrl:} adoptados se contrastan con una copia local del vocabulario ODRL~2.2
(\cod{internal/dsl/odrlvocab.go}), nunca con la red. Se comprueba que el término exista y que su categoría sea la
declarada. Un término adoptado que no sea de ODRL tiene que estar creado por otro perfil del mismo \emph{bundle}.
\end{bajoElCapo}

\section{Términos creados}

\subsection{Acciones}

\begin{lstlisting}[style=yaml, firstnumber=34]
terms:
  - id: forecastDemand
    kind: action
    included_in: odrl:use
    label: { en: Forecast demand, es: Predecir la demanda }
    definition:
      en: To train or run load-forecasting models on the metering data.

  - id: settleImbalance
    kind: action
    included_in: odrl:use
    label: { en: Settle imbalance, es: Liquidar desvíos }
    definition:
      en: To use the metering data in the imbalance settlement between market parties.

  - id: profileConsumer
    kind: action
    included_in: odrl:use
    label: { en: Profile consumer, es: Perfilar al consumidor }
    definition:
      en: To derive a behavioural profile of an individual supply point.
\end{lstlisting}

Todo término tiene \cod{id}, \cod{kind}, \cod{label} y \cod{definition}. El \cod{id} es un nombre local, letras,
dígitos y guion bajo, que se une al espacio de nombres: \cod{forecastDemand} pasa a ser
\cod{enr:forecastDemand}.

\cod{included_in} solo existe en las acciones y dice de qué acción más general es un caso. Si falta, el
compilador avisa con \cod{ACTION_INCLUDED_IN_MISSING}, porque un motor de políticas no sabría que permitir
\cod{odrl:use} cubre también la acción nueva.

\subsection{Operandos izquierdos y sus firmas}

\begin{lstlisting}[style=yaml, firstnumber=57]
  - id: minAggregation
    kind: left_operand
    label: { en: Minimum aggregation, es: Agregación mínima }
    definition:
      en: Minimum number of supply points that every delivered record must aggregate.
    signature:
      datatype: xsd:integer
      operators: [gteq, gt]

  - id: retentionDays
    kind: left_operand
    label: { en: Retention in days, es: Retención en días }
    definition:
      en: Maximum number of days the consumer may keep the data.
    signature:
      datatype: xsd:integer
      operators: [lteq, lt]

  - id: processingRegion
    kind: left_operand
    label: { en: Processing region, es: Región de tratamiento }
    definition:
      en: Jurisdiction where the data may be stored and processed.
    signature:
      value_kind: iri
      values: [enr:regionES, enr:regionEU, enr:regionEEA]
      operators: [eq, isAnyOf]
\end{lstlisting}

La \textbf{firma} es la parte más valiosa del lenguaje. Un operando izquierdo sin firma es solo un nombre: se
puede escribir \texttt{minAggregation eq "muchos"} y nadie protesta. Con firma, el operando tiene un álgebra.

\noindent\begin{tabularx}{\linewidth}{@{}L{3.1cm}Y@{}}
\toprule
\textbf{Campo de la firma} & \textbf{Significado} \\
\midrule
\cod{operators} & Obligatorio, al menos uno. Operadores admitidos: los doce de ODRL o el CURIE de un operador creado por el perfil. \\
\cod{datatype} & Tipo XSD del operando derecho, como \cod{xsd:integer}. Implica un valor literal. \\
\cod{value_kind} & \cod{literal} o \cod{iri}. Normalmente se deduce. \\
\cod{values} & Enumeración cerrada de valores admitidos. \\
\cod{scheme} & Alternativa a \cod{values}: los valores son conceptos de un esquema SKOS. Incompatible con \cod{values}, y hoy no es practicable (hallazgo H7 del capítulo~\ref{ch:shacl}). \\
\cod{allow_reference} & Si es \cod{true}, admite \cod{rightOperandReference}. Por defecto es \cod{false}. \\
\bottomrule
\end{tabularx}

\medskip
El tipo del valor, cuando no se escribe \cod{value_kind}, se deduce en este orden: si hay \cod{datatype}, es
literal; si todos los \cod{values} parecen CURIE, es IRI; si hay \cod{scheme}, es IRI; en otro caso, literal.

No todos los operadores casan con todos los tipos. El compilador aplica esta tabla:

\noindent\begin{tabularx}{\linewidth}{@{}L{4.3cm}Y@{}}
\toprule
\textbf{Operadores} & \textbf{Exigen} \\
\midrule
\cod{eq}, \cod{neq} & Nada: valen con cualquier tipo. \\
\cod{lt}, \cod{lteq}, \cod{gt}, \cod{gteq} & Un tipo ordenado: \cod{xsd:integer}, \cod{decimal}, \cod{double}, \cod{float}, \cod{date}, \cod{dateTime} o \cod{time}. Con \cod{xsd:duration} hay aviso, porque el orden es parcial. \\
\cod{isA}, \cod{hasPart}, \cod{isPartOf} & Valores IRI. \\
\cod{isAnyOf}, \cod{isAllOf}, \cod{isNoneOf} & Nada, pero convierten el operando derecho en multivaluado: se admite una lista. \\
\bottomrule
\end{tabularx}

\subsection{Operandos derechos}

\begin{lstlisting}[style=yaml, firstnumber=85]
  - id: gridBalancing
    kind: right_operand
    label: { en: Grid balancing }
    definition: { en: Operating the distribution grid and its flexibility markets. }
  - id: billing
    kind: right_operand
    label: { en: Billing }
    definition: { en: Invoicing the energy supplied to a supply point. }
  # ... demandForecast, regionES, regionEU, regionEEA
\end{lstlisting}

Los valores IRI de una firma deberían ser operandos derechos del propio perfil. Si un valor no lo es, el
compilador avisa con \cod{SIG_VALUE_KIND}: la política sería válida, pero nadie habría definido qué significa.

\subsection{Los once tipos de término}

\noindent\begin{tabularx}{\linewidth}{@{}L{4.4cm}Y@{}}
\toprule
\textbf{\texttt{kind}} & \textbf{Forma RDF que genera} \\
\midrule
\cod{action} & \cod{odrl:Action}, con \cod{odrl:includedIn}. \\
\cod{left_operand}, \cod{operator}, \cod{right_operand}, \cod{conflict_term} & Individuo de la clase ODRL correspondiente. \\
\cod{party_function}, \cod{asset_relation}, \cod{logical_operand} & Subpropiedad de \cod{odrl:function}, \cod{odrl:relation} u \cod{odrl:operand}. \\
\cod{rule_subclass}, \cod{policy_subclass}, \cod{constraint_subclass} & Clase OWL. Exigen \cod{subclass_of}. \\
\bottomrule
\end{tabularx}

\section{Ciclo de vida de un término}

Un término tiene \cod{status}: \cod{active} por defecto, \cod{deprecated} o \cod{retired}. Junto a él,
\cod{replaced_by} dice qué usar en su lugar. En la versión 1.0.0 todo está activo; el capítulo~\ref{ch:evolucion}
usa estos campos para retirar \cod{profileConsumer} sin romper a nadie de golpe.

\begin{ideaClave}
Un término \textbf{nunca se borra}. Se marca obsoleto en una versión y retirado en otra posterior. El compilador
lo exige comparando cada versión con la anterior.
\end{ideaClave}

\section{Cierre del vocabulario}

\begin{lstlisting}[style=yaml, firstnumber=112]
rules:
  actions: closed
  left_operands: closed
\end{lstlisting}

Con la lista de acciones \textbf{cerrada}, una política solo puede usar las acciones del perfil, creadas o
adoptadas, más \cod{odrl:use} y \cod{odrl:transfer}, que son el núcleo de ODRL y un perfil no puede prohibir. Con
la lista de operandos izquierdos cerrada, una restricción solo puede usar los del perfil.

Ambas listas son cerradas por defecto. Abrir una (\cod{open}) es un cambio menor; cerrarla después es una
ruptura.

\section{El bloque \texttt{enforcement}}

Un operando izquierdo admite además un bloque \cod{enforcement} con \cod{cedar_path}, \cod{pdp_level} y
\cod{eval_timing}, que describe cómo lo evaluaría un motor de decisión. El lenguaje lo valida: exige que el
operando tenga firma, y falla con \cod{ENFORCEMENT_WITHOUT_SIGNATURE} si no la tiene.

\begin{notaPractica}[Hoy no genera nada]
La proyección a Cedar está desconectada (\texttt{cedarEnabled = false} en \cod{internal/dsl/cedar.go}). El bloque se
analiza y se versiona, pero no produce \cod{cedar-mapping.json} ni \cod{cedar.schema}, y un \emph{binding} que
pida esos roles falla con \cod{ARTIFACT_MISSING}. Por eso el perfil de Flexgrid no lo usa.
\end{notaPractica}

\section{Un YAML más estricto de lo habitual}

El lenguaje usa YAML restringido al modelo de datos de JSON. No se admiten \textbf{anclas}, \textbf{alias} ni
\textbf{etiquetas} explícitas, las claves duplicadas se rechazan y las claves de un mapa tienen que ser cadenas.

La razón es práctica: dos personas que leen el mismo fichero tienen que ver el mismo documento. Un ancla definida
cien líneas más arriba cambia el significado de una línea sin que se note en un \emph{diff}.

\begin{retos}
\begin{enumerate}[leftmargin=*, itemsep=2pt]
  \item Añade un operando \cod{meteringResolution} que exprese la resolución temporal mínima. Pruébalo con
        \cod{xsd:duration} y \cod{gteq}, y luego con \cod{xsd:integer} en minutos. ¿Qué dice el compilador en cada caso?
  \item Cambia \texttt{operators: [lteq, lt]} de \cod{retentionDays} por \texttt{[lteq, isA]} y lee el error.
  \item Quita \cod{included_in} de una acción y comprueba que el aviso no impide compilar.
\end{enumerate}
\end{retos}

\chapter{Compilar un perfil y leer lo que sale}
\label{ch:compilar}

\section{Los dos subcomandos}

\cod{alexandriac} tiene dos subcomandos. \cod{validate} ejecuta todas las comprobaciones e imprime los
hallazgos. \cod{compile} hace lo mismo y además escribe los artefactos.

\begin{lstlisting}[style=sh]
cd tutorial/manual/examples
../../../bin/alexandriac validate dsl/metering-1.0.0.yaml
\end{lstlisting}
\begin{lstlisting}[style=out]
ok: 1 documents valid
\end{lstlisting}

\begin{lstlisting}[style=sh]
../../../bin/alexandriac compile -o /tmp/metering dsl/metering-1.0.0.yaml
\end{lstlisting}
\begin{lstlisting}[style=out]
/tmp/metering/profiles/odrl/metering/1.0.0/spec.html
/tmp/metering/profiles/odrl/metering/1.0.0/vocab.ttl
/tmp/metering/profiles/odrl/metering/1.0.0/context.jsonld
/tmp/metering/profiles/odrl/metering/1.0.0/prof.ttl
/tmp/metering/profiles/odrl/metering/1.0.0/schema.json
/tmp/metering/profiles/odrl/metering/1.0.0/shapes.ttl
\end{lstlisting}

\noindent\begin{tabularx}{\linewidth}{@{}L{4.6cm}Y@{}}
\toprule
\textbf{Opción} & \textbf{Para qué} \\
\midrule
\texttt{-o <dir>} & Directorio de salida de \cod{compile}. Por defecto, \cod{out}. \\
\texttt{-vendor <dir>} & Copias locales de perfiles externos, para \emph{bindings} vendorizados (capítulo~\ref{ch:spaceprofile}). \\
\texttt{-previous <fichero>} & Fuente de la versión anterior, para comprobar el versionado (capítulo~\ref{ch:evolucion}). Repetible. \\
\cod{-log-level}, \cod{-log-format} & Registro estructurado por la salida de error. \\
\bottomrule
\end{tabularx}

\medskip
El código de salida es \cod{0} si no hay errores, \cod{1} si los hay y \cod{2} si el comando está mal invocado.
Los avisos no cambian el código de salida.

El directorio de un perfil es \texttt{profiles/\allowbreak{}odrl/\allowbreak{}<token>/\allowbreak{}<versión>/\allowbreak{}}, donde el \emph{token} es el último
segmento del \cod{id} del perfil: \cod{metering}.

\section{Los seis artefactos de un perfil}

De 114 líneas de YAML salen unos 37~KB repartidos en seis ficheros. Todos dicen lo mismo en idiomas distintos,
para consumidores distintos.

\noindent\begin{tabularx}{\linewidth}{@{}L{2.9cm}L{2.9cm}Y@{}}
\toprule
\textbf{Fichero} & \textbf{Rol} & \textbf{Para quién} \\
\midrule
\cod{vocab.ttl} & \cod{vocabulary} & Quien razona sobre el vocabulario: ontología OWL, colecciones SKOS y firmas. \\
\cod{context.jsonld} & \cod{mapping} & Quien escribe mensajes: un alias corto por término. \\
\cod{schema.json} & \cod{json-schema} & El gate que valida la forma compacta de una política. \\
\cod{shapes.ttl} & \cod{shacl} & El gate que validaría el grafo expandido (capítulo~\ref{ch:shacl}). \\
\cod{prof.ttl} & \cod{prof} & Catálogos y buscadores: manifiesto W3C PROF del perfil. \\
\cod{spec.html} & \cod{specification} & Personas: una tabla legible de términos. \\
\bottomrule
\end{tabularx}

\subsection{\texttt{vocab.ttl}: la ontología}

Cada término creado se emite en la forma RDF de su tipo. Una acción:

\begin{lstlisting}[style=ttl]
enr:forecastDemand
    a odrl:Action, skos:Concept ;
    odrl:includedIn odrl:use ;
    rdfs:isDefinedBy <https://flexgrid.example/odrl/profile/metering> ;
    rdfs:label "Forecast demand"@en, "Predecir la demanda"@es ;
    skos:definition "To train or run load-forecasting models on the metering data."@en .
\end{lstlisting}

Los términos se agrupan en colecciones SKOS por tipo, lo que permite preguntar «qué operandos izquierdos tiene
este perfil» sin conocer su estructura:

\begin{lstlisting}[style=ttl]
<https://flexgrid.example/odrl/profile/metering#leftOperands>
    a skos:Collection ;
    skos:prefLabel "Constraint Left Operands"@en ;
    skos:member odrl:dateTime, odrl:purpose, enr:minAggregation, enr:processingRegion, enr:retentionDays .
\end{lstlisting}

Y cada firma queda reificada, para que una máquina pueda leer el álgebra de un operando:

\begin{lstlisting}[style=ttl]
<https://flexgrid.example/odrl/profile/metering#sig-minAggregation>
    a sig:Signature ;
    sig:leftOperand enr:minAggregation ;
    sig:operator odrl:gteq, odrl:gt ;
    sig:valueKind sig:Literal ;
    sig:allowReference false ;
    sig:datatype xsd:integer .
\end{lstlisting}

De los términos adoptados solo se dice que son miembros. El perfil no afirma nada sobre IRI que no son suyas.

\subsection{\texttt{context.jsonld}: los alias}

\begin{lstlisting}[style=json]
{"@context":{"@version":1.1,
  "aggregate":"http://www.w3.org/ns/odrl/2/aggregate",
  "anonymize":"http://www.w3.org/ns/odrl/2/anonymize",
  "billing":"enr:billing",
  "dateTime":"http://www.w3.org/ns/odrl/2/dateTime",
  "enr":"https://flexgrid.example/ns/metering#",
  "forecastDemand":"enr:forecastDemand",
  "minAggregation":"enr:minAggregation",
  "purpose":"http://www.w3.org/ns/odrl/2/purpose",
  "regionES":"enr:regionES", … }}
\end{lstlisting}

Hay un alias por término, creado o adoptado. Gracias a él una política escribe \texttt{"action": "forecastDemand"}
en lugar de la IRI completa.

\begin{trampaHabitual}[Dos términos no pueden compartir alias]
Si creas un término \cod{purpose} cuando ya has adoptado \cod{odrl:purpose}, los dos querrían el alias
\cod{purpose} y el contexto no sabría distinguirlos. El compilador lo rechaza:

\begin{lstlisting}[style=out]
error   ALIAS_COLLISION   n2-alias.yaml#1:24 /adopted/left_operands/0 enr:purpose and odrl:purpose share the compact alias "purpose"
\end{lstlisting}
\end{trampaHabitual}

El contexto no redefine \cod{use} ni los doce operadores de ODRL, porque el contexto de DSP ya los define y los
protege.

\subsection{\texttt{schema.json}: el gate que hoy se ejecuta}

Es un JSON Schema Draft 2020-12 \textbf{autocontenido}: no tiene ningún \texttt{\$ref} remoto, así que un gate lo
evalúa sin tocar la red. Valida una política ODRL en la forma compacta que viaja en los mensajes DSP.

Comprueba cuatro cosas. La primera, que la política \textbf{declara el perfil}:

\begin{lstlisting}[style=json]
"profile": {"anyOf": [
  {"enum": ["https://flexgrid.example/odrl/profile/metering",
            "https://flexgrid.example/odrl/profile/metering/1.0.0"]},
  {"type": "array", "contains": {"enum": [ … los mismos dos … ]}}
]}
\end{lstlisting}

La segunda, que las \textbf{acciones} están en la lista cerrada, en cualquiera de sus tres grafías: alias, CURIE
o IRI completa. La tercera, que los \textbf{operandos izquierdos} están en su lista. Y la cuarta, que cada
restricción respeta la \textbf{firma} de su operando, con un condicional por firma:

\begin{lstlisting}[style=json]
{
  "if":   {"required": ["leftOperand"],
           "properties": {"leftOperand": {"enum": ["minAggregation", "enr:minAggregation",
                          "https://flexgrid.example/ns/metering#minAggregation"]}}},
  "then": {"required": ["rightOperand"],
           "not": {"required": ["rightOperandReference"]},
           "properties": {
             "operator": {"enum": ["gteq", "odrl:gteq", "http://www.w3.org/ns/odrl/2/gteq",
                                   "gt",   "odrl:gt",   "http://www.w3.org/ns/odrl/2/gt"]},
             "rightOperand": {"anyOf": [
               {"type": "object", "required": ["@value", "@type"],
                "properties": {"@type": {"enum": ["xsd:integer",
                                         "http://www.w3.org/2001/XMLSchema#integer"]},
                               "@value": {"type": "string", "pattern": "^[+-]?[0-9]+$"}}},
               {"type": "integer"}]}}}
}
\end{lstlisting}

Se lee así: \emph{si} el operando izquierdo es \cod{minAggregation}, \emph{entonces} el operador es \cod{gteq} o
\cod{gt}, el operando derecho es obligatorio y es un entero, y no se admite una referencia. Si el entero va
como valor tipado, su \cod{@value} tiene que tener además la forma léxica de un entero: el patrón rechaza
\texttt{"abc"} aunque diga ser \cod{xsd:integer}.

Fuera de las firmas, toda regla tiene que llevar \cod{action} cuando la lista de acciones es cerrada, y toda
restricción atómica tiene que llevar \cod{leftOperand}, \cod{operator} y un operando derecho o una referencia.

\subsection{\texttt{shapes.ttl}: lo mismo, sobre el grafo}

El mismo perfil en SHACL. La firma anterior es esta \emph{shape}:

\begin{lstlisting}[style=ttl]
gen:Sig_minAggregation a sh:NodeShape ;
  sh:targetSubjectsOf odrl:leftOperand ;
  sh:message "Constraint enr:minAggregation: signature not respected." ;
  sh:or (
    [ sh:not [ sh:path odrl:leftOperand ; sh:hasValue enr:minAggregation ] ]
    [
      sh:property [ sh:path odrl:operator ; sh:minCount 1 ; sh:maxCount 1 ; sh:in ( odrl:gteq odrl:gt ) ] ;
      sh:property [ sh:path odrl:rightOperand ; sh:minCount 1 ; sh:maxCount 1 ;
                    sh:nodeKind sh:Literal ; sh:datatype xsd:integer ] ;
      sh:property [ sh:path odrl:rightOperandReference ; sh:maxCount 0 ]
    ]
  ) .
\end{lstlisting}

Es el mismo condicional, escrito como «o no es este operando, o cumple estas tres propiedades». Junto a las
\emph{shapes} de firma, el fichero lleva \cod{gen:ProfileShape}, \cod{gen:ActionShape},
\cod{gen:LeftOperandShape} y \cod{gen:ConstraintShape}, que repiten las demás comprobaciones del esquema. El
capítulo~\ref{ch:shacl} ejecuta estas \emph{shapes} con un motor de verdad y compara los resultados con los del
esquema.

\subsection{\texttt{prof.ttl} y \texttt{spec.html}}

\cod{prof.ttl} es el manifiesto W3C PROF: dice que el perfil es un perfil de ODRL Core, cuál es su \emph{token}
y qué recurso cumple cada rol. \cod{spec.html} es una página con la tabla de términos, su tipo, su estado y su
definición. No tienen efecto en la validación.

\section{Escribir una política contra el perfil}

Con el perfil compilado, esta es una oferta válida. Es lo que la distribuidora ofrece al agregador:

\begin{lstlisting}[style=json]
{
  "@type": "Offer",
  "@id": "urn:uuid:offer-forecast-01",
  "profile": "https://flexgrid.example/odrl/profile/metering",
  "target": "urn:uuid:dataset-qh-load-madrid-2026",
  "permission": [{
    "action": "forecastDemand",
    "constraint": [
      {"leftOperand": "purpose",          "operator": "eq",      "rightOperand": {"@id": "enr:demandForecast"}},
      {"leftOperand": "minAggregation",   "operator": "gteq",    "rightOperand": 50},
      {"leftOperand": "retentionDays",    "operator": "lteq",    "rightOperand": 90},
      {"leftOperand": "processingRegion", "operator": "isAnyOf",
       "rightOperand": [{"@id": "enr:regionES"}, {"@id": "enr:regionEU"}]},
      {"leftOperand": "dateTime",         "operator": "lt",
       "rightOperand": {"@value": "2027-01-01T00:00:00Z", "@type": "xsd:dateTime"}}
    ]
  }]
}
\end{lstlisting}

En castellano: se permite prever la demanda, solo con esa finalidad, sobre registros que agreguen al menos 50
puntos de suministro, guardándolos como mucho 90 días, tratándolos en España o en la UE, hasta final de año.

Tres detalles de forma que conviene interiorizar:

\begin{itemize}[leftmargin=*, itemsep=2pt]
  \item Un valor \textbf{IRI} se escribe como objeto \texttt{\{"@id": …\}}. Una cadena suelta es un literal, y la
        firma lo rechaza.
  \item Un \textbf{entero} se escribe como número JSON, o como valor tipado con \cod{@value} y \cod{@type}.
  \item Una \textbf{fecha} tiene que ir tipada. Una cadena suelta con una fecha dentro es \cod{xsd:string}.
\end{itemize}

\section{Qué acepta y qué rechaza}
\label{sec:tabla-ofertas}

La tabla resume qué hace \cod{schema.json} con variaciones de esa oferta. Todas las filas están ejecutadas; el
capítulo~\ref{ch:shacl} repite la tabla añadiendo la columna de SHACL.

\begin{longtable}{@{}L{8.6cm}C{1.6cm}L{5.2cm}@{}}
\toprule
\textbf{Variación} & \textbf{Esquema} & \textbf{Por qué} \\
\midrule
\endhead
La oferta tal cual & \pasa & \\
Acción \cod{odrl:sell} & \falla & No está en la lista cerrada. \\
\cod{minAggregation} con operador \cod{lt} & \falla & La firma solo admite \cod{gteq} y \cod{gt}. \\
\texttt{minAggregation gteq "fifty"} & \falla & No es un entero. \\
\texttt{minAggregation gteq 50.5} & \falla & No es un entero. \\
\texttt{purpose eq enr:marketing} & \falla & Valor fuera de la enumeración. \\
\texttt{purpose eq "enr:billing"}, como cadena & \falla & Tiene que ser \texttt{\{"@id": …\}}. \\
\texttt{processingRegion isAnyOf} con dos regiones & \pasa & Operador de conjunto: admite lista. \\
\texttt{retentionDays lteq [30, 90]} & \falla & Sin operador de conjunto no hay listas. \\
Acción \cod{use}, del núcleo de ODRL & \pasa & Un perfil no puede prohibirla. \\
Operando \cod{odrl:spatial} & \falla & No está en la lista cerrada de operandos. \\
Restricción lógica \cod{and} con dos hojas válidas & \pasa & Las restricciones lógicas se validan por dentro. \\
Restricción lógica con una hoja inválida & \falla & \\
\texttt{dateTime lt "2027-01-01T00:00:00Z"}, sin tipar & \falla & Falta \cod{@type}. \\
\cod{duty} con una acción ajena & \falla & Los deberes anidados también se validan. \\
Sin el campo \cod{profile} & \falla & La política tiene que declarar el perfil. \\
\cod{rightOperandReference} en \cod{minAggregation} & \falla & \cod{allow_reference} es \cod{false}. \\
Una regla sin \cod{action} & \falla & La lista de acciones es cerrada. \\
Fecha tipada con valor \texttt{"mañana"} & \falla & No tiene la forma léxica del tipo. \\
\bottomrule
\end{longtable}

\section{Determinismo}

Las mismas fuentes producen los mismos bytes, siempre. Los prefijos se ordenan, los términos se ordenan por IRI
y el JSON se emite en forma canónica.

No es un detalle estético. El hub publica cada fichero por su resumen SHA-256 y firma el manifiesto que los
lista. Si compilar dos veces diera bytes distintos, dos hubs con la misma fuente publicarían paquetes distintos.

\begin{lstlisting}[style=sh]
../../../bin/alexandriac compile -o /tmp/a dsl/metering-1.0.0.yaml > /dev/null
../../../bin/alexandriac compile -o /tmp/b dsl/metering-1.0.0.yaml > /dev/null
diff -r /tmp/a /tmp/b && echo identicos
\end{lstlisting}

\begin{retos}
\begin{enumerate}[leftmargin=*, itemsep=2pt]
  \item Cambia \cod{rules.actions} a \cod{open}, recompila y busca la definición \cod{Action} en
        \cod{schema.json}. ¿Qué ha pasado con la enumeración? ¿Y con \cod{gen:ActionShape} en \cod{shapes.ttl}?
  \item Abre \cod{spec.html} en un navegador. ¿Qué columna distingue un término adoptado de uno creado?
\end{enumerate}
\end{retos}

\chapter{El \texttt{SpaceProfile}: cerrar la pila}
\label{ch:spaceprofile}

Un \cod{OdrlProfile} define un vocabulario, pero no dice dónde se aplica. Eso lo dice el \cod{SpaceProfile}: qué
perfil gobierna cada slot, con qué obligación y en qué gate se comprueba. Es además el documento que lleva la
\textbf{versión del espacio}, que es la que publica el hub.

\section{El documento}

\begin{lstlisting}[style=yaml]
dsl: alexandria.profiles/v1
kind: SpaceProfile

space:
  id: https://flexgrid.example/spaces/flexgrid
  version: 1.0.0
  title:
    en: Flexgrid Energy Data Space
    es: Espacio de datos energético Flexgrid
  operator: https://flexgrid.example/party/authority
  base_url: https://flexgrid.example/spaces/flexgrid

lifecycle:
  grace_window: P30D
  propagation_window: P2D
  freshness:
    ttl: PT24H
    max_staleness: P7D
  degraded:
    policy_intake: reject
    dsp_message: continue

bindings:
  - slot: odrl
    profile: https://flexgrid.example/odrl/profile/metering
    source: internal
    version: "^1.0"
    obligation: mandatory
    conformance:
      - { role: json-schema, applies_to: policy_intake }
      - { role: shacl, applies_to: policy_intake }

  - slot: dsp
    profile: https://w3id.org/dspace/2025/1
    source: vendored
    version: "=2025-1"
    obligation: mandatory
    conformance:
      - { role: json-schema, applies_to: dsp_message }
\end{lstlisting}

\section{La sección \texttt{lifecycle}}

Estos parámetros viajan dentro del manifiesto firmado de cada versión, así que todos los clientes aplican los
mismos valores. Lo que no se escribe toma el valor por defecto del hub.

\noindent\begin{tabularx}{\linewidth}{@{}L{5.1cm}C{1.5cm}Y@{}}
\toprule
\textbf{Parámetro} & \textbf{Defecto} & \textbf{Significado} \\
\midrule
\cod{propagation_window} & \cod{P2D} & Tiempo entre que se publica una versión y que empieza a valer. Da margen para sincronizar. \\
\cod{grace_window} & \cod{P30D} & Cuánto convive una versión de ruptura con su predecesora. \\
\cod{min_grace_window} & \cod{P7D} & Gracia mínima que puede pedir un \emph{steward} al aprobar. Es del hub: el \cod{SpaceProfile} no puede cambiarla. \\
\cod{freshness.ttl} & \cod{PT24H} & Validez de un listado. El hub lo reemite al menos cada mitad de este tiempo. \\
\cod{freshness.max_staleness} & \cod{P7D} & Cuánto aguanta un cliente con un listado caducado antes de cerrarse. \\
\cod{sync_interval} & \cod{PT6H} & Intervalo máximo entre sincronizaciones del cliente. \\
\cod{clock_skew} & \cod{PT5S} & Diferencia de reloj tolerada. \\
\cod{lease_margin} & \cod{PT10S} & Margen que el cliente añade alrededor de cada frontera de vigencia. \\
\cod{max_concurrent_versions} & \cod{4} & Máximo de versiones vigentes a la vez. \\
\cod{degraded} & \texttt{\{\}} & Por gate, qué hacer con el listado caducado: \cod{continue} o \cod{reject}. \\
\bottomrule
\end{tabularx}

\medskip
En Flexgrid, \cod{degraded} dice que si el cliente pierde el contacto con el hub sigue validando la envoltura
DSP con lo último que tenía, pero deja de admitir políticas nuevas. El capítulo~\ref{ch:resiliencia} lo provoca.

El compilador comprueba los invariantes entre parámetros. Por ejemplo, con una gracia de tres días y un
\cod{max_staleness} de una hora:

\begin{lstlisting}[style=out]
error   LIFECYCLE_GRACE_TOO_SHORT     s-lifecycle.yaml#2:129 /lifecycle/grace_window grace_window must be at least min_grace_window
error   LIFECYCLE_STALENESS_LT_TTL    s-lifecycle.yaml#2:133 /lifecycle/freshness/max_staleness max_staleness must be greater than ttl
\end{lstlisting}

\section{Los \emph{bindings}}

\noindent\begin{tabularx}{\linewidth}{@{}L{2.6cm}Y@{}}
\toprule
\textbf{Campo} & \textbf{Significado} \\
\midrule
\cod{slot} & \cod{odrl}, \cod{dcat}, \cod{dsp} o \cod{dcp}. Solo un \emph{binding} por slot. \\
\cod{profile} & IRI del perfil que gobierna el slot. \\
\cod{source} & \cod{internal}: un \cod{OdrlProfile} del mismo \emph{bundle}. \cod{vendored}: una copia local de una especificación externa. \\
\cod{version} & Rango SemVer si es interno; versión exacta con \cod{=} si es vendorizado. \\
\cod{obligation} & \cod{mandatory}, \cod{recommended} u \cod{optional}. \\
\cod{via} & Cadena de perfiles intermedios, para perfiles de perfiles. \\
\cod{conformance} & Lista de pares \cod{role} y \cod{applies_to}: qué artefacto se usa en qué gate. \\
\bottomrule
\end{tabularx}

\subsection{Obligación}

La obligación decide qué pasa cuando un mensaje no conforma:

\begin{itemize}[leftmargin=*, itemsep=2pt]
  \item \cod{mandatory}: el gate falla. Un \emph{binding} obligatorio tiene que tener al menos un gate; si no,
        no habría forma de verificarlo.
  \item \cod{recommended}: el gate pasa, y el veredicto lleva el código informativo
        \cod{VERDICT_RECOMMENDED_NONCONFORMANT}.
  \item \cod{optional}: no se evalúa.
\end{itemize}

\subsection{Rangos de versión}

Un \emph{binding} interno no fija una versión del perfil: da un rango, y el compilador elige la mayor versión
del \emph{bundle} que lo cumple.

\noindent\begin{tabularx}{\linewidth}{@{}L{2.6cm}Y L{4.6cm}@{}}
\toprule
\textbf{Rango} & \textbf{Acepta} & \textbf{Ejemplo} \\
\midrule
\cod{=1.2.3} & Exactamente esa versión. & \\
\texttt{\textasciicircum{}1.0} & Misma versión mayor, igual o superior. & \cod{1.0.0}, \cod{1.4.2}; no \cod{2.0.0} \\
\texttt{\textasciicircum{}0.3} & Con mayor 0, misma versión menor. & \cod{0.3.0}, \cod{0.3.9}; no \cod{0.4.0} \\
\texttt{\textasciitilde{}1.2} & Misma mayor y menor. & \cod{1.2.0}, \cod{1.2.7}; no \cod{1.3.0} \\
\cod{1} & Sin operador, las partes que faltan son comodines. & cualquier \cod{1.x.y} \\
\bottomrule
\end{tabularx}

\medskip
Las versiones preliminares, como \cod{1.1.0-rc1}, no satisfacen ningún rango.

\subsection{Roles y gates}

\cod{conformance} une un \textbf{rol} de artefacto con un \textbf{gate}. Un perfil interno ofrece los roles
\cod{json-schema} y \cod{shacl}. Un perfil vendorizado ofrece los que declare su \cod{vendor.json}.

\noindent\begin{tabularx}{\linewidth}{@{}L{5.2cm}Y C{2.4cm}@{}}
\toprule
\textbf{Gate} & \textbf{Qué se valida} & \textbf{¿Lo aplica el cliente?} \\
\midrule
\cod{dsp_message} & La envoltura del mensaje DSP. & \si \\
\cod{policy_intake} & La política de una oferta, una petición o un acuerdo. & \si \\
\cod{catalog_entry} & Las respuestas de catálogo. & \si \\
\cod{pdp_load} & La política al cargarse en el motor de decisión. & \no \\
\cod{credential_presentation} & Las credenciales presentadas. & \no \\
\bottomrule
\end{tabularx}

\begin{trampaHabitual}[Dos gates del lenguaje no los aplica nadie]
El lenguaje admite \cod{pdp_load} y \cod{credential_presentation}, y un \emph{binding} en ellos compila. Pero
\cod{alexandria-client} solo conoce \cod{dsp_message}, \cod{policy_intake}, \cod{catalog_entry} y
\cod{agreement_admission}. Un gate que ningún componente evalúa es documentación, no control.
\end{trampaHabitual}

\begin{trampaHabitual}[Un gate obligatorio cubierto solo por SHACL]
Hoy ningún gate ejecuta SHACL: el cliente solo ejecuta \cod{json-schema}. Si un \emph{binding} \cod{mandatory}
enlaza \emph{únicamente} \cod{shacl} a un gate, ese gate acepta cualquier mensaje. El compilador lo avisa:

\begin{lstlisting}[style=out]
warning GATE_SHACL_ONLY   shacl-only.yaml#2:145 /bindings/0/conformance gate policy_intake of mandatory slot odrl is covered only by shacl, which no gate executes yet: it will accept any message
\end{lstlisting}

Por eso Flexgrid enlaza los dos roles a \cod{policy_intake}: el esquema valida hoy, y las \emph{shapes} quedan
publicadas para cuando haya motor.
\end{trampaHabitual}

\section{Perfiles vendorizados}

DSP no es un perfil que escriba Flexgrid: es una especificación externa que el espacio \textbf{fija}. Como nada se
resuelve contra la red, hay que tener una copia local. El compilador la busca en
\texttt{<vendor>/\allowbreak{}<slot>/\allowbreak{}<versión>/\allowbreak{}vendor.json}:

\begin{lstlisting}[style=json]
{
  "profile": "https://w3id.org/dspace/2025/1",
  "version": "2025-1",
  "artifacts": [
    { "path": "dsp-messages.schema.json", "role": "json-schema",
      "media_type": "application/schema+json" },
    { "path": "context.jsonld", "role": "context",
      "media_type": "application/ld+json",
      "url": "https://w3id.org/dspace/2025/1/context.jsonld" }
  ]
}
\end{lstlisting}

El campo \cod{profile} tiene que coincidir con el del \emph{binding}. Cada artefacto declara su ruta, su rol y
su tipo. El campo \cod{url} es importante en los contextos JSON-LD: es la URL con la que los mensajes los citan,
y el cliente solo resolverá las URL que el paquete declare.

\begin{notaPractica}[La copia de este tutorial es reducida]
Los dos ficheros de \cod{examples/vendor/dsp/2025-1/} están escritos a mano para el laboratorio: un esquema que
comprueba los campos obligatorios de tres mensajes y un contexto con los términos que usan los ejemplos. En un
espacio real se copian ahí los ficheros oficiales de DSP y se fija su versión.
\end{notaPractica}

Sin la opción \cod{-vendor}, el \emph{binding} no resuelve:

\begin{lstlisting}[style=sh]
../../../bin/alexandriac validate dsl/flexgrid-1.0.0.yaml
\end{lstlisting}
\begin{lstlisting}[style=out]
error   BINDING_UNRESOLVED   flexgrid-1.0.0.yaml#2:148 /bindings/1 no vendored copy of https://w3id.org/dspace/2025/1 2025-1 for slot dsp
alexandriac: 1 errors
\end{lstlisting}

\section{El \emph{bundle} y su compilación}

Un \textbf{bundle} es un flujo YAML: un \cod{SpaceProfile} y los \cod{OdrlProfile} que enlaza, separados por
\cod{---}. \cod{examples/dsl/flexgrid-1.0.0.yaml} es exactamente el perfil del capítulo anterior seguido del
espacio de este.

\begin{lstlisting}[style=sh]
../../../bin/alexandriac compile -o /tmp/flexgrid -vendor vendor dsl/flexgrid-1.0.0.yaml | sort
\end{lstlisting}
\begin{lstlisting}[style=out]
/tmp/flexgrid/profiles/dsp/vendored/2025-1/context.jsonld
/tmp/flexgrid/profiles/dsp/vendored/2025-1/dsp-messages.schema.json
/tmp/flexgrid/profiles/odrl/metering/1.0.0/context.jsonld
/tmp/flexgrid/profiles/odrl/metering/1.0.0/prof.ttl
/tmp/flexgrid/profiles/odrl/metering/1.0.0/schema.json
/tmp/flexgrid/profiles/odrl/metering/1.0.0/shapes.ttl
/tmp/flexgrid/profiles/odrl/metering/1.0.0/spec.html
/tmp/flexgrid/profiles/odrl/metering/1.0.0/vocab.ttl
/tmp/flexgrid/space/prof.ttl
/tmp/flexgrid/space/shacl-union.ttl
/tmp/flexgrid/space/space-lock.json
/tmp/flexgrid/space/spec.html
\end{lstlisting}

Aparecen dos cosas nuevas: los ficheros vendorizados, copiados bajo \texttt{profiles/\allowbreak{}<slot>/\allowbreak{}vendored/\allowbreak{}<versión>/\allowbreak{}}, y
cuatro artefactos del espacio.

\noindent\begin{tabularx}{\linewidth}{@{}L{4.8cm}Y@{}}
\toprule
\textbf{Artefacto del espacio} & \textbf{Contenido} \\
\midrule
\cod{space/space-lock.json} & La resolución: versión elegida de cada \emph{binding}, sus artefactos con SHA-256 y gates, los parámetros de ciclo de vida y los perfiles declarables. \\
\cod{space/prof.ttl} & Manifiesto PROF del espacio: un \cod{prof:isProfileOf} por \emph{binding} y un descriptor por artefacto. \\
\cod{space/shacl-union.ttl} & Las \emph{shapes} de los \emph{bindings} obligatorios, concatenadas por gate. \\
\cod{space/spec.html} & Tabla legible de \emph{bindings}. \\
\bottomrule
\end{tabularx}

\medskip
El \cod{space-lock.json} es el equivalente a un fichero de bloqueo de dependencias. Donde el YAML decía
\texttt{"\textasciicircum{}1.0"}, el bloqueo dice qué se eligió y con qué bytes:

\begin{lstlisting}[style=json]
{ "slot": "odrl", "source": "internal", "obligation": "mandatory",
  "profile": "https://flexgrid.example/odrl/profile/metering",
  "version": "1.0.0",
  "artifacts": [
    { "path": "profiles/odrl/metering/1.0.0/schema.json", "role": "json-schema",
      "media_type": "application/schema+json", "applies_to": ["policy_intake"],
      "sha256": "fa70ada7f8c3e45426e804c1883da0a338d3e718425211ad86626d702cdc5e09" },
    { "path": "profiles/odrl/metering/1.0.0/shapes.ttl", "role": "shacl",
      "media_type": "text/turtle", "applies_to": ["policy_intake"],
      "sha256": "43d5fb036d58eea976c45b4af11575ec0e012162ce2b29ad35185900eb4929c8" },
    … ] }
\end{lstlisting}

\begin{ideaClave}
Cada artefacto lleva los gates de su \emph{binding}. \cod{schema.json} del slot \cod{odrl} se aplica en
\cod{policy_intake} porque el \cod{SpaceProfile} lo dice, no porque el artefacto lo sepa. Cambiar dónde se valida
algo es editar una línea del espacio, sin tocar el perfil.
\end{ideaClave}

\section{Errores propios del espacio}

\noindent\begin{tabularx}{\linewidth}{@{}L{4.6cm}Y@{}}
\toprule
\textbf{Código} & \textbf{Causa} \\
\midrule
\cod{SLOT_DUP} & Dos \emph{bindings} para el mismo slot. \\
\cod{BINDING_UNRESOLVED} & Ninguna versión del \emph{bundle} cumple el rango, o falta la copia vendorizada. \\
\cod{BINDING_BASE_MISMATCH} & Un slot distinto de \cod{odrl} declarado como \cod{internal}: el lenguaje solo tiene documentos \cod{OdrlProfile}. \\
\cod{ARTIFACT_MISSING} & Se pide un rol que el perfil no produce, como \cod{cedar-mapping}. \\
\cod{MANDATORY_UNVERIFIABLE} & Un \emph{binding} obligatorio sin gates. El esquema lo detiene antes, como error \cod{SCHEMA}. \\
\cod{VIA_CYCLE} & \cod{via} contiene el propio perfil. \\
\bottomrule
\end{tabularx}

\medskip
Las salidas reales de los dos primeros:

\begin{lstlisting}[style=out]
error   SLOT_DUP             s-slot-dup.yaml#2:156 /bindings/2/slot slot odrl has more than one binding
error   BINDING_UNRESOLVED   s-unresolved.yaml#2:139 /bindings/0 no release of https://flexgrid.example/odrl/profile/metering satisfies ^2.0 in the bundle
\end{lstlisting}

En todos los mensajes, \texttt{\#2} es el número de documento dentro del flujo YAML, y lo que sigue a los dos puntos
es la línea del fichero.

\begin{retos}
\begin{enumerate}[leftmargin=*, itemsep=2pt]
  \item Cambia la obligación del slot \cod{odrl} a \cod{recommended} y compara el \cod{space-lock.json}.
  \item Añade un \emph{binding} \cod{dcat} vendorizado sin crear su directorio y lee el error. Después crea
        \cod{vendor/dcat/3.0.0/vendor.json} con un artefacto de rol \cod{shacl} y enlázalo a \cod{catalog_entry}.
  \item Sube \cod{max_concurrent_versions} a 1 y piensa qué implicaría para una versión de ruptura.
\end{enumerate}
\end{retos}

\chapter{Los cinco niveles de validación}
\label{ch:validacion}

El compilador valida en cinco niveles, del más superficial al más profundo. Saber en qué nivel falla algo dice
mucho de qué está mal.

\begin{center}
\begin{tikzpicture}[font=\small, node distance=3mm,
  lvl/.style={draw=alxNavy, fill=alxLightBlue, rounded corners=2pt, minimum height=10mm, minimum width=25mm, align=center},
  arr/.style={-{Stealth[length=2mm]}, alxNavy, thick}]
  \node[lvl] (n1) {\textbf{N1}\\[-2pt]\scriptsize sintaxis};
  \node[lvl, right=of n1] (n2) {\textbf{N2}\\[-2pt]\scriptsize invariantes};
  \node[lvl, right=of n2] (n3) {\textbf{N3}\\[-2pt]\scriptsize referencias};
  \node[lvl, right=of n3] (n4) {\textbf{N4}\\[-2pt]\scriptsize versionado};
  \node[lvl, right=of n4, fill=alxLightTeal, draw=alxTeal] (em) {emisión};
  \node[lvl, right=of em] (n5) {\textbf{N5}\\[-2pt]\scriptsize salidas};
  \draw[arr] (n1) -- (n2); \draw[arr] (n2) -- (n3); \draw[arr] (n3) -- (n4); \draw[arr] (n4) -- (em); \draw[arr] (em) -- (n5);
\end{tikzpicture}
\end{center}

\noindent\begin{tabularx}{\linewidth}{@{}C{1cm}L{3.2cm}Y@{}}
\toprule
\textbf{Nivel} & \textbf{Pregunta} & \textbf{Qué comprueba} \\
\midrule
N1 & ¿Es un documento del lenguaje? & YAML restringido al modelo JSON y el esquema normativo (\cod{internal/dsl/schema.json}). \\
N2 & ¿Es coherente consigo mismo? & Identificadores únicos, prefijos declarados, tabla operador–tipo, un \emph{binding} por slot, invariantes de ciclo de vida. \\
N3 & ¿Existe lo que cita? & Términos adoptados, resolución de rangos, copias vendorizadas, artefactos de cada gate. \\
N4 & ¿Puede suceder a la versión anterior? & Diferencias contra la versión previa y el salto de versión mínimo. \\
N5 & ¿Sirve lo generado? & El JSON es válido, el esquema compila, el contexto expande cada alias a su IRI y el Turtle está bien formado. \\
\bottomrule
\end{tabularx}

\medskip
Si hay errores hasta N4, no se emite nada. Si N5 falla, los artefactos se descartan.

\section{Cómo leer un hallazgo}

Cada línea tiene el mismo formato:

\begin{lstlisting}[style=out]
error   SIG_OPERATOR_TYPE    n2-operator-type.yaml#1:73 /terms/4/signature/operators/1 isA requires value_kind iri
\end{lstlisting}

\noindent\begin{tabularx}{\linewidth}{@{}L{6.4cm}Y@{}}
\toprule
\textbf{Parte} & \textbf{Significado} \\
\midrule
\cod{error} & Severidad: \cod{error} bloquea, \cod{warning} no. \\
\cod{SIG_OPERATOR_TYPE} & Código estable, pensado para buscarlo. \\
\texttt{n2-operator-type.yaml\#1:73} & Fichero, número de documento dentro del flujo y línea. \\
\cod{/terms/4/signature/operators/1} & Puntero JSON al nodo culpable. \\
El resto & Mensaje para personas. \\
\bottomrule
\end{tabularx}

\medskip
Estos códigos son del lenguaje. No pertenecen al catálogo cerrado de 77 códigos del hub y del cliente, aunque
cuando el hub compila un borrador los devuelve tal cual en la respuesta.

\section{N1: sintaxis}

\subsection{Un ancla}

\begin{lstlisting}[style=out]
error   YAML     n1-anchor.yaml#1:1  line 93: anchors are not allowed (§6)
\end{lstlisting}

\subsection{Un campo que no existe}

\begin{lstlisting}[style=out]
error   SCHEMA   n1-unknown-field.yaml#1:1  'oneOf' failed, none matched
error   SCHEMA   n1-unknown-field.yaml#1:1  validation failed
error   SCHEMA   n1-unknown-field.yaml#1:1  missing properties 'space', 'bindings'
error   SCHEMA   n1-unknown-field.yaml#1:1  additional properties 'namespaces', 'profile', 'rules', 'terms', 'adopted' not allowed
error   SCHEMA   n1-unknown-field.yaml#1:2 /kind value must be 'SpaceProfile'
error   SCHEMA   n1-unknown-field.yaml#1:113 /rules validation failed
alexandriac: 6 errors
\end{lstlisting}

\begin{trampaHabitual}[Seis errores para una errata]
El esquema de un documento es «o es un \cod{OdrlProfile}, o es un \cod{SpaceProfile}». Cuando una rama falla, el
validador informa de por qué no encaja en \emph{ninguna}, y las cinco primeras líneas explican por qué tu perfil
no es un \cod{SpaceProfile}, lo cual es cierto e inútil.

La línea que importa es \textbf{la última}, la que tiene puntero y número de línea profundos: \cod{/rules} en la
línea 113. Ahí está el campo mal escrito.
\end{trampaHabitual}

\section{N2: invariantes}

\begin{longtable}{@{}L{6.3cm}L{9.3cm}@{}}
\toprule
\textbf{Código} & \textbf{Causa y mensaje real} \\
\midrule
\endhead
\cod{TERM_DUP_ID} & Dos términos con el mismo \cod{id}. \newline {\scriptsize\ttfamily term id gridBalancing is repeated} \\
\cod{ALIAS_COLLISION} & Un término creado y uno adoptado comparten nombre local. \newline {\scriptsize\ttfamily enr:purpose and odrl:purpose share the compact alias "purpose"} \\
\cod{PREFIX_UNDECLARED} & Un CURIE usa un prefijo que no está en \cod{namespaces}. \newline {\scriptsize\ttfamily prefix of geo:ES is not declared} \\
\cod{SIG_OPERATOR_TYPE} & Un operador no casa con el tipo. \newline {\scriptsize\ttfamily gteq requires an ordered datatype, not ""} \newline {\scriptsize\ttfamily isA requires value\_kind iri} \\
\cod{SIG_OPERATOR_UNKNOWN} & El operador no es de ODRL ni un \cod{operator} del perfil. \\
\cod{SIG_VALUE_KIND} & \cod{datatype} que no es \cod{xsd:}, o \cod{datatype} junto a \texttt{value\_kind: iri}. Como aviso: un valor IRI que no es operando derecho del perfil. \\
\cod{ENFORCEMENT_WITHOUT_SIGNATURE} & Bloque \cod{enforcement} sin firma. \newline {\scriptsize\ttfamily enforcement of enr:retentionDays requires a signature (§10)} \\
\cod{PROFILE_ID_IS_NAMESPACE} & El \cod{id} del perfil coincide con un espacio de nombres. \\
\cod{SLOT_DUP}, \cod{VIA_CYCLE} & Del espacio, vistos en el capítulo anterior. \\
\cod{LIFECYCLE_*} & Parámetros de ciclo de vida incoherentes. \\
\bottomrule
\end{longtable}

Los avisos de N2 no bloquean, pero conviene no ignorarlos:

\noindent\begin{tabularx}{\linewidth}{@{}L{6.3cm}Y@{}}
\toprule
\textbf{Aviso} & \textbf{Qué te está diciendo} \\
\midrule
\cod{ACTION_INCLUDED_IN_MISSING} & Una acción sin \cod{included_in}: un motor no sabrá de qué acción es un caso. \\
\cod{TERM_NO_REPLACEMENT} & Un término obsoleto o retirado sin \cod{replaced_by}: quien lo usa no sabe a qué migrar. \\
\cod{SIG_OPERATOR_TYPE} & Un operador de orden sobre \cod{xsd:duration}, que solo está parcialmente ordenado. \\
\cod{SIG_SCHEME_UNSUPPORTED} & Se usa \cod{scheme}, que hoy no es practicable. \\
\cod{EXTRA_TURTLE_UNCHECKED} & \cod{extra_turtle} se inyecta tal cual: el compilador no puede analizarlo. \\
\cod{GATE_SHACL_ONLY} & Un gate de un \emph{binding} obligatorio solo tiene \cod{shacl}, que nadie ejecuta: lo aceptará todo. \\
\cod{CONFORMANCE_REQUIRED_IGNORED} & \cod{lifecycle.conformance_required} se acepta y se ignora. \\
\bottomrule
\end{tabularx}

\section{N3: referencias}

\begin{lstlisting}[style=out]
error   ADOPTED_UNKNOWN         n3-adopted-unknown.yaml#1:22 /adopted/actions/1 odrl:pseudonymize is not in the ODRL 2.2 vocabulary
error   ADOPTED_KIND_MISMATCH   n3-adopted-kind.yaml#1:22 /adopted/actions/1 odrl:spatial is a left_operand in ODRL, not a action
\end{lstlisting}

El primero es el error típico de quien supone que ODRL tiene un término y no lo tiene: ODRL define
\cod{anonymize}, pero no \cod{pseudonymize}. Si lo necesitas, lo creas tú en \cod{terms}.

El segundo detecta la categoría equivocada: \cod{odrl:spatial} existe, pero es un operando izquierdo, no una
acción.

De N3 son también \cod{BINDING_UNRESOLVED}, \cod{BINDING_BASE_MISMATCH} y \cod{ARTIFACT_MISSING}, del capítulo
anterior. Este último tiene una salida reveladora:

\begin{lstlisting}[style=out]
error   ARTIFACT_MISSING   s-artifact.yaml#2:146 /bindings/0/conformance/1 https://flexgrid.example/odrl/profile/metering 1.0.0 produces no cedar-mapping artefact
\end{lstlisting}

\section{N5: salidas}

N5 revisa lo que el propio compilador acaba de generar. Puede parecer redundante, pero es la red de seguridad
contra los fallos de los emisores y contra \cod{extra_turtle}, que es texto del autor.

\noindent\begin{tabularx}{\linewidth}{@{}L{3.6cm}Y@{}}
\toprule
\textbf{Artefacto} & \textbf{Comprobación} \\
\midrule
\cod{schema.json} & Se compila como JSON Schema. \\
Otros JSON & Son JSON válido. \\
\cod{context.jsonld} & Se procesa sin red y cada alias se expande a la IRI de su término. \\
Turtle & Corchetes y paréntesis equilibrados, cadenas e IRI cerradas, punto final. \\
\bottomrule
\end{tabularx}

\begin{hallazgo}{La comprobación de Turtle es solo estructural}
N5 no analiza el Turtle: cuenta corchetes. Tampoco ejecuta las \emph{shapes}. Un \cod{shapes.ttl} que fuera
SHACL sin sentido, pero con los paréntesis en su sitio, pasaría. Es el hallazgo H2 del capítulo~\ref{ch:shacl}.
\end{hallazgo}

\section{N4 tiene capítulo propio}

El versionado compara dos versiones y por eso necesita la anterior. Es el tema del capítulo siguiente.

\begin{retos}
\begin{enumerate}[leftmargin=*, itemsep=2pt]
  \item Provoca los seis códigos de error de la tabla de N2 modificando una copia de
        \cod{metering-1.0.0.yaml}. Todos se consiguen cambiando una sola línea.
  \item Escribe \texttt{version: v1} en el perfil y localiza, entre las seis líneas de error, la única que sirve.
\end{enumerate}
\end{retos}

\chapter{Evolucionar un perfil sin romper a nadie}
\label{ch:evolucion}

Un vocabulario publicado es un contrato. Hay políticas escritas contra él, acuerdos firmados con él y
conectores configurados para él. El nivel N4 existe para que cambiarlo sea una decisión consciente y no un
accidente.

\section{Las tres reglas}

\begin{ideaClave}
\begin{enumerate}[leftmargin=*, itemsep=1pt]
  \item Una versión publicada \textbf{no cambia nunca}. Mismo número, mismo contenido.
  \item Un término \textbf{no se borra}. Pasa de activo a obsoleto y, en una versión posterior, a retirado.
  \item El número de versión tiene que \textbf{decir la verdad}: ampliar es un cambio menor, restringir es uno mayor.
\end{enumerate}
\end{ideaClave}

Para comprobarlas, el compilador necesita la fuente de la versión anterior. En la línea de comandos se le da con
\cod{-previous}. En el hub no hay que hacer nada: cada paquete publicado guarda su fuente en
\cod{sources/draft.yaml}, y el hub compara el borrador nuevo con ella.

\section{Qué salto exige cada cambio}

\subsection{En un perfil}

\begin{longtable}{@{}L{9.4cm}C{2cm}@{}}
\toprule
\textbf{Cambio} & \textbf{Salto mínimo} \\
\midrule
\endhead
Título, descripción o editor del perfil & PARCHE \\
Etiqueta, definición, nota o \cod{replaced_by} de un término & PARCHE \\
\midrule
Término nuevo, creado o adoptado & MENOR \\
Un término pasa de activo a obsoleto & MENOR \\
Operador añadido a una firma & MENOR \\
Valor añadido a una enumeración & MENOR \\
Se permiten referencias donde no se permitían & MENOR \\
Se quita la firma de un operando & MENOR \\
Se abre una lista cerrada & MENOR \\
Cambia \cod{included_in}, \cod{subclass_of} o \cod{extra_turtle} & MENOR \\
Cambia la licencia o \cod{base_url} & MENOR \\
\midrule
Un término obsoleto pasa a retirado & MAYOR \\
Operador quitado de una firma & MAYOR \\
Valor quitado de una enumeración & MAYOR \\
Cambia el tipo de dato o el tipo de valor de una firma & MAYOR \\
Se prohíben referencias que se permitían & MAYOR \\
Se añade una firma a un operando que no la tenía & MAYOR \\
Se cierra una lista abierta & MAYOR \\
Se deja de adoptar un término & MAYOR \\
Cambia el tipo (\cod{kind}) de un término & MAYOR \\
\bottomrule
\end{longtable}

La lógica es siempre la misma: si toda política que antes era válida lo sigue siendo, el cambio es menor. Si
alguna deja de serlo, es mayor.

\subsection{En un espacio}

\begin{longtable}{@{}L{9.4cm}C{2cm}@{}}
\toprule
\textbf{Cambio} & \textbf{Salto mínimo} \\
\midrule
\endhead
Título o descripción del espacio & PARCHE \\
\midrule
Operador, \cod{base_url} o cualquier parámetro de \cod{lifecycle} & MENOR \\
\emph{Binding} no obligatorio añadido o quitado & MENOR \\
Se baja una obligación que no era \cod{mandatory} & MENOR \\
Cambian los gates de un \emph{binding} & MENOR \\
\midrule
\emph{Binding} obligatorio añadido o quitado & MAYOR \\
Cambia el perfil o el origen de un slot & MAYOR \\
Se sube una obligación, o se baja desde \cod{mandatory} & MAYOR \\
Se añade un gate a un slot obligatorio & MAYOR \\
Cambia la versión de un perfil vendorizado & MAYOR \\
\bottomrule
\end{longtable}

Además, cuando el perfil interno de un slot cambia de versión, el espacio tiene que saltar \textbf{al menos lo
mismo}. Si el perfil pasa de 1.1.0 a 2.0.0, el espacio no puede quedarse en un salto menor.

\section{La versión 1.1.0: ampliar y avisar}

La universidad entra en Flexgrid y hace falta una finalidad nueva. Al mismo tiempo, el regulador anuncia que el
perfilado individual dejará de ser admisible. Son cuatro cambios sobre el 1.0.0:

\begin{lstlisting}[style=out]
6c6
<   version: 1.0.0
---
>   version: 1.1.0
27c27
<         values: [enr:gridBalancing, enr:billing, enr:demandForecast]
---
>         values: [enr:gridBalancing, enr:billing, enr:demandForecast, enr:research]
55a56,57
>     status: deprecated
>     replaced_by: enr:forecastDemand
98a101,104
>   - id: research
>     kind: right_operand
>     label: { en: Research }
>     definition: { en: Non-commercial scientific research on energy systems. }
121c127
<   version: 1.0.0
---
>   version: 1.1.0
\end{lstlisting}

Un valor añadido, un término nuevo y una obsolescencia: los tres son cambios menores, y 1.0.0 a 1.1.0 es un
salto menor. El \emph{binding} del espacio no se toca, porque \texttt{"\textasciicircum{}1.0"} ya admite el 1.1.0.

\begin{lstlisting}[style=sh]
../../../bin/alexandriac validate -vendor vendor \
  -previous dsl/flexgrid-1.0.0.yaml dsl/flexgrid-1.1.0.yaml
\end{lstlisting}
\begin{lstlisting}[style=out]
ok: 2 documents valid
\end{lstlisting}

En los artefactos, la obsolescencia aparece en dos sitios. En \cod{vocab.ttl}:

\begin{lstlisting}[style=ttl]
enr:profileConsumer
    a odrl:Action, skos:Concept ;
    odrl:includedIn odrl:use ;
    rdfs:isDefinedBy <https://flexgrid.example/odrl/profile/metering> ;
    rdfs:label "Profile consumer"@en, "Perfilar al consumidor"@es ;
    skos:definition "To derive a behavioural profile of an individual supply point."@en ;
    dct:isReplacedBy enr:forecastDemand ;
    owl:deprecated true ;
    alx:status "deprecated" .
\end{lstlisting}

Y en \cod{shapes.ttl}, como una \emph{shape} de severidad \cod{sh:Warning}:

\begin{lstlisting}[style=ttl]
gen:DeprecatedActionShape a sh:NodeShape ;
  sh:targetObjectsOf odrl:permission, odrl:prohibition, odrl:obligation, odrl:duty ;
  sh:property [ sh:path odrl:action ; sh:severity sh:Warning ;
    sh:message "Deprecated term of profile metering." ;
    sh:not [ sh:in ( enr:profileConsumer ) ] ] .
\end{lstlisting}

\begin{trampaHabitual}[El aviso de obsolescencia solo existe en SHACL]
\cod{schema.json} no distingue un término obsoleto de uno activo: la acción sigue en la enumeración y la
política pasa sin más. El aviso vive solo en \cod{shapes.ttl}.

Como el cliente no ejecuta SHACL, hoy \textbf{nadie recibe ese aviso}. Una comercializadora que siga usando
\cod{profileConsumer} bajo la 1.1.0 no ve nada hasta que la 2.0.0 la rechaza. Es la consecuencia más visible de
la parte de SHACL sin hacer, y el motivo de que el código \cod{VERDICT_SHACL_WARNING} sea inalcanzable.
\end{trampaHabitual}

\section{La versión 2.0.0: retirar}

\begin{lstlisting}[style=out]
6c6
<   version: 1.1.0
---
>   version: 2.0.0
56c56
<     status: deprecated
---
>     status: retired
127c127
<   version: 1.1.0
---
>   version: 2.0.0
148c148
<     version: "^1.0"
---
>     version: "^2.0"
\end{lstlisting}

Retirar es un cambio mayor: una política con \cod{profileConsumer} deja de ser válida. El perfil sube a 2.0.0,
el \emph{binding} pasa a \texttt{"\textasciicircum{}2.0"} y, como el perfil ha dado un salto mayor, el espacio también.

El término retirado sigue en \cod{vocab.ttl}, con \texttt{alx:status "retired"}: su IRI no se reutiliza y quien
encuentre una política antigua puede seguir sabiendo qué significaba. Lo que desaparece es su presencia en las
listas cerradas de \cod{schema.json} y \cod{shapes.ttl}.

\begin{lstlisting}[style=sh]
../../../bin/alexandriac validate -vendor vendor \
  -previous dsl/flexgrid-1.1.0.yaml dsl/flexgrid-2.0.0.yaml
\end{lstlisting}
\begin{lstlisting}[style=out]
ok: 2 documents valid
\end{lstlisting}

\section{Lo que el compilador no deja hacer}

\subsection{Retirar sin haber avisado}

Si se intenta validar la 2.0.0 directamente contra la 1.0.0, saltándose la 1.1.0:

\begin{lstlisting}[style=sh]
../../../bin/alexandriac validate -vendor vendor \
  -previous dsl/flexgrid-1.0.0.yaml dsl/flexgrid-2.0.0.yaml
\end{lstlisting}
\begin{lstlisting}[style=out]
error   TERM_RETIRED_WITHOUT_DEPRECATION   flexgrid-2.0.0.yaml#1:56 /terms/2/status enr:profileConsumer must be deprecated in a previous release before it is retired
alexandriac: 1 errors
\end{lstlisting}

Hay dos formas correctas de llegar a la 2.0.0. La primera es publicar antes la 1.1.0, que es lo que hace este
tutorial. La segunda, si no se quiere publicar una versión intermedia del espacio, es \textbf{incluir el perfil
1.1.0 en el mismo \emph{bundle}} que el 2.0.0. El \emph{binding} elige el 2.0.0, el 1.1.0 no genera artefactos,
pero su presencia demuestra que la obsolescencia existió. Así lo hace el escenario F6 del repositorio
(\cod{conformance/f6/drafts/v2.0.0.yaml}).

\subsection{Restringir con un salto menor}

Se quita una región de la enumeración y se sube solo la versión menor:

\begin{lstlisting}[style=out]
error   VERSION_BUMP_TOO_LOW   bad-bump.yaml#1:6 /profile/version 1.0.0 → 1.1.0 is a MINOR bump; the changes require MAJOR: [value removed from enr:processingRegion]
\end{lstlisting}

El mensaje dice el salto que se ha hecho, el que hacía falta y por qué. Cuando hay varios motivos, los lista
todos.

\subsection{Borrar un término}

\begin{lstlisting}[style=out]
error   TERM_REMOVED   bad-removed.yaml#1:36 /terms term enr:profileConsumer was removed; terms are deprecated and retired, never removed
\end{lstlisting}

\subsection{Cambiar una versión ya publicada}

Se corrige una errata en una definición sin cambiar el número de versión:

\begin{lstlisting}[style=out]
error   PROFILE_VERSION_CHANGED   changed.yaml#1:6 /profile/version https://flexgrid.example/odrl/profile/metering 1.0.0 is already released with other content
\end{lstlisting}

Una errata es un cambio de PARCHE: 1.0.1.

\section{Las clases de versión del hub}

El hub clasifica cada versión \textbf{del espacio} por su número, y la clase decide sus fechas:

\noindent\begin{tabularx}{\linewidth}{@{}L{2.6cm}L{4.4cm}Y@{}}
\toprule
\textbf{Clase} & \textbf{Cuándo} & \textbf{Efecto} \\
\midrule
\cod{initial} & La primera versión. & Entra en vigor tras la ventana de propagación. \\
\cod{compatible} & Salto menor o de parche. & Sustituye a la anterior tras la ventana de propagación. No hay convivencia porque no hace falta. \\
\cod{breaking} & Salto mayor. & Convive con la anterior durante la ventana de gracia. \\
\bottomrule
\end{tabularx}

\medskip
N4 es lo que garantiza que esa clasificación sea honesta. Sin él, alguien podría retirar un término en una
versión «compatible» y dejar a los participantes sin ventana de gracia.

\begin{retos}
\begin{enumerate}[leftmargin=*, itemsep=2pt]
  \item Sobre la 2.0.0, añade el operador \cod{neq} a \cod{odrl:purpose}. ¿Qué número de versión toca? Compruébalo.
  \item Ahora quita \cod{isAnyOf} de \cod{processingRegion} y vuelve a preguntarte lo mismo.
  \item Cambia \cod{degraded.policy_intake} a \cod{continue} en el espacio. ¿Es un cambio menor o mayor? ¿Estás
        de acuerdo con el compilador?
\end{enumerate}
\end{retos}

\chapter{Un segundo perfil: uso secundario de datos de salud}
\label{ch:salud}

El perfil de medida es deliberadamente regular. Este segundo ejemplo, \cod{examples/dsl/health-0.3.0.yaml},
enseña las piezas del lenguaje que aquel no necesita. El dominio es el uso secundario de datos de salud: un
organismo de acceso emite un permiso de datos, y los datos se tratan en un entorno seguro, seudonimizados y con
finalidades tasadas.

\section{Las piezas nuevas}

\subsection{Un operador propio}

Los doce operadores de ODRL comparan valores. Aquí hace falta decir algo distinto: que un uso está \emph{cubierto}
por un permiso. Eso es un operador nuevo:

\begin{lstlisting}[style=yaml]
  - id: withinPermit
    kind: operator
    label: { en: Within permit }
    definition: { en: The left operand is covered by the data permit given as right operand. }
\end{lstlisting}

Y se usa en una firma citándolo por su CURIE, junto a un operador estándar:

\begin{lstlisting}[style=yaml]
  - id: dataPermit
    kind: left_operand
    label: { en: Data permit, es: Permiso de datos }
    definition: { en: Data permit issued by the health data access body that covers the use. }
    signature:
      datatype: xsd:anyURI
      operators: [eq, hds:withinPermit]
      allow_reference: true
\end{lstlisting}

El compilador comprueba que \cod{hds:withinPermit} exista y sea de tipo \cod{operator}. Si no, falla con
\cod{SIG_OPERATOR_UNKNOWN}.

\subsection{Referencias en el operando derecho}

\texttt{allow\_reference: true} permite que la restricción no lleve un valor, sino una referencia a un recurso
externo con \cod{rightOperandReference}. Tiene sentido para un permiso de datos, que es un documento que vive en
otro sitio. La \emph{shape} generada pierde entonces dos restricciones: el operando derecho deja de ser
obligatorio y la referencia deja de estar prohibida. Que haya al menos uno de los dos lo exige una
\emph{shape} general, \cod{gen:ConstraintShape}, igual que hace el esquema.

\begin{lstlisting}[style=ttl]
gen:Sig_dataPermit a sh:NodeShape ;
  sh:targetSubjectsOf odrl:leftOperand ;
  sh:message "Constraint hds:dataPermit: signature not respected." ;
  sh:or (
    [ sh:not [ sh:path odrl:leftOperand ; sh:hasValue hds:dataPermit ] ]
    [
      sh:property [ sh:path odrl:operator ; sh:minCount 1 ; sh:maxCount 1 ; sh:in ( odrl:eq hds:withinPermit ) ] ;
      sh:property [ sh:path odrl:rightOperand ; sh:maxCount 1 ; sh:nodeKind sh:Literal ; sh:datatype xsd:anyURI ]
    ]
  ) .
\end{lstlisting}

\subsection{Una enumeración de literales}

En el perfil de medida las enumeraciones eran de IRI. Cuando los valores no son CURIE, el compilador deduce que
son literales:

\begin{lstlisting}[style=yaml]
  - id: environment
    kind: left_operand
    label: { en: Processing environment, es: Entorno de tratamiento }
    definition: { en: Kind of environment in which the data may be processed. }
    signature:
      values: [secure-processing-environment, federated-node]
      operators: [eq]
\end{lstlisting}

En la política, el valor es entonces una cadena: \texttt{"rightOperand": "federated-node"}.

\begin{notaPractica}[¿IRI o literal?]
Una enumeración de IRI obliga a definir cada valor como término, con etiqueta y definición, y eso es lo que la
hace interoperable: otro espacio puede adoptar tus valores. Una enumeración de literales es más rápida de
escribir y no significa nada fuera de tu esquema. Para finalidades, regiones o categorías de datos, IRI. Para
etiquetas internas, literales.
\end{notaPractica}

\subsection{Una subclase de regla}

\begin{lstlisting}[style=yaml]
  - id: SecondaryUsePermission
    kind: rule_subclass
    subclass_of: odrl:Permission
    label: { en: Secondary use permission }
    definition: { en: A permission granted under a data permit for secondary use. }
\end{lstlisting}

Los tres tipos \cod{*_subclass} exigen \cod{subclass_of} y generan una clase OWL:

\begin{lstlisting}[style=ttl]
hds:SecondaryUsePermission
    a owl:Class ;
    rdfs:subClassOf odrl:Permission ;
    rdfs:isDefinedBy <https://salud.example/odrl/profile/secondary-use> ;
    rdfs:label "Secondary use permission"@en ;
    skos:definition "A permission granted under a data permit for secondary use."@en .
\end{lstlisting}

\subsection{Una lista abierta}

\begin{lstlisting}[style=yaml]
rules:
  actions: closed
  left_operands: open
\end{lstlisting}

Las acciones siguen cerradas, pero una política puede llevar operandos izquierdos que el perfil no conoce. Los
que sí conoce se siguen validando contra su firma. En \cod{shapes.ttl} desaparece \cod{gen:LeftOperandShape}, y
en \cod{schema.json} desaparece la enumeración de \cod{leftOperand}.

Es la opción razonable para un perfil joven, versión \cod{0.3.0}, que todavía no sabe qué restricciones va a
necesitar. Cerrar la lista más adelante será un cambio mayor.

\section{Los avisos}

\begin{lstlisting}[style=sh]
../../../bin/alexandriac validate dsl/health-0.3.0.yaml
\end{lstlisting}
\begin{lstlisting}[style=out]
warning SIG_OPERATOR_TYPE            health-0.3.0.yaml#1:28 /adopted/left_operands/1/signature/operators/0 lteq on xsd:duration is only partially ordered
warning ACTION_INCLUDED_IN_MISSING   health-0.3.0.yaml#1:37 /terms/1 action reidentify has no included_in
ok: 1 documents valid
\end{lstlisting}

Los dos son deliberados y los dos merecen leerse.

El primero viene de adoptar \cod{odrl:elapsedTime} con \cod{xsd:duration} y \cod{lteq}. Las duraciones solo
están parcialmente ordenadas: no hay respuesta a si \cod{P1M} es menor o igual que \cod{P30D}, porque depende del
mes. Si el motor que evalúa la política va a comparar duraciones, conviene restringirse a días y segundos, o
cambiar a un entero con unidad fija, como hizo Flexgrid con \cod{retentionDays}.

El segundo viene de la acción \cod{reidentify}, que no tiene \cod{included_in}. Aquí es intencionado: es una
acción que solo aparecerá en prohibiciones, y no se quiere que un permiso general de \cod{odrl:use} la cubra.

\section{Una política contra este perfil}

\begin{lstlisting}[style=json]
{
  "@type": "Offer",
  "@id": "urn:uuid:offer-cohort-diabetes-01",
  "profile": "https://salud.example/odrl/profile/secondary-use",
  "target": "urn:uuid:dataset-primary-care-2019-2024",
  "permission": [{
    "action": "linkRecords",
    "constraint": [
      {"leftOperand": "purpose", "operator": "isAnyOf",
       "rightOperand": [{"@id": "hds:research"}, {"@id": "hds:publicHealth"}]},
      {"leftOperand": "dataPermit", "operator": "hds:withinPermit",
       "rightOperandReference": {"@id": "https://hdab.salud.example/permits/2026/0417"}},
      {"leftOperand": "kAnonymity", "operator": "gteq", "rightOperand": 10},
      {"leftOperand": "environment", "operator": "eq", "rightOperand": "secure-processing-environment"}
    ]
  }],
  "prohibition": [{ "action": "reidentify" }]
}
\end{lstlisting}

\section{Qué enseña la comparación}

\noindent\begin{tabularx}{\linewidth}{@{}L{4.2cm}Y Y@{}}
\toprule
\textbf{Decisión} & \textbf{Medida (Flexgrid)} & \textbf{Salud} \\
\midrule
Madurez & 1.0.0, vocabulario estable & 0.3.0, aún en exploración \\
Lista de operandos & Cerrada & Abierta \\
Enumeraciones & De IRI, con términos definidos & De IRI y de literales \\
Operadores & Solo los de ODRL & Uno propio \\
Referencias & No & Sí, para el permiso \\
Avisos del compilador & Ninguno & Dos, ambos asumidos \\
\bottomrule
\end{tabularx}

\begin{retos}
\begin{enumerate}[leftmargin=*, itemsep=2pt]
  \item Compila el perfil y busca en \cod{schema.json} cómo se valida \cod{environment}. ¿Qué dos formas admite
        el valor?
  \item Añade \texttt{included\_in: odrl:use} a \cod{reidentify} y piensa qué significaría para una oferta que
        permite \cod{odrl:use} y prohíbe \cod{reidentify}.
  \item Escribe un \cod{SpaceProfile} para este perfil. Recuerda que \texttt{"\textasciicircum{}0.3"} con versión mayor 0 fija la
        versión menor.
\end{enumerate}
\end{retos}


\part{El hub: alexandria-server}
\chapter{Arrancar el hub y gobernar un \emph{namespace}}
\label{ch:hub}

\cod{alexandria-server} es el hub. Tiene dos caras: una API de gestión autenticada bajo \cod{/v1} y un árbol de
publicación público e inmutable bajo \cod{/pub}. Este capítulo lo arranca, crea el \emph{namespace} de Flexgrid
y le asigna una \emph{steward}.

\section{El laboratorio}

Todo ocurre en \cod{tutorial/manual/examples/lab}. El fichero \cod{env.sh} define las variables y unas pocas
funciones que envuelven a \cod{curl}; ninguna esconde una llamada.

\begin{lstlisting}[style=sh]
cd tutorial/manual/examples/lab
source env.sh
\end{lstlisting}

\noindent\begin{tabularx}{\linewidth}{@{}L{3.3cm}L{5.1cm}Y@{}}
\toprule
\textbf{Variable} & \textbf{Valor} & \textbf{Qué es} \\
\midrule
\texttt{\$BIN} & \cod{../../../../bin} & Los binarios compilados. \\
\texttt{\$HUB} & \cod{https://127.0.0.1:18443} & El hub. \\
\texttt{\$IDP} & \cod{http://127.0.0.1:19400} & El emisor de tokens. \\
\texttt{\$CONNECTOR} & \cod{http://127.0.0.1:18282} & El conector de mentira. \\
\texttt{\$INGRESS} & \cod{http://127.0.0.1:19080} & La entrada de tráfico del cliente. \\
\texttt{\$ADMINAPI} & \cod{http://127.0.0.1:19090} & La API de administración del cliente. \\
\bottomrule
\end{tabularx}

\subsection{Un reloj que se puede mover}

El ciclo de vida de una versión dura semanas. Para recorrerlo en minutos, todos los procesos aceptan
\cod{-clock-file}: leen la hora de un fichero en lugar del reloj del sistema. La función \cod{clock} escribe en
él.

\begin{lstlisting}[style=sh]
clock 2026-09-01T09:00:00Z
\end{lstlisting}

\begin{trampaHabitual}[Todos con el mismo reloj]
El hub, el cliente y el emisor de tokens tienen que leer el \textbf{mismo} fichero. Un token emitido con la hora
real del sistema está, para un hub que cree vivir en septiembre, caducado o aún no vigente, y la respuesta es
\cod{401}. Por eso \cod{tok} pasa también \cod{--clock-file}.
\end{trampaHabitual}

\section{La configuración del hub}

\cod{lab/server.yaml} usa la disposición por secciones. Lo que lo distingue de \cod{config/server.yaml} son tres
cosas: los puertos, el compilador y el directorio de copias vendorizadas. Aquí se muestra condensado.

\begin{lstlisting}[style=yaml]
common_config:
  hosts:
    http: { protocol: "https", url: "127.0.0.1", port: "18443", internal_port: "18443" }
  db:
    db_type: "SQLite"
    database: "file:var/hub/hub.db"
    blobs: "file:var/hub/blobs"
  connection:
    tls: { cert: "file:var/tls/hub.crt", key: "file:var/tls/hub.key" }

auth_config:
  enabled: true
  issuer: "http://127.0.0.1:19400"
  audience: "alexandria-vhub"
  roles_claim: "roles"
  admin_role: "vhub-admin"

did_config:
  type: "key"
  signing_key: "file:var/hub/hub.key"

vhub_config:
  rdf_context: "self-hosted"
  compiler: "dsl"
  vendor_dir: "file:../vendor"
\end{lstlisting}

\noindent\begin{tabularx}{\linewidth}{@{}L{5.3cm}Y@{}}
\toprule
\textbf{Campo} & \textbf{Significado} \\
\midrule
\cod{common_config.hosts.http} & Dirección pública y puerto. De aquí sale la URL base, prefijo de toda IRI que emite el hub. \\
\cod{common_config.db} & Base de datos SQLite y directorio de blobs direccionados por contenido. \\
\cod{connection.tls} & Certificado y clave. Si no existen, se genera un certificado de laboratorio. \\
\cod{auth_config} & El hub es un servidor de recursos OAuth 2.0: emisor aceptado, audiencia y la reclamación que marca al administrador. \\
\cod{did_config.signing_key} & Clave Ed25519 con la que se firma. Si no existe, se genera. El DID del hub es por defecto el \cod{did:key} de esa clave. \\
\cod{vhub_config.compiler} & \cod{poc}, por defecto: borradores JSON. \cod{dsl}: borradores YAML del lenguaje de perfiles. \\
\cod{vhub_config.vendor_dir} & Copias locales para los \emph{bindings} vendorizados. \\
\cod{vhub_config.rdf_context} & \cod{self-hosted}: el hub sirve su propio contexto JSON-LD en \cod{/pub/_vhub/v1/context.jsonld}. \\
\bottomrule
\end{tabularx}

\medskip
Las rutas \cod{file:} se resuelven desde el directorio del fichero de configuración. El documento se valida al
arrancar contra \cod{deploy/server.schema.json}; con un error, el hub no arranca.

\begin{notaPractica}[Dos disposiciones, el mismo contenido]
Los dos binarios aceptan YAML o JSON, y dos disposiciones equivalentes: la de secciones, que usa este tutorial,
y la plana de la especificación (\cod{deploy/server.json}), donde los mismos campos se llaman \cod{listen},
\cod{base}, \cod{signing.key}, \cod{auth.issuers} o \cod{storage.database}. En la plana, \cod{compiler} y
\cod{vendor_dir} van en la raíz.
\end{notaPractica}

\section{Arrancar}

\begin{lstlisting}[style=sh]
$BIN/fake-idp serve --addr 127.0.0.1:19400 --issuer $IDP --key idp.key --clock-file clock > idp.log 2>&1 &
$BIN/fake-connector -listen 127.0.0.1:18282 > connector.log 2>&1 &
$BIN/alexandria-server -config server.yaml -clock-file clock > server.log 2>&1 &
\end{lstlisting}

En el primer arranque, el hub crea lo que le falta y lo dice:

\begin{lstlisting}[style=out]
time=2026-09-01T09:00:00.000Z level=INFO msg="generated lab TLS certificate" … cert_path=var/tls/hub.crt note="PoC: not for production"
time=2026-09-01T09:00:00.000Z level=INFO msg="signing key verified in hub DID document" … hub_did=did:key:z6Mkkfgq…
time=2026-09-01T09:00:00.000Z level=INFO msg="server listening" … listen=127.0.0.1:18443 base=https://127.0.0.1:18443 tls_enabled=true
\end{lstlisting}

Tres cosas a notar. Las marcas de tiempo son las del reloj de fichero. El hub comprueba que su clave de firma
es un método de aserción de su propio documento DID antes de aceptar tráfico: nunca firma con una clave que no
haya publicado. Y el \cod{hub_did} hay que apuntarlo, porque el cliente solo aceptará credenciales de ese DID.

\begin{lstlisting}[style=sh]
hubdid
\end{lstlisting}
\begin{lstlisting}[style=out]
did:key:z6MkkfgqPKJK19zucGBNfLJQHUFwEA3mc9skvYsPDbkoNP8e
\end{lstlisting}

El tuyo será distinto: la clave se genera en cada laboratorio nuevo.

\section{Quién puede hacer qué}

\subsection{Tokens}

\texttt{fake-idp token} firma un token de acceso. Con \cod{--admin} añade el rol de administrador del hub.

\begin{lstlisting}[style=sh]
tok --sub ana --admin > admin.token
tok --sub lucia > steward.token
\end{lstlisting}

Ana es la administradora del hub. Lucía, de la distribuidora, será la \emph{steward} de Flexgrid. Por ahora su
token solo dice quién es; el hub decidirá qué puede hacer.

El hub valida cada token en cinco pasos, y el primero que falla responde \cod{401}:

\noindent\begin{tabularx}{\linewidth}{@{}C{1.1cm}Y@{}}
\toprule
\textbf{Paso} & \textbf{Comprobación} \\
\midrule
A1 & Es un JWT de tipo \cod{at+jwt} firmado con un algoritmo asimétrico permitido. \\
A2 & Su emisor es uno de los configurados. \\
A3 & La firma verifica contra las claves públicas del emisor, que se descubren por OIDC. \\
A4 & La audiencia incluye la del hub. \\
A5 & Está dentro de su periodo de validez, con la tolerancia de reloj. \\
\bottomrule
\end{tabularx}

\subsection{La tabla de autorización}

\noindent\begin{tabularx}{\linewidth}{@{}Y C{1.6cm} C{2.3cm} C{1.5cm} C{1.2cm}@{}}
\toprule
\textbf{Operación} & \textbf{Anónimo} & \textbf{Autenticado} & \textbf{Steward} & \textbf{Admin} \\
\midrule
Leer \emph{namespaces}, vocabularios, listado y registros; todo \cod{/pub} & \si & \si & \si & \si \\
Crear un \emph{namespace}; asignar o quitar \emph{stewards}; leer el registro de eventos global & 401 & 403 & 403 & \si \\
Listar los \emph{stewards} de un \emph{namespace} & 401 & 403 & \si & \si \\
Vocabularios, borradores, candidatos, aprobación, revocación y eventos de un \emph{namespace} & 401 & 403 & \si & 403 \\
\bottomrule
\end{tabularx}

\begin{trampaHabitual}[El administrador no gobierna el vocabulario]
La última fila sorprende: el administrador recibe \cod{403} en las operaciones de un \emph{namespace}. Ser
administrador del hub no da poder sobre el contenido de un espacio. Si Ana quisiera publicar versiones de
Flexgrid, tendría que asignarse como \emph{steward}.
\end{trampaHabitual}

Las respuestas reales de los dos casos de error:

\begin{lstlisting}[style=sh]
curl -sk -X POST $HUB/v1/namespaces -d '{"slug":"x"}'
\end{lstlisting}
\begin{lstlisting}[style=out]
{"@context":"https://127.0.0.1:18443/pub/_vhub/v1/context.jsonld",
 "type":"https://w3id.org/alexandria-vhub/problems#API_UNAUTHENTICATED","title":"API_UNAUTHENTICATED",
 "status":401,"detail":"this operation requires a bearer token","code":"API_UNAUTHENTICATED","severity":"error"}
\end{lstlisting}

Un administrador intentando crear un borrador devuelve la misma forma con \cod{API_FORBIDDEN_ROLE} y \cod{403}.
Toda respuesta de \cod{/v1}, errores incluidos, es JSON-LD compacto con un \cod{code} del catálogo cerrado.

\section{Crear el \emph{namespace} y su \emph{steward}}

\begin{lstlisting}[style=sh]
curl -sk -X POST $HUB/v1/namespaces -H "Authorization: Bearer $(cat admin.token)" \
  -d '{"slug":"flexgrid","title":"Flexgrid Energy Data Space"}'
\end{lstlisting}
\begin{lstlisting}[style=out]
{"@context":"https://127.0.0.1:18443/pub/_vhub/v1/context.jsonld","created_at":"2026-09-01T09:00:00Z",
 "id":"https://127.0.0.1:18443/v1/namespaces/flexgrid","slug":"flexgrid",
 "title":"Flexgrid Energy Data Space","type":"Namespace"}
\end{lstlisting}

El \cod{slug} es único e inmutable y forma parte de todas las URL. \cod{_vhub} está reservado.

\begin{lstlisting}[style=sh]
curl -sk -X POST $HUB/v1/namespaces/flexgrid/stewards -H "Authorization: Bearer $(cat admin.token)" \
  -d '{"iss":"http://127.0.0.1:19400","sub":"lucia","label":"Lucía (DSO)"}'
\end{lstlisting}
\begin{lstlisting}[style=out]
{"@context":"…","assigned_at":"2026-09-01T09:00:00Z","assigned_by":"http://127.0.0.1:19400#ana",
 "id":"https://127.0.0.1:18443/v1/namespaces/flexgrid/stewards/ed6edbdc291414db",
 "iss":"http://127.0.0.1:19400","label":"Lucía (DSO)","sub":"lucia","type":"Steward"}
\end{lstlisting}

Una \emph{steward} se identifica por el par emisor y sujeto, no por un nombre. La pertenencia se evalúa en cada
petición: el token de Lucía no ha cambiado, pero desde este instante el hub le deja operar Flexgrid.

\section{El registro de vocabularios}

Además de versiones, un \emph{namespace} puede alojar documentos sueltos con revisiones inmutables: una lista de
códigos, un contexto de terceros fijado, una ontología auxiliar. Hay dos tipos.

\noindent\begin{tabularx}{\linewidth}{@{}L{2cm}Y@{}}
\toprule
\textbf{Tipo} & \textbf{Comportamiento} \\
\midrule
\cod{hosted} & El contenido se sube. Cada subida distinta crea una revisión nueva; subir los mismos bytes no crea nada. \\
\cod{proxy} & El hub descarga un documento externo al crear el vocabulario y al pedir \cod{:refresh}. Nunca contacta con el origen al servir \cod{/pub}. \\
\bottomrule
\end{tabularx}

\medskip
Flexgrid aloja su esquema de regiones, \cod{examples/dsl/regions.ttl}:

\begin{lstlisting}[style=sh]
S=$(cat steward.token)
curl -sk -X POST $HUB/v1/namespaces/flexgrid/vocabularies -H "Authorization: Bearer $S" \
  -d '{"name":"regions","kind":"hosted","media_type":"text/turtle"}'

curl -sk -i -X PUT $HUB/v1/namespaces/flexgrid/vocabularies/regions/content \
  -H "Authorization: Bearer $S" -H 'Content-Type: text/turtle' --data-binary @../dsl/regions.ttl
\end{lstlisting}
\begin{lstlisting}[style=out]
HTTP/2 201
{"@context":"…","created_at":"…","created_by":"http://127.0.0.1:19400#lucia",
 "digest":"ni:///sha-256;zmzyAQLYqxTjnY6aXPy_HRSkqmgxrO7MhJsJkU6ohOI",
 "id":"https://127.0.0.1:18443/pub/flexgrid/vocab/regions/r/1","n":1,"size":432,"type":"Revision"}
\end{lstlisting}

Repetir la misma subida devuelve \cod{200} y no crea revisión. Subir algo que no es Turtle se rechaza:

\begin{lstlisting}[style=out]
{"…","title":"VOCAB_UNPARSEABLE","status":422,"detail":"content is not valid text/turtle","code":"VOCAB_UNPARSEABLE","severity":"error"}
\end{lstlisting}

La revisión es pública de inmediato. La URL sin revisión redirige a la última; la URL con revisión es inmutable:

\begin{lstlisting}[style=sh]
curl -sk -i $HUB/pub/flexgrid/vocab/regions
curl -sk -i $HUB/pub/flexgrid/vocab/regions/r/1
\end{lstlisting}
\begin{lstlisting}[style=out]
HTTP/2 303
cache-control: no-cache
location: /pub/flexgrid/vocab/regions/r/1

HTTP/2 200
cache-control: public, max-age=31536000, immutable
content-type: text/turtle
etag: "ce6cf20102d8ab14e39d8e9a5cfcbf1d14a4aa6831aceecc849b09914ea884e2"
\end{lstlisting}

\begin{ideaClave}
Este patrón se repite en todo el hub: una URL \textbf{estable} que redirige y no se cachea, y una URL
\textbf{inmutable} que se cachea para siempre y cuyo \cod{ETag} es el SHA-256 del contenido. Quien quiera «lo
último» usa la primera; quien necesite reproducibilidad cita la segunda.
\end{ideaClave}

\begin{retos}
\begin{enumerate}[leftmargin=*, itemsep=2pt]
  \item Emite un token para un tercero, \texttt{tok -{}-sub eve}, e intenta crear un vocabulario. ¿Qué código devuelve?
  \item Sube una segunda versión de \cod{regions.ttl} con una región más y comprueba que
        \cod{/pub/flexgrid/vocab/regions} redirige ahora a \cod{r/2} mientras que \cod{r/1} sigue igual.
  \item Intenta crear el \emph{namespace} \cod{flexgrid} otra vez, y luego uno llamado \cod{Flexgrid!}.
\end{enumerate}
\end{retos}

\chapter{Publicar una versión}
\label{ch:release}

Publicar no es subir un fichero. Es un proceso con tres estados, y cada paso deja rastro.

\begin{center}
\begin{tikzpicture}[font=\small, node distance=16mm,
  st/.style={draw=alxNavy, fill=alxLightBlue, rounded corners=2pt, minimum height=10mm, minimum width=27mm, align=center},
  arr/.style={-{Stealth[length=2.2mm]}, alxNavy, thick}]
  \node[st] (d) {Borrador\\[-2pt]\scriptsize editable};
  \node[st, right=of d] (c) {Candidato\\[-2pt]\scriptsize inmutable};
  \node[st, right=of c, fill=alxLightTeal, draw=alxTeal] (v) {Versión\\[-2pt]\scriptsize firmada};
  \node[st, right=of v, fill=alxLightRed, draw=alxRed] (r) {Revocada};
  \draw[arr] (d) -- node[above]{\scriptsize\texttt{:propose}} (c);
  \draw[arr] (c) -- node[above]{\scriptsize\texttt{:approve}} (v);
  \draw[arr] (v) -- node[above]{\scriptsize\texttt{:revoke}} (r);
  \draw[arr] (d) edge[loop above] node{\scriptsize\texttt{PUT}, \texttt{:validate}} (d);
  \draw[arr] (c) -- ++(0,-11mm) node[below]{\scriptsize\texttt{:reject}};
\end{tikzpicture}
\end{center}

\section{El borrador}

Con \texttt{compiler: dsl}, un borrador es un \emph{bundle} YAML, exactamente el fichero que \cod{alexandriac} ya
ha validado. Se avanza el reloj una hora y se sube:

\begin{lstlisting}[style=sh]
clock 2026-09-01T10:00:00Z
curl -sk -i -X POST $HUB/v1/namespaces/flexgrid/drafts -H "Authorization: Bearer $(cat steward.token)" \
  -H 'Content-Type: application/yaml' --data-binary @../dsl/flexgrid-1.0.0.yaml
\end{lstlisting}
\begin{lstlisting}[style=out]
HTTP/2 201
content-type: application/ld+json
etag: "05511c6385475d68-1"
location: https://127.0.0.1:18443/v1/namespaces/flexgrid/drafts/05511c6385475d68

{"@context":"…","id":"https://127.0.0.1:18443/v1/namespaces/flexgrid/drafts/05511c6385475d68",
 "namespace":"https://127.0.0.1:18443/v1/namespaces/flexgrid","rev":1}
\end{lstlisting}

El borrador es lo único editable del proceso. Para reemplazarlo se usa \cod{PUT} con la cabecera \cod{If-Match}
y el \cod{ETag} actual. Sin ella, el hub responde \cod{428}; con una desfasada, \cod{412}. Así dos
\emph{stewards} no se pisan el trabajo.

\subsection{Validar sin efectos}

\begin{lstlisting}[style=sh]
DRAFT=$HUB/v1/namespaces/flexgrid/drafts/05511c6385475d68
curl -sk -X POST "$DRAFT:validate" -H "Authorization: Bearer $(cat steward.token)"
\end{lstlisting}
\begin{lstlisting}[style=out]
{"@context":"…","class":"initial","findings":[],"valid":true}
\end{lstlisting}

\cod{:validate} ejecuta los cinco niveles y no cambia nada. Devuelve la clase que tendría la versión y la lista
de hallazgos. Cuando algo está mal, los hallazgos son los mismos que daría \cod{alexandriac}, con su código, su
puntero y su línea. Por ejemplo, un borrador 2.1.0 que quita una región con un salto menor:

\begin{lstlisting}[style=out]
{ "class": "",
  "findings": [
    { "Code": "VERSION_BUMP_TOO_LOW",
      "Detail": "draft.yaml#1: 2.0.0 → 2.1.0 is a MINOR bump; the changes require MAJOR: [operator added to enr:minAggregation value removed from enr:processingRegion]",
      "Path": "/profile/version", "Line": 6 } ],
  "valid": false }
\end{lstlisting}

Intentar proponerlo devuelve \cod{422} con el mismo código. Nótese que el hub conocía la versión anterior sin
que nadie se la diera: la tomó de \cod{sources/draft.yaml} del paquete vigente.

\begin{notaPractica}[Las claves de los hallazgos van en mayúscula]
En \cod{findings}, los campos se llaman \cod{Code}, \cod{Detail}, \cod{Path} y \cod{Line}, con mayúscula
inicial, mientras que el resto de la API usa minúsculas (\cod{code}, \cod{detail}). Es una asimetría de la
implementación actual; un \emph{script} que lea ambos tiene que contar con ella.
\end{notaPractica}

\section{El candidato}

\begin{lstlisting}[style=sh]
curl -sk -X POST "$DRAFT:propose" -H "Authorization: Bearer $(cat steward.token)"
\end{lstlisting}
\begin{lstlisting}[style=out]
{
 "class": "initial",
 "diff": [
  "added profile/slots.json",
  "added profiles/dsp/vendored/2025-1/context.jsonld",
  "added profiles/dsp/vendored/2025-1/dsp-messages.schema.json",
  "added profiles/odrl/metering/1.0.0/context.jsonld",
  "added profiles/odrl/metering/1.0.0/prof.ttl",
  "added profiles/odrl/metering/1.0.0/schema.json",
  "added profiles/odrl/metering/1.0.0/shapes.ttl",
  "added profiles/odrl/metering/1.0.0/spec.html",
  "added profiles/odrl/metering/1.0.0/vocab.ttl",
  "added sources/draft.yaml",
  "added space/prof.ttl",
  "added space/shacl-union.ttl",
  "added space/spec.html"
 ],
 "id": "https://127.0.0.1:18443/v1/namespaces/flexgrid/candidates/5e6789e09b415f9d",
 "lifecycle": {
  "freshness": { "ttl": "P1D", "max_staleness": "P7D" },
  "min_grace_window": "P7D", "grace_window": "P30D", "sync_interval": "PT6H",
  "clock_skew": "PT5S", "lease_margin": "PT10S", "propagation_window": "P2D",
  "max_concurrent_versions": 4,
  "degraded": { "dsp_message": "continue", "policy_intake": "reject" }
 },
 "planned": {
  "version": "1.0.0", "class": "initial",
  "published_at": "2026-09-01T10:00:00Z",
  "new_content_from": "2026-09-01T10:00:00Z",
  "enforced_from": "2026-09-01T10:00:00Z",
  "revoked_from": null
 },
 "proposed_at": "2026-09-01T10:00:00Z",
 "state": "open",
 "version": "1.0.0"
}
\end{lstlisting}

Al proponer, el hub compila el borrador y congela el resultado. El candidato es lo que una \emph{steward} revisa
antes de aprobar, y contiene cuatro cosas:

\begin{itemize}[leftmargin=*, itemsep=2pt]
  \item \cod{class}: la clase de la versión.
  \item \cod{diff}: qué ficheros se añaden, cambian o desaparecen respecto de la versión vigente.
  \item \cod{lifecycle}: los parámetros efectivos, con los valores por defecto del hub ya rellenados.
  \item \cod{planned}: las fechas que tendría si se aprobara en este instante.
\end{itemize}

Editar el borrador después no cambia el candidato. Un candidato solo sale de \cod{open} aprobado o rechazado.

\section{Aprobar}

\begin{lstlisting}[style=sh]
curl -sk -X POST "$HUB/v1/namespaces/flexgrid/candidates/5e6789e09b415f9d:approve" \
  -H "Authorization: Bearer $(cat steward.token)" -d '{}'
\end{lstlisting}
\begin{lstlisting}[style=out]
{
  "@context": ["https://www.w3.org/ns/credentials/v2",
               "https://127.0.0.1:18443/pub/_vhub/v1/context.jsonld"],
  "type": ["VerifiableCredential", "VersionRecord"],
  "issuer": "did:key:z6MkkfgqPKJK19zucGBNfLJQHUFwEA3mc9skvYsPDbkoNP8e",
  "validFrom": "2026-09-01T10:00:00Z",
  "credentialSubject": {
    "id": "https://127.0.0.1:18443/v1/namespaces/flexgrid/versions/1.0.0",
    "namespace": "https://127.0.0.1:18443/v1/namespaces/flexgrid",
    "version": "1.0.0", "class": "initial",
    "manifest": "ni:///sha-256;cCYTVfOvidDw_pxQY5o1DHnHmfRjcSzDDAp3xLmxw2k",
    "published_at": "2026-09-01T10:00:00Z",
    "new_content_from": "2026-09-01T10:00:00Z",
    "enforced_from": "2026-09-01T10:00:00Z",
    "revoked_from": null
  },
  "proof": {
    "type": "DataIntegrityProof", "cryptosuite": "eddsa-jcs-2022",
    "proofPurpose": "assertionMethod", "created": "2026-09-01T10:00:00Z",
    "verificationMethod": "did:key:z6Mkkfgq…#z6Mkkfgq…",
    "proofValue": "z5AkwWzBn1jYoaQJZocbagUD1fs7xWdhSLL9RiBXh6ptx39QiPy57dfeux1kJmEGeHFbnkVx1i1kdb1A63WC3VDb8"
  }
}
\end{lstlisting}

La respuesta es el \textbf{registro de versión}: una credencial verificable W3C~VC~2.0 firmada por el hub con
\cod{eddsa-jcs-2022}. Quien la tenga puede comprobar, sin preguntar a nadie, que el hub dijo que la 1.0.0 de
Flexgrid tiene ese manifiesto y esas fechas.

\begin{notaPractica}[Con las funciones del laboratorio]
Los tres pasos anteriores están envueltos en \cod{publish} y \cod{approve} de \cod{env.sh}. A partir de aquí el
tutorial los usa:

\begin{lstlisting}[style=sh]
publish ../dsl/flexgrid-1.0.0.yaml
approve '{}'
\end{lstlisting}
\end{notaPractica}

\section{Las tres credenciales}

\noindent\begin{tabularx}{\linewidth}{@{}L{3.1cm}L{5.4cm}Y@{}}
\toprule
\textbf{Credencial} & \textbf{Dónde} & \textbf{Qué afirma} \\
\midrule
\cod{VersionListing} & \texttt{/\allowbreak{}v1/\allowbreak{}namespaces/\allowbreak{}\{ns\}/\allowbreak{}versions} & Qué versiones existen ahora, con un número de secuencia y una caducidad. \\
\cod{VersionRecord} & \texttt{/\allowbreak{}v1/\allowbreak{}namespaces/\allowbreak{}\{ns\}/\allowbreak{}versions/\allowbreak{}\{v\}} & Las fechas y el manifiesto de una versión. \\
\cod{PackageManifest} & \texttt{/\allowbreak{}pub/\allowbreak{}packages/\allowbreak{}\{resumen\}/\allowbreak{}manifest.json} & Qué ficheros componen el paquete, con su resumen. \\
\bottomrule
\end{tabularx}

\medskip
El \textbf{listado} es la pieza que el cliente sincroniza:

\begin{lstlisting}[style=sh]
curl -sk $HUB/v1/namespaces/flexgrid/versions
\end{lstlisting}
\begin{lstlisting}[style=out]
{ "type": ["VerifiableCredential", "VersionListing"],
  "issuer": "did:key:z6Mkkfgq…",
  "validFrom": "2026-09-01T10:00:00Z",
  "validUntil": "2026-09-02T10:00:00Z",
  "credentialSubject": {
    "id": "https://127.0.0.1:18443/v1/namespaces/flexgrid/versions",
    "sequence": 1,
    "versions": [ { "version": "1.0.0", "class": "initial", "manifest": "ni:///sha-256;cCYT…",
                    "new_content_from": "2026-09-01T10:00:00Z", "enforced_from": "2026-09-01T10:00:00Z",
                    "revoked_from": null, … } ] },
  "proof": { … } }
\end{lstlisting}

Dos campos sostienen la seguridad del conjunto. \cod{sequence} crece con cada emisión: un cliente que reciba una
secuencia menor que la última aceptada sabe que le están sirviendo algo viejo. \cod{validUntil} es la caducidad:
el hub reemite el listado al menos cada medio \cod{ttl}, y un cliente con un listado caducado sabe que ha perdido
el contacto.

\section{Dentro del paquete}

\begin{lstlisting}[style=sh]
M=$(curl -sk $HUB/v1/namespaces/flexgrid/versions/1.0.0 |
    python3 -c "import json,sys;print(json.load(sys.stdin)['credentialSubject']['manifest'])")
curl -sk "$HUB/pub/packages/${M#ni:///sha-256;}/manifest.json"
\end{lstlisting}

El manifiesto lista los ficheros del paquete. Resumido en una tabla:

\begin{longtable}{@{}L{2.5cm}L{1.1cm}L{2.9cm}L{8.3cm}@{}}
\toprule
\textbf{Rol} & \textbf{Slot} & \textbf{Gates} & \textbf{Ruta} \\
\midrule
\endhead
\cod{json-schema} & \cod{odrl} & \cod{policy_intake} & \cod{profiles/odrl/metering/1.0.0/schema.json} \\
\cod{shacl} & \cod{odrl} & \cod{policy_intake} & \cod{profiles/odrl/metering/1.0.0/shapes.ttl} \\
\cod{other} & \cod{odrl} & & \cod{profiles/odrl/metering/1.0.0/context.jsonld} \newline con \cod{url}: \cod{https://flexgrid.example/odrl/profile/metering/context.jsonld} \\
\cod{other} & \cod{odrl} & & \cod{vocab.ttl}, \cod{prof.ttl}, \cod{spec.html} \\
\cod{json-schema} & \cod{dsp} & \cod{dsp_message} & \cod{profiles/dsp/vendored/2025-1/dsp-messages.schema.json} \\
\cod{context} & \cod{dsp} & & \cod{profiles/dsp/vendored/2025-1/context.jsonld} \newline con \cod{url}: \cod{https://w3id.org/dspace/2025/1/context.jsonld} \\
\cod{other} & & & \cod{space/prof.ttl}, \cod{space/shacl-union.ttl}, \cod{space/spec.html} \\
\cod{other} & & & \cod{profile/slots.json} \\
\cod{other} & & & \cod{sources/draft.yaml} \\
\bottomrule
\end{longtable}

Tres observaciones:

\begin{itemize}[leftmargin=*, itemsep=2pt]
  \item El manifiesto solo conoce cuatro roles: \cod{json-schema}, \cod{shacl}, \cod{context} y \cod{other}.
        Los roles más finos del lenguaje, como \cod{vocabulary}, \cod{mapping} o \cod{prof}, se publican como
        \cod{other}.
  \item \cod{profile/slots.json} lleva la obligación y los gates de cada slot. Es lo que el cliente lee para
        saber qué evaluar.
  \item \cod{sources/draft.yaml} es el borrador tal cual. Es la «versión anterior» con la que el hub comparará
        el siguiente.
\end{itemize}

Cada fichero se sirve inmutable bajo la versión:

\begin{lstlisting}[style=sh]
curl -sk -i $HUB/pub/flexgrid/1.0.0/profiles/odrl/metering/1.0.0/schema.json
\end{lstlisting}
\begin{lstlisting}[style=out]
HTTP/2 200
cache-control: public, max-age=31536000, immutable
content-type: application/schema+json
etag: "fa70ada7f8c3e45426e804c1883da0a338d3e718425211ad86626d702cdc5e09"
\end{lstlisting}

El \cod{ETag} es el mismo SHA-256 que aparecía en el \cod{space-lock.json} que generó \cod{alexandriac} en el
capítulo~\ref{ch:spaceprofile}. Los bytes que compilaste en local y los que publica el hub son los mismos.

\noindent\begin{tabularx}{\linewidth}{@{}L{6.4cm}Y@{}}
\toprule
\textbf{Ruta de \texttt{/\allowbreak{}pub}} & \textbf{Respuesta} \\
\midrule
\texttt{/\allowbreak{}pub/\allowbreak{}\{ns\}} & \cod{303} al primer fichero de rol \cod{context} de la versión cabeza. \cod{404} mientras ninguna versión esté vigente. \\
\texttt{/\allowbreak{}pub/\allowbreak{}\{ns\}/\allowbreak{}\{versión\}/\allowbreak{}\{ruta\}} & Un fichero de un paquete. Inmutable. \\
\texttt{/\allowbreak{}pub/\allowbreak{}\{ns\}/\allowbreak{}vocab/\allowbreak{}\{nombre\}} & \cod{303} a la última revisión. \\
\texttt{/\allowbreak{}pub/\allowbreak{}\{ns\}/\allowbreak{}vocab/\allowbreak{}\{nombre\}/\allowbreak{}r/\allowbreak{}\{n\}} & Una revisión. Inmutable. \\
\texttt{/\allowbreak{}pub/\allowbreak{}packages/\allowbreak{}\{resumen\}/\allowbreak{}manifest.json} & Manifiesto firmado. Inmutable. \\
\texttt{/\allowbreak{}pub/\allowbreak{}blobs/\allowbreak{}sha256/\allowbreak{}\{hex\}} & Un blob por su resumen. Inmutable. \\
\cod{/pub/_vhub/v1/context.jsonld} & El contexto JSON-LD del propio hub. \\
\bottomrule
\end{tabularx}

\begin{trampaHabitual}[A dónde redirige \texttt{/\allowbreak{}pub/\allowbreak{}flexgrid}]
Con una versión vigente, \cod{/pub/flexgrid} redirige a
\cod{/pub/flexgrid/2.0.0/profiles/dsp/vendored/2025-1/context.jsonld}: el contexto de DSP, no el del perfil. El
motivo es que el contexto del perfil se publica con rol \cod{other} (era \cod{mapping}), y el único fichero con
rol \cod{context} es el vendorizado. Si esperabas el contexto de tu perfil en la URL del \emph{namespace}, no
llega.
\end{trampaHabitual}

\section{Las fechas}
\label{sec:fechas}

Una versión tiene tres instantes. \cod{published_at} es cuando se aprueba. \cod{new_content_from} es cuando su
contenido empieza a aceptarse. \cod{enforced_from} es cuando pasa a ser la única que vale.

Con $P$ el instante de publicación, $pw$ la ventana de propagación y $gw$ la de gracia:

\noindent\begin{tabularx}{\linewidth}{@{}L{2.3cm}L{3.9cm}L{3.9cm}Y@{}}
\toprule
\textbf{Clase} & \textbf{\texttt{grace\_window} al aprobar} & \textbf{\texttt{new\_content\_from}} & \textbf{\texttt{enforced\_from}} \\
\midrule
\cod{initial} & No se admite & $P$ & $P$ \\
\cod{compatible} & No se admite & $P + pw$ & $P + pw$ \\
\cod{breaking} & Omitida & $P + pw$ & $P + pw + gw$ \\
\cod{breaking} & Igual o mayor que la mínima & $P + pw$ & $P + pw + $ la pedida \\
\cod{breaking} & Menor que la mínima & \multicolumn{2}{l}{Rechazada: \cod{DATES_GRACE_TOO_SHORT}} \\
\bottomrule
\end{tabularx}

\medskip
Una versión compatible hereda la fecha de entrada en vigor de su predecesora cuando esa fecha es posterior: no
puede adelantar a una versión que todavía está en su ventana.

Para la 1.0.0 (\cod{initial}), aprobada el 1 de septiembre a las 10:00, tanto \cod{new_content_from} como
\cod{enforced_from} valen ese mismo instante: al no existir ninguna versión previa en el \emph{namespace}, entra
en vigor desde el principio. En cambio, para las transiciones posteriores (\cod{compatible} y \cod{breaking}) se
suma la ventana de propagación ($pw$): nadie se entera de una regla nueva sobre un espacio en marcha en el mismo
instante en que empieza a aplicarse.

El capítulo~\ref{ch:migracion} publica una versión compatible y una de ruptura y enseña cómo se ven esas fechas
desde el cliente.

\section{La auditoría}

Toda operación con efecto añade un evento con un identificador creciente. Las tablas que los guardan no
admiten modificación ni borrado: lo impiden disparadores de la base de datos.

\begin{lstlisting}[style=sh]
curl -sk "$HUB/v1/namespaces/flexgrid/events?limit=100" -H "Authorization: Bearer $(cat steward.token)"
\end{lstlisting}

Resumido, tras publicar la 1.0.0, con el identificador, el instante, el tipo, quién lo hizo y sobre qué
(el \emph{principal} completo es \texttt{http://127.0.0.1:19400\#lucia}):

\begin{lstlisting}[style=out]
3  2026-09-01T09:00:00Z  draft.created       lucia   05511c6385475d68
4  2026-09-01T10:00:00Z  candidate.proposed  lucia   5e6789e09b415f9d
5  2026-09-01T10:00:00Z  candidate.approved  lucia   5e6789e09b415f9d
6  2026-09-01T10:00:00Z  release.published   lucia   …/versions/1.0.0
7  2026-09-01T10:00:00Z  listing.issued      vhub    …/versions
\end{lstlisting}

Hay quince tipos de evento, entre ellos \cod{namespace.created}, \cod{steward.assigned},
\cod{vocabulary.revised}, \cod{candidate.rejected}, \cod{release.aborted} y \cod{version.revoked}. El de un
\emph{namespace} lo leen sus \emph{stewards}; el global, \cod{/v1/events}, la administradora.

\begin{retos}
\begin{enumerate}[leftmargin=*, itemsep=2pt]
  \item Crea un borrador, reemplázalo con \cod{PUT} sin \cod{If-Match} y luego con un \cod{ETag} inventado.
  \item Propón un borrador y recházalo con \cod{:reject}. ¿Qué pasa si después intentas aprobarlo?
  \item Verifica a mano que el \cod{ETag} de un fichero de \cod{/pub} es su SHA-256 con \texttt{shasum -a 256}.
\end{enumerate}
\end{retos}


\part{El borde: alexandria-client}
\chapter{Arrancar el cliente y entender su estado}
\label{ch:cliente}

\cod{alexandria-client} se coloca delante del conector de un participante. Hace dos cosas, y conviene verlas
por separado: \textbf{sincroniza} las versiones del \emph{namespace}, verificando cada firma, y \textbf{decide}
sobre cada mensaje con lo que tiene en local. La segunda nunca depende de que el hub responda.

\section{La configuración}

\begin{lstlisting}[style=yaml]
common_config:
  connection:
    tls:
      ca: "file:var/tls/hub.crt"

did_config:
  type: "key"
  hub_did: "HUB_DID_HERE"
  did_ttl: "PT6H"

observability:
  events_sink: "file:var/client/events.jsonl"
  audit_window: 100

client_config:
  trust:
    hub: "https://127.0.0.1:18443"
    namespace: "flexgrid"
    context: "https://127.0.0.1:18443/pub/_vhub/v1/context.jsonld"
  upstreams:
    dsp: "http://127.0.0.1:18282/dsp"
  listeners:
    dsp-ingress: { bind: "127.0.0.1:19080" }
    dsp-egress:  { bind: "127.0.0.1:19081" }
    admin:       { bind: "127.0.0.1:19090" }
  rules:
    seed: "file:../../../../deploy/rules.seed.json"
  default_action: "forward"
  unknown_agreement: "forward"
\end{lstlisting}

\noindent\begin{tabularx}{\linewidth}{@{}L{4.4cm}Y@{}}
\toprule
\textbf{Campo} & \textbf{Significado} \\
\midrule
\cod{trust.hub} & URL del hub. \\
\cod{trust.namespace} & El \emph{namespace} que sigue. Un cliente sigue exactamente uno. \\
\cod{did_config.hub_did} & El DID cuyas claves se aceptan para firmar listados y manifiestos. Es el ancla de confianza. \\
\cod{trust.context} & Contexto JSON-LD esperado en las credenciales. Con un hub \cod{self-hosted}, la URL del propio hub. \\
\cod{connection.tls.ca} & Raíz TLS del hub. Solo en laboratorio, donde el certificado es autofirmado. \\
\cod{upstreams.dsp} & El punto DSP del conector que se protege. \\
\cod{listeners} & Dónde escucha cada tipo de tráfico. \\
\cod{rules.seed} & Reglas iniciales. Solo se cargan si el almacén de reglas está vacío. \\
\cod{default_action} & Qué hacer con una petición que ninguna regla cubre: \cod{forward} o \cod{reject}. \\
\cod{unknown_agreement} & Qué hacer con una transferencia cuyo acuerdo nunca se capturó. \\
\cod{events_sink} & Fichero de eventos, una línea JSON por evento. \\
\bottomrule
\end{tabularx}

\begin{ideaClave}
Toda la confianza del cliente cuelga de un solo dato: \cod{hub_did}. El cliente no confía en el servidor que le
responde, ni en el certificado TLS, ni en la URL. Confía en que el listado y los manifiestos estén firmados por
una clave de ese DID. Un hub suplantado no puede hacerle aceptar nada.
\end{ideaClave}

\subsection{Los cuatro \emph{listeners}}

\noindent\begin{tabularx}{\linewidth}{@{}L{2.8cm}L{2.8cm}Y@{}}
\toprule
\textbf{\emph{Listener}} & \textbf{Tipo} & \textbf{Tráfico} \\
\midrule
\cod{dsp-ingress} & Proxy inverso & Mensajes DSP entrantes, hacia \cod{upstreams.dsp}. \\
\cod{dsp-egress} & Proxy directo & Mensajes DSP salientes del conector hacia otros participantes. \\
\cod{mgmt-ingress} & Proxy inverso & La API de gestión del conector. Opcional. \\
\cod{admin} & API & Estado, reglas, acuerdos, auditoría, salud y métricas. Nunca se monta en un \emph{listener} de tráfico. \\
\bottomrule
\end{tabularx}

\medskip
Peticiones y respuestas se reenvían byte a byte. Lo único que el cliente añade jamás es la cabecera
\cod{VHub-Verdict}.

\section{Arrancar}

Primero se escribe el DID del hub en la configuración, y después se arranca con el mismo reloj:

\begin{lstlisting}[style=sh]
sed "s/HUB_DID_HERE/$(hubdid)/" client.yaml > client.local.yaml
$BIN/alexandria-client -config client.local.yaml -state-dir var/client/state \
  -clock-file clock > client.log 2>&1 &
\end{lstlisting}

\noindent\begin{tabularx}{\linewidth}{@{}L{4.4cm}Y@{}}
\toprule
\textbf{Opción} & \textbf{Para qué} \\
\midrule
\texttt{-config <ruta>} & Fichero de configuración. \\
\texttt{-state-dir <dir>} & Estado persistente: continuidad, paquetes y acuerdos. \\
\texttt{-clock-file <ruta>} & Reloj de fichero. Solo laboratorio. \\
\texttt{-hub-ca <ruta>} & Alternativa a \cod{tls.ca} para dar la raíz TLS del hub. \\
\cod{-log-level}, \cod{-log-format} & Sobrescriben la configuración de registro. \\
\bottomrule
\end{tabularx}

\medskip
El registro cuenta lo que ha pasado en el primer segundo:

\begin{lstlisting}[style=out]
level=INFO msg="ruleset initialized" … component=rules revision=1 rules=16 seeded=true
level=INFO msg="listener started" … listener=dsp-ingress bind=127.0.0.1:19080
level=INFO msg="listener started" … listener=admin bind=127.0.0.1:19090
level=INFO msg="listener started" … listener=dsp-egress bind=127.0.0.1:19081
level=INFO msg="hub listing accepted" … component=sync sequence=1 versions=1 fresh_packages=1
\end{lstlisting}

Ha cargado las 16 reglas de la semilla, ha abierto sus \emph{listeners} y ha aceptado el listado de secuencia 1,
descargando un paquete.

\section{El estado}

\begin{lstlisting}[style=sh]
curl -s $ADMINAPI/status
\end{lstlisting}
\begin{lstlisting}[style=out]
{
  "namespace": "https://127.0.0.1:18443/v1/namespaces/flexgrid",
  "state": "Active",
  "versions": [
    { "version": "1.0.0", "manifest": "ni:///sha-256;cCYTVfOvidDw_pxQY5o1DHnHmfRjcSzDDAp3xLmxw2k",
      "local_role": "scheduled" }
  ],
  "last_accepted_listing": {
    "sequence": 1, "valid_from": "2026-09-01T10:00:00Z", "valid_until": "2026-09-02T10:00:00Z" },
  "hub_did": { "did": "did:key:z6Mkkfgq…", "resolved_at": "2026-09-01T10:00:00Z" },
  "rules": { "revision": 1, "count": 16 },
  "clock_skew_estimate_ms": 0
}
\end{lstlisting}

El estado es \cod{Active} y la 1.0.0 tiene el papel \cod{scheduled}: el cliente la conoce y tiene su paquete,
pero todavía no está vigente. Estamos a 1 de septiembre y empieza a valer el 3.

\noindent\begin{tabularx}{\linewidth}{@{}L{2.4cm}Y@{}}
\toprule
\textbf{Papel local} & \textbf{Significado} \\
\midrule
\cod{scheduled} & Conocida, aún no vigente. \\
\cod{head} & Vigente y la más reciente: contra ella se evalúa primero. \\
\cod{in_force} & Vigente pero no la más reciente. Solo ocurre durante una ventana de gracia. \\
\cod{retained} & Ya no vigente; su paquete se conserva porque un acuerdo la referencia. \\
\bottomrule
\end{tabularx}

\subsection{La máquina de estados}

El estado es una función pura del último listado aceptado y del instante local.

\begin{center}
\begin{tikzpicture}[font=\small,
  st/.style={draw=alxNavy, fill=alxLightBlue, rounded corners=2pt, minimum height=9mm, minimum width=22mm, align=center},
  bad/.style={st, fill=alxLightAmber, draw=alxAmber},
  dead/.style={st, fill=alxLightRed, draw=alxRed},
  lab/.style={font=\scriptsize, align=center},
  arr/.style={-{Stealth[length=2.2mm]}, alxNavy, thick}]
  \node[st, fill=alxSlate, draw=alxBorder] (on) at (0,0) {Onboarding};
  \node[st] (ac) at (4.4,0) {Active};
  \node[st] (tr) at (4.4,2.1) {Transition};
  \node[bad] (dg) at (8.8,0) {Degraded};
  \node[dead] (fc) at (13.2,0) {FailClosed};
  \draw[arr] (on) -- node[lab, above]{primer\\listado} (ac);
  \draw[arr] ([xshift=-3mm]ac.north) -- node[lab, left]{varias\\vigentes} ([xshift=-3mm]tr.south);
  \draw[arr] ([xshift=3mm]tr.south) -- node[lab, right]{una\\vigente} ([xshift=3mm]ac.north);
  \draw[arr] (ac) -- node[lab, above]{listado\\caducado} (dg);
  \draw[arr] (dg) -- node[lab, above]{pasa\\\texttt{max\_staleness}} (fc);
  \draw[arr, dashed] (fc.south) to[bend left=22] node[lab, below]{listado nuevo aceptado} (ac.south);
\end{tikzpicture}
\end{center}

\noindent\begin{tabularx}{\linewidth}{@{}L{2.4cm}L{5.2cm}Y@{}}
\toprule
\textbf{Estado} & \textbf{Condición} & \textbf{Efecto en los gates} \\
\midrule
\cod{Onboarding} & Ningún listado aceptado todavía. & \cod{CLIENT_NOT_READY}. \\
\cod{Active} & Listado válido, como mucho una versión vigente. & Clasificación normal. \\
\cod{Transition} & Listado válido, varias versiones vigentes. & Clasificación normal; puede haber avisos. \\
\cod{Degraded} & Listado caducado desde hace menos de \cod{max_staleness}. & Por gate, según \cod{degraded}: rechazar o seguir. \\
\cod{FailClosed} & Listado caducado desde hace \cod{max_staleness} o más. & \cod{CLIENT_FAIL_CLOSED} en todos. \\
\bottomrule
\end{tabularx}

\medskip
\texttt{GET /\allowbreak{}readyz} responde \cod{200} en \cod{Active}, \cod{Transition} y \cod{Degraded}, y \cod{503} en los otros
dos. Es lo que debe consultar un orquestador.

\section{La sincronización}

El cliente sincroniza al arrancar y, como mínimo, cada \cod{sync_interval}. Tras un fallo reintenta al minuto.
Cada pasada son seis fases, y cualquiera puede descartar el listado entero:

\noindent\begin{tabularx}{\linewidth}{@{}L{2.2cm}Y L{5.7cm}@{}}
\toprule
\textbf{Fase} & \textbf{Qué se comprueba} & \textbf{Códigos} \\
\midrule
Descarga & Se pide el listado. & \cod{SYNC_UNREACHABLE}, \cod{SYNC_NAMESPACE_NOT_FOUND} \\
Forma y firma & Esquema, contexto, emisor y \emph{namespace} esperados; la firma verifica contra el DID del hub. & \cod{SYNC_MALFORMED}, \cod{SYNC_SIGNATURE_INVALID}, \cod{SYNC_DID_UNRESOLVABLE} \\
Frescura & El instante local es anterior a \cod{validUntil}. & \cod{SYNC_EXPIRED} \\
Continuidad & La secuencia no retrocede; con igual secuencia, bytes idénticos; las revocaciones no desaparecen; las fechas emitidas no cambian. & \cod{SYNC_ROLLBACK}, \cod{SYNC_SEQUENCE_CONFLICT}, \cod{SYNC_REVOCATION_DROPPED}, \cod{SYNC_DATES_CHANGED} \\
Paquetes & Por cada versión nueva: el resumen del manifiesto coincide con el listado, su firma verifica y cada fichero coincide con su resumen. & \cod{SYNC_MANIFEST_DIGEST_MISMATCH}, \cod{SYNC_FILE_HASH_MISMATCH}, \cod{SYNC_FILE_SET_MISMATCH} \\
Fechas & Cada registro respeta la ventana de propagación y la gracia mínima. & \cod{SYNC_DATES_INVALID} \\
\bottomrule
\end{tabularx}

\begin{ideaClave}
\textbf{Todo o nada.} Si una sola versión del listado falla, se descarta el listado completo y sigue en uso la
instantánea anterior. Un paquete se usa entero o no se usa. Y un fallo de sincronización nunca cambia el estado
por sí mismo: lo que lo cambia es el paso del tiempo sobre el último listado bueno.
\end{ideaClave}

\begin{lstlisting}[style=sh]
curl -s $ADMINAPI/versions
\end{lstlisting}
\begin{lstlisting}[style=out]
[ { "version": "1.0.0", "class": "initial",
    "published_at": "2026-09-01T10:00:00Z",
    "new_content_from": "2026-09-03T10:00:00Z", "enforced_from": "2026-09-03T10:00:00Z",
    "revoked_from": null, "state": "Scheduled", "retained": true, … } ]
\end{lstlisting}

\section{Mover el tiempo}

Para que la 1.0.0 entre en vigor se avanza el reloj. El hub reemite el listado al ver que ha pasado más de medio
\cod{ttl}, y el cliente lo recoge. \cod{syncat} mueve el reloj y espera a que el cliente haya aceptado un
listado emitido en ese instante:

\begin{lstlisting}[style=sh]
syncat 2026-09-03T12:00:00Z
curl -s $ADMINAPI/status | python3 -c "import json,sys;d=json.load(sys.stdin);print(d['state'], d['versions'], d['last_accepted_listing'])"
\end{lstlisting}
\begin{lstlisting}[style=out]
Active [{'version': '1.0.0', 'manifest': 'ni:///sha-256;cCYT…', 'local_role': 'head'}]
       {'sequence': 2, 'valid_from': '2026-09-03T12:00:00Z', 'valid_until': '2026-09-04T12:00:00Z'}
\end{lstlisting}

La 1.0.0 es ahora \cod{head}, y la secuencia del listado ha subido a 2.

\begin{trampaHabitual}[Mover el reloj no basta: hay que esperar al listado]
Si se escribe una fecha en el fichero y se envía un mensaje en el mismo instante, el cliente puede decidir con
el listado anterior, que para la hora nueva está caducado, y responder en estado \cod{Degraded}. Tras cada salto
de reloj grande hay que esperar a la sincronización. Para eso existe \cod{syncat}.
\end{trampaHabitual}

\begin{retos}
\begin{enumerate}[leftmargin=*, itemsep=2pt]
  \item Arranca un segundo cliente con un \cod{hub_did} inventado y otro \cod{-state-dir}. ¿En qué estado se
        queda? Busca el código en su fichero de eventos.
  \item Consulta \cod{/readyz} y \cod{/healthz}. ¿En qué se diferencian?
  \item Para el cliente, bórrale el directorio de estado y arráncalo otra vez. ¿Qué vuelve a hacer?
\end{enumerate}
\end{retos}

\chapter{Reglas, gates y tráfico real}
\label{ch:reglas}

El cliente no sabe nada de DSP por sí mismo. Lo que sabe está en su \textbf{conjunto de reglas}: cada regla
dice qué peticiones cubre, qué parte del mensaje se valida, con qué gates y qué se hace con el resultado.

\section{Una regla}

\begin{lstlisting}[style=sh]
curl -s $ADMINAPI/rules/contract-request
\end{lstlisting}
\begin{lstlisting}[style=json]
{
  "id": "contract-request",
  "priority": 110,
  "listener": "dsp-ingress",
  "match": {
    "method": "POST",
    "path": "/negotiations/{providerPid}/request",
    "body": [ { "pointer": "/@type",
                "one_of": ["ContractRequestMessage", "dspace:ContractRequestMessage"] } ]
  },
  "message": "ContractRequestMessage",
  "verify": { "gates": ["dsp_message", "policy_intake"] },
  "mode": "enforce"
}
\end{lstlisting}

\noindent\begin{tabularx}{\linewidth}{@{}L{4cm}Y@{}}
\toprule
\textbf{Campo} & \textbf{Significado} \\
\midrule
\cod{priority} & Las de valor más bajo se prueban antes. Única por \emph{listener}. \\
\cod{listener} & En qué \emph{listener} aplica. \\
\cod{match.method}, \cod{match.path} & Método y plantilla de ruta. Una variable entre llaves casa con un segmento no vacío. \\
\cod{match.body} & Condiciones opcionales sobre el cuerpo, por puntero JSON. \\
\cod{message} & El tipo de mensaje DSP que la regla espera. \\
\cod{verify.gates} & De uno a cuatro gates. \\
\cod{verify.select} & Puntero JSON al nodo que se valida. Vacío, el cuerpo entero. \\
\cod{mode} & \cod{observe}, por defecto, \cod{shadow} o \cod{enforce}. \\
\cod{validate_response} & Validación pasiva de respuestas de catálogo. \\
\bottomrule
\end{tabularx}

\subsection{Cómo se elige la regla}

Para una petición, los candidatos son las reglas del \emph{listener} cuyo método y ruta casan, en orden de
prioridad. Se elige la primera cuyas condiciones de cuerpo se cumplen. Si no casa ninguna, se aplica
\cod{default_action}.

Se elige \textbf{una sola regla} por mensaje, y cada mensaje se decide con una única revisión del conjunto.

\subsection{Los tres modos}

\noindent\begin{tabularx}{\linewidth}{@{}L{2.6cm}Y Y Y@{}}
\toprule
\textbf{Clasificación} & \textbf{\texttt{observe}} & \textbf{\texttt{shadow}} & \textbf{\texttt{enforce}} \\
\midrule
\ok & reenvía & reenvía & reenvía \\
\warn & reenvía & reenvía y emite evento & reenvía con \cod{VHub-Verdict} y emite evento \\
\viol & reenvía & reenvía y emite evento & \textbf{rechaza} con \cod{400} y emite evento \\
\bottomrule
\end{tabularx}

\begin{ideaClave}
Los tres modos son un camino de despliegue. \cod{observe} no evalúa nada: la regla existe, pero está callada.
\cod{shadow} evalúa y cuenta lo que \emph{se habría} rechazado, sin tocar el tráfico. \cod{enforce} actúa. Una
regla se mueve entre ellos en caliente, y el cambio vale para el mensaje siguiente.
\end{ideaClave}

\subsection{La semilla}

La semilla trae las 16 reglas del enlace DSP 2025-1. Solo \cod{contract-request} viene en \cod{enforce}; las
otras quince, en \cod{observe}.

\begin{longtable}{@{}L{4.5cm}L{6.2cm}L{4.3cm}@{}}
\toprule
\textbf{Regla} & \textbf{Ruta} & \textbf{Gates} \\
\midrule
\endhead
\cod{catalog-request} & \texttt{POST /\allowbreak{}catalog/\allowbreak{}request} & \cod{dsp_message} \\
\cod{dataset-request} & \texttt{GET /\allowbreak{}catalog/\allowbreak{}datasets/\allowbreak{}\{id\}} & \cod{dsp_message} \\
\cod{contract-request-initial} & \texttt{POST /\allowbreak{}negotiations/\allowbreak{}request} & \cod{dsp_message}, \cod{policy_intake} \\
\cod{contract-request} & \texttt{POST /\allowbreak{}negotiations/\allowbreak{}\{providerPid\}/\allowbreak{}request} & \cod{dsp_message}, \cod{policy_intake} \\
\cod{contract-offer-initial} & \texttt{POST /\allowbreak{}negotiations/\allowbreak{}offers} & \cod{dsp_message}, \cod{policy_intake} \\
\cod{contract-offer} & \texttt{POST /\allowbreak{}negotiations/\allowbreak{}\{consumerPid\}/\allowbreak{}offers} & \cod{dsp_message}, \cod{policy_intake} \\
\cod{contract-agreement-in} & \texttt{POST /\allowbreak{}negotiations/\allowbreak{}\{consumerPid\}/\allowbreak{}agreement} & \cod{dsp_message}, \cod{policy_intake} \\
\cod{contract-agreement-out} & La misma, en \cod{dsp-egress} & \cod{dsp_message}, \cod{policy_intake} \\
\cod{agreement-verification} & \texttt{POST …/\allowbreak{}agreement/\allowbreak{}verification} & \cod{dsp_message} \\
\cod{negotiation-event} & \texttt{POST /\allowbreak{}negotiations/\allowbreak{}\{pid\}/\allowbreak{}events} & \cod{dsp_message} \\
\cod{negotiation-termination} & \texttt{POST /\allowbreak{}negotiations/\allowbreak{}\{pid\}/\allowbreak{}termination} & \cod{dsp_message} \\
\cod{transfer-request} & \texttt{POST /\allowbreak{}transfers/\allowbreak{}request} & \cod{dsp_message}, \cod{agreement_admission} \\
\cod{transfer-start} y otras tres & \texttt{POST /\allowbreak{}transfers/\allowbreak{}\{pid\}/\allowbreak{}…} & \cod{dsp_message} \\
\bottomrule
\end{longtable}

\section{El mensaje}

Este es \cod{examples/msgs/m1-forecast.json}: el agregador pide a la distribuidora la oferta del
capítulo~\ref{ch:compilar}, ahora dentro de un mensaje DSP.

\begin{lstlisting}[style=json]
{
  "@context": ["https://w3id.org/dspace/2025/1/context.jsonld",
               "https://flexgrid.example/odrl/profile/metering/context.jsonld"],
  "@type": "ContractRequestMessage",
  "providerPid": "urn:uuid:9c1f7a52-dso-neg-0042",
  "consumerPid": "urn:uuid:4be0d6c1-agg-neg-0017",
  "callbackAddress": "https://aggregator.flexgrid.example/dsp/callback",
  "offer": {
    "@type": "Offer",
    "@id": "urn:uuid:offer-forecast-01",
    "profile": "https://flexgrid.example/odrl/profile/metering",
    "target": "urn:uuid:dataset-qh-load-madrid-2026",
    "permission": [ { "action": "forecastDemand", "constraint": [ … las cinco restricciones … ] } ]
  }
}
\end{lstlisting}

Las dos URL de \cod{@context} son las que el paquete declara con \cod{url}: el contexto de DSP, vendorizado, y
el del perfil, generado. Ninguna se descarga: el cliente las resuelve dentro del paquete.

\section{Preguntar antes de enviar: \texttt{:explain}}

\texttt{POST /\allowbreak{}rules:explain} responde qué haría el proxy con una petición descrita, \textbf{sin efectos}: no
reenvía, no registra y no emite eventos. Recibe \cod{listener}, \cod{method}, \cod{path} y \cod{body}. La
función \cod{ex} lo construye a partir de un fichero de mensaje, y \cod{short} condensa la respuesta.

Con la regla de la semilla:

\begin{lstlisting}[style=sh]
ex m1-forecast | short
\end{lstlisting}
\begin{lstlisting}[style=out]
rule=contract-request state=Active action=reject code=VERDICT_VIOLATION
  gate=dsp_message    class=ok        codes=[] matched=1.0.0
  gate=policy_intake  class=violation codes=['VERDICT_VIOLATION'] matched=None
     - Violation json-schema  missing property 'profile'
     - Info loader  POC_EVAL_UNSUPPORTED_ROLE: shacl artefact profiles/odrl/metering/1.0.0/shapes.ttl is not evaluated
  verdict: violation;gate=policy_intake;code="VERDICT_VIOLATION";rule="contract-request"
\end{lstlisting}

Se rechaza un mensaje que es correcto. El hallazgo dice por qué: falta la propiedad \cod{profile}. Pero el
mensaje la tiene; está dentro de \cod{offer}.

\begin{trampaHabitual}[Una regla, un nodo]
\cod{schema.json} valida \textbf{una política}. La regla de la semilla no tiene \cod{verify.select}, así que el
gate \cod{policy_intake} recibe el mensaje entero, y un \cod{ContractRequestMessage} no es una política: la
política está en \cod{/offer}.

La solución es \texttt{"select": "/\allowbreak{}offer"}. Pero \cod{select} es \textbf{uno por regla} y vale para todos sus
gates, y solo se elige una regla por mensaje. Con \texttt{select: "/\allowbreak{}offer"}, el gate \cod{dsp_message} recibiría
también la oferta en lugar de la envoltura. Hoy no se pueden validar en un mismo mensaje la envoltura con un
esquema y la política con otro: hay que elegir.
\end{trampaHabitual}

Para una petición de contrato, lo que importa es la política. La regla adaptada está en
\cod{examples/rules/contract-request.json}:

\begin{lstlisting}[style=json]
  "verify": { "select": "/offer", "gates": ["policy_intake"] },
  "mode": "enforce"
\end{lstlisting}

\section{Cambiar una regla}

Todo cambio de reglas exige \cod{If-Match} con la revisión actual, que es el \cod{ETag} de \texttt{GET /\allowbreak{}rules}.

\begin{lstlisting}[style=sh]
curl -s -i -X PUT $ADMINAPI/rules/contract-request --data-binary @../rules/contract-request.json | head -1
curl -s -i -X PUT $ADMINAPI/rules/contract-request -H 'If-Match: "1"' \
  --data-binary @../rules/contract-request.json | grep -i -E '^HTTP|^etag'
\end{lstlisting}
\begin{lstlisting}[style=out]
HTTP/1.1 428 Precondition Required

HTTP/1.1 200 OK
Etag: "2"
\end{lstlisting}

Sin la cabecera, \cod{428}. Con una revisión desfasada sería \cod{412}. El conjunto entero se revalida en cada
cambio, y un cambio rechazado deja intacto el conjunto anterior.

\noindent\begin{tabularx}{\linewidth}{@{}C{1.8cm}Y L{5.7cm}@{}}
\toprule
\textbf{Invariante} & \textbf{Comprobación} & \textbf{Código} \\
\midrule
RULE-1 & El conjunto valida contra su esquema. & \cod{RULES_SCHEMA} \\
RULE-2 & Los identificadores son únicos. & \cod{RULES_DUPLICATE_ID} \\
RULE-3 & Las prioridades son únicas por \emph{listener}. & \cod{RULES_DUPLICATE_PRIORITY} \\
RULE-4 & \cod{agreement_admission} solo en peticiones de transferencia; \cod{validate_response} solo en catálogo. & \cod{RULES_GATE_MESSAGE_MISMATCH} \\
RULE-5 & Las reglas de \cod{mgmt-ingress} necesitan \cod{upstreams.management}. & \cod{RULES_UPSTREAM_MISSING} \\
\bottomrule
\end{tabularx}

\medskip
De paso se adaptan las otras dos reglas que usará el capítulo siguiente. \cod{setrule} lee la revisión actual y
hace el \cod{PUT}:

\begin{lstlisting}[style=sh]
setrule contract-agreement-in
setrule transfer-request
\end{lstlisting}
\begin{lstlisting}[style=out]
contract-agreement-in -> 200
transfer-request -> 200
\end{lstlisting}

\noindent\begin{tabularx}{\linewidth}{@{}L{4.7cm}Y@{}}
\toprule
\textbf{Regla} & \textbf{Cambio respecto de la semilla} \\
\midrule
\cod{contract-request} & \texttt{select: /\allowbreak{}offer}, solo \cod{policy_intake}. Ya estaba en \cod{enforce}. \\
\cod{contract-agreement-in} & \texttt{select: /\allowbreak{}agreement}, solo \cod{policy_intake}, pasa a \cod{enforce}. \\
\cod{transfer-request} & Pasa a \cod{enforce}. Conserva \cod{dsp_message} y \cod{agreement_admission} sobre el cuerpo entero. \\
\bottomrule
\end{tabularx}

\section{Seis mensajes}

Con la regla adaptada, estos son los veredictos a 3 de septiembre, con la 1.0.0 vigente:

\begin{lstlisting}[style=sh]
for m in m1-forecast m2-profiling m3-research m4-invalid m5-foreign-context m6-no-callback; do
  echo "--- $m"; ex $m | short
done
\end{lstlisting}
\begin{lstlisting}[style=out]
--- m1-forecast
rule=contract-request state=Active action=forward code=None
  gate=policy_intake  class=ok        codes=[] matched=1.0.0
--- m2-profiling
rule=contract-request state=Active action=forward code=None
  gate=policy_intake  class=ok        codes=[] matched=1.0.0
--- m3-research
rule=contract-request state=Active action=reject code=VERDICT_VIOLATION
  gate=policy_intake  class=violation codes=['VERDICT_VIOLATION'] matched=None
     - Violation json-schema /permission/0/constraint/0 validation failed
--- m4-invalid
rule=contract-request state=Active action=reject code=VERDICT_VIOLATION
  gate=policy_intake  class=violation codes=['VERDICT_VIOLATION'] matched=None
     - Violation json-schema /permission/0/action validation failed
     - Violation json-schema /permission/1/constraint/0 validation failed
--- m5-foreign-context
rule=contract-request state=Active action=reject code=CONTEXT_NOT_IN_PACKAGE
  gate=policy_intake  class=violation codes=['CONTEXT_NOT_IN_PACKAGE'] matched=None
     - Violation loader  @context https://schema.org/ is not in the package
--- m6-no-callback
rule=contract-request state=Active action=forward code=None
  gate=policy_intake  class=ok        codes=[] matched=1.0.0
\end{lstlisting}

En todas las respuestas aparece además la línea informativa \cod{POC_EVAL_UNSUPPORTED_ROLE}, omitida aquí.

\noindent\begin{tabularx}{\linewidth}{@{}L{3.9cm}Y C{1.9cm}@{}}
\toprule
\textbf{Mensaje} & \textbf{Qué lleva} & \textbf{Resultado} \\
\midrule
\cod{m1-forecast} & Previsión de demanda con cinco restricciones válidas. & \ok \\
\cod{m2-profiling} & Perfilado de consumidores. En la 1.0.0 es una acción activa. & \ok \\
\cod{m3-research} & Finalidad \cod{enr:research}, que la 1.0.0 no conoce. & \viol \\
\cod{m4-invalid} & Acción \cod{odrl:sell} y \texttt{minAggregation lt 5}. Dos fallos, dos hallazgos. & \viol \\
\cod{m5-foreign-context} & Un tercer contexto, \cod{https://schema.org/}. & \viol \\
\cod{m6-no-callback} & Sin \cod{callbackAddress}, obligatorio en DSP. & \ok \\
\bottomrule
\end{tabularx}

\medskip
Tres comentarios.

\textbf{m3} es una violación \emph{hoy}. El capítulo siguiente publica la 1.1.0 y el mismo mensaje pasa.

\textbf{m5} no se rechaza por su política, que es válida, sino por citar un contexto que el paquete no contiene.
El cliente nunca descarga un contexto para validar: si lo hiciera, quien controlara esa URL podría cambiar el
significado de un mensaje después de firmado el vocabulario. Esta comprobación la dispara el artefacto
\cod{shacl} del \emph{binding}: con solo \cod{json-schema}, los contextos no se miran.

\textbf{m6} pasa aunque le falta un campo obligatorio de DSP. Es el precio de la trampa anterior: esta regla ya
no evalúa \cod{dsp_message}.

\section{Tráfico de verdad}

\cod{row} envía un mensaje por \cod{dsp-ingress} e imprime el estado y la cabecera de veredicto.

\begin{lstlisting}[style=sh]
curl -s -i -X POST "$INGRESS/negotiations/$NEG/request" -H 'Content-Type: application/json' \
  --data-binary @../msgs/m1-forecast.json
\end{lstlisting}
\begin{lstlisting}[style=out]
HTTP/1.1 200 OK
Content-Length: 90
Content-Type: application/json

{"@type":"Ack","received":"POST /dsp/negotiations/urn:uuid:9c1f7a52-dso-neg-0042/request"}
\end{lstlisting}

Es la respuesta del conector, sin tocar. El conector de laboratorio recuerda lo que recibe:

\begin{lstlisting}[style=sh]
curl -s $CONNECTOR/_received
\end{lstlisting}

El cuerpo que aparece ahí tiene los mismos 1627 bytes que el fichero enviado.

Ahora el mensaje inválido:

\begin{lstlisting}[style=sh]
curl -s -i -X POST "$INGRESS/negotiations/$NEG/request" -H 'Content-Type: application/json' \
  --data-binary @../msgs/m4-invalid.json
\end{lstlisting}
\begin{lstlisting}[style=out]
HTTP/1.1 400 Bad Request
Content-Type: application/json
VHub-Verdict: violation;gate=policy_intake;code="VERDICT_VIOLATION";rule="contract-request"

{"@context":["https://w3id.org/dspace/2025/1/context.jsonld"],
 "@type":"ContractNegotiationError","code":"VERDICT_VIOLATION",
 "consumerPid":"urn:uuid:4be0d6c1-agg-neg-0017","providerPid":"urn:uuid:9c1f7a52-dso-neg-0042",
 "reason":["validation failed","validation failed"]}
\end{lstlisting}

El conector no ha recibido nada. La respuesta es el mensaje de error que DSP define para ese intercambio, un
\cod{ContractNegotiationError} con los identificadores del proceso copiados de la petición. Un conector al otro
lado lo entiende sin saber que existe Alexandria.

\cod{reason} lleva un elemento por hallazgo, dos en este caso. Los hallazgos informativos, como el aviso de que
el artefacto \cod{shacl} no se ejecuta, no se envían a la contraparte: quedan en \cod{:explain} y en el evento.

\subsection{La cabecera \texttt{VHub-Verdict}}

Es una lista de campos estructurados (RFC~9651), con un elemento por cada gate en \cod{enforce} cuya
clasificación no es \ok:

\begin{lstlisting}[style=out]
VHub-Verdict: violation;gate=policy_intake;code="VERDICT_VIOLATION";rule="contract-request"
VHub-Verdict: warning;gate=policy_intake;code="VERDICT_NOT_IN_HEAD";rule="contract-request";version="1.1.0";sunset=@1793613600
\end{lstlisting}

\section{Eventos y métricas}

Cada decisión que merece contarse se añade al fichero de eventos:

\begin{lstlisting}[style=sh]
tail -n 1 var/client/events.jsonl
\end{lstlisting}
\begin{lstlisting}[style=json]
{"type":"vhub.verdict","ts":"2026-09-03T12:00:00Z","listener":"dsp-ingress",
 "rule":"contract-request","rules_revision":2,"message":"ContractRequestMessage",
 "gate":"policy_intake","mode":"enforce","state":"Active",
 "classification":"violation","action":"reject","evaluated":["1.0.0"],"matched":null,
 "codes":["VERDICT_VIOLATION"],
 "findings":[
   {"slot":"odrl","source":"json-schema","severity":"Violation","path":"/permission/0/action","message":"validation failed"},
   {"slot":"odrl","source":"json-schema","severity":"Violation","path":"/permission/1/constraint/0","message":"validation failed"},
   … ]}
\end{lstlisting}

El evento dice qué regla, en qué revisión, contra qué versiones se evaluó, cuál aceptó el mensaje y, por
hallazgo, el slot, el artefacto y el puntero al nodo culpable.

\noindent\begin{tabularx}{\linewidth}{@{}L{3cm}Y@{}}
\toprule
\textbf{Tipo} & \textbf{Cuándo} \\
\midrule
\cod{vhub.verdict} & Un gate produce una clasificación que hay que contar. \\
\cod{vhub.sync} & Falla una sincronización. También lleva los avisos del evaluador. \\
\cod{vhub.state} & Cambia el estado, o el desfase de reloj supera el tolerado. \\
\cod{vhub.agreement} & Se programa la revocación de la versión de un acuerdo capturado. \\
\cod{vhub.rules} & Cambia el conjunto de reglas. \\
\cod{vhub.audit} & Se produce un informe de auditoría pasiva. \\
\cod{vhub.gc} & Pasa el recolector de paquetes. \\
\bottomrule
\end{tabularx}

\medskip
Y las métricas, en formato OpenMetrics:

\begin{lstlisting}[style=sh]
curl -s $ADMINAPI/metrics | grep -v '^#'
\end{lstlisting}
\begin{lstlisting}[style=out]
vhub_classifications_total{rule="contract-request",gate="policy_intake",version="",result="violation"} 1
vhub_classifications_total{rule="contract-request",gate="policy_intake",version="1.0.0",result="ok"} 1
vhub_listing_age_seconds 0
vhub_request_latency_seconds_sum 0.058964584
vhub_request_latency_seconds_count 2
\end{lstlisting}

\cod{vhub_listing_age_seconds} es la métrica sobre la que poner una alerta: si crece, el cliente ha dejado de
hablar con el hub, y cuando pase de \cod{ttl} entrará en \cod{Degraded}.

\begin{retos}
\begin{enumerate}[leftmargin=*, itemsep=2pt]
  \item Pasa \cod{contract-request} a \cod{shadow} y envía \cod{m4-invalid}. ¿Qué recibe el emisor? ¿Qué queda
        en el fichero de eventos?
  \item Añade con \texttt{POST /\allowbreak{}rules} una copia de una regla existente. ¿Qué invariante salta?
  \item Cambia \cod{default_action} a \cod{reject} en la configuración, reinicia el cliente y envía una petición
        a una ruta sin regla.
\end{enumerate}
\end{retos}

\chapter{Una migración de principio a fin}
\label{ch:migracion}

Este capítulo recorre dos meses y medio de vida de Flexgrid. Se publican la 1.1.0 y la 2.0.0, se firma un
acuerdo, se retira una versión y se revoca. En cada paso se envían los mismos tres mensajes y se mira qué
cambia.

\noindent\begin{tabularx}{\linewidth}{@{}L{2.8cm}Y@{}}
\toprule
\textbf{Mensaje} & \textbf{Qué lleva y bajo qué versiones es válido} \\
\midrule
\cod{m1-forecast} & Previsión de demanda. Válido en las tres versiones. \\
\cod{m2-profiling} & Perfilado de consumidores. Válido en 1.0.0 y 1.1.0; inválido en 2.0.0. \\
\cod{m3-research} & Finalidad de investigación. Inválido en 1.0.0; válido desde 1.1.0. \\
\bottomrule
\end{tabularx}

\medskip
\cod{rows} envía los tres por \cod{dsp-ingress} e imprime, por cada uno, la línea de estado HTTP y la cabecera
\cod{VHub-Verdict}, seguidas del estado del cliente y sus versiones.

\section{El punto de partida}

3 de septiembre, con la 1.0.0 recién entrada en vigor:

\begin{lstlisting}[style=sh]
rows
\end{lstlisting}
\begin{lstlisting}[style=out]
m1-forecast    HTTP/1.1 200 OK | -
m2-profiling   HTTP/1.1 200 OK | -
m3-research    HTTP/1.1 400 Bad Request | VHub-Verdict: violation;gate=policy_intake;code="VERDICT_VIOLATION";rule="contract-request"
   Active [('1.0.0', 'head')]
\end{lstlisting}

\section{Una versión compatible}

El 20 de septiembre Lucía publica la 1.1.0, que añade la finalidad de investigación y marca obsoleto el
perfilado.

\begin{lstlisting}[style=sh]
syncat 2026-09-20T10:00:00Z
publish ../dsl/flexgrid-1.1.0.yaml
\end{lstlisting}
\begin{lstlisting}[style=out]
draft:     ffd85632051f7891
{"@context":"…","class":"compatible","findings":[],"valid":true}
candidate: 85089995ebc525b7
\end{lstlisting}

El hub la clasifica como \cod{compatible} por su número, y N4 ha comprobado contra la fuente de la 1.0.0 que lo
es de verdad. El \cod{diff} del candidato muestra el cambio en términos de ficheros: aparecen los seis de
\cod{profiles/odrl/metering/1.1.0/}, desaparecen los seis de \cod{1.0.0/} y cambian los cuatro del espacio y la
fuente.

Una versión compatible no admite ventana de gracia, porque no hay nada que hacer convivir:

\begin{lstlisting}[style=sh]
approve '{"grace_window":"P10D"}'
\end{lstlisting}
\begin{lstlisting}[style=out]
HTTP/2 422
{"…","title":"DATES_GRACE_ON_COMPATIBLE","status":422,
 "detail":"grace_window is only allowed on breaking versions","code":"DATES_GRACE_ON_COMPATIBLE","severity":"error"}
\end{lstlisting}

\begin{lstlisting}[style=sh]
approve '{}'
\end{lstlisting}
\begin{lstlisting}[style=out]
HTTP/2 201
{'version': '1.1.0', 'class': 'compatible', 'published_at': '2026-09-20T10:00:00Z',
 'new_content_from': '2026-09-22T10:00:00Z', 'enforced_from': '2026-09-22T10:00:00Z'}
\end{lstlisting}

Tres días después:

\begin{lstlisting}[style=sh]
syncat 2026-09-23T12:00:00Z
rows
\end{lstlisting}
\begin{lstlisting}[style=out]
m1-forecast    HTTP/1.1 200 OK | -
m2-profiling   HTTP/1.1 200 OK | -
m3-research    HTTP/1.1 200 OK | -
   Active [('1.1.0', 'head')]
\end{lstlisting}

La universidad ya puede negociar. El estado sigue siendo \cod{Active}, con una sola versión: la 1.1.0 ha
sustituido limpiamente a la 1.0.0, que pasa a \cod{Superseded}.

Y \cod{m2-profiling} pasa sin ninguna señal, aunque usa un término obsoleto. Es la trampa del
capítulo~\ref{ch:evolucion} vista en producción: el aviso existe en \cod{shapes.ttl}, y nadie lo ejecuta.

\section{Una versión de ruptura}

El 1 de octubre se publica la 2.0.0, que retira el perfilado.

\begin{lstlisting}[style=sh]
syncat 2026-10-01T10:00:00Z
publish ../dsl/flexgrid-2.0.0.yaml
\end{lstlisting}
\begin{lstlisting}[style=out]
draft:     a5fad501c89c41e8
{"@context":"…","class":"breaking","findings":[],"valid":true}
candidate: 5552a13ebfe21682
\end{lstlisting}

Lucía intenta acortar la convivencia a tres días:

\begin{lstlisting}[style=sh]
approve '{"grace_window":"P3D"}'
\end{lstlisting}
\begin{lstlisting}[style=out]
HTTP/2 422
{"…","title":"DATES_GRACE_TOO_SHORT","status":422,
 "detail":"grace_window P3D is below min_grace_window P7D","code":"DATES_GRACE_TOO_SHORT","severity":"error"}
\end{lstlisting}

El mínimo son siete días y es del hub: ninguna \emph{steward} puede dejar a los participantes sin tiempo para
migrar. Con la gracia por defecto, de treinta:

\begin{lstlisting}[style=sh]
approve '{}'
\end{lstlisting}
\begin{lstlisting}[style=out]
HTTP/2 201
{'version': '2.0.0', 'class': 'breaking', 'published_at': '2026-10-01T10:00:00Z',
 'new_content_from': '2026-10-03T10:00:00Z', 'enforced_from': '2026-11-02T10:00:00Z'}
\end{lstlisting}

Ahora las dos fechas son distintas, y entre ellas está la ventana de gracia.

\begin{center}
\begin{tikzpicture}[font=\scriptsize, x=0.156cm, y=1cm]
  \draw[alxBorder] (0,-0.5) -- (80,-0.5);
  \foreach \d/\l in {2/3 sep, 21/22 sep, 32/3 oct, 62/2 nov, 72/12 nov}
    { \draw[alxBorder] (\d,-0.5) -- (\d,-0.65); \node[below] at (\d,-0.68) {\l}; }
  \fill[alxAmber, opacity=0.13] (32,0.25) rectangle (62,1.75);
  \node[alxAmber!80!black] at (47,1.0) {ventana de gracia: conviven};
  \draw[alxBlue, dashed, line width=1.2pt] (0,2.2) -- (2,2.2);
  \draw[alxBlue, line width=3.2pt] (2,2.2) -- (21,2.2);
  \node[left] at (0,2.2) {1.0.0};
  \draw[alxTeal, dashed, line width=1.2pt] (19,1.4) -- (21,1.4);
  \draw[alxTeal, line width=3.2pt] (21,1.4) -- (62,1.4);
  \draw[alxTeal!45, line width=3.2pt] (62,1.4) -- (72,1.4);
  \node[left] at (19,1.4) {1.1.0};
  \node[above, alxTeal!70!black] at (67,1.45) {retenida};
  \node[alxRed, right] at (72,1.4) {revocada};
  \draw[alxPurple, dashed, line width=1.2pt] (30,0.6) -- (32,0.6);
  \draw[alxPurple, line width=3.2pt] (32,0.6) -- (80,0.6);
  \node[left] at (30,0.6) {2.0.0};
\end{tikzpicture}

{\scriptsize El tramo discontinuo es la ventana de propagación: publicada, pero aún no vigente.}
\end{center}

\section{Lo que ve el borde, instante a instante}

\begin{lstlisting}[style=sh]
syncat 2026-10-02T12:00:00Z; rows
syncat 2026-10-03T06:00:00Z; clock 2026-10-03T09:59:55Z; rows
syncat 2026-10-10T12:00:00Z; rows
syncat 2026-11-02T06:00:00Z; clock 2026-11-02T10:00:05Z; rows
syncat 2026-11-05T12:00:00Z; rows
\end{lstlisting}

\begin{longtable}{@{}L{2.7cm}C{1.7cm}C{1.7cm}C{1.7cm}L{3.3cm}L{3.2cm}@{}}
\toprule
\textbf{Instante} & \textbf{m1} & \textbf{m2} & \textbf{m3} & \textbf{Estado y versiones} & \textbf{Por qué} \\
\midrule
\endhead
3 sep, 12:00 & \ok & \ok & \viol & \cod{Active}: 1.0.0 & Solo la 1.0.0. \\
23 sep, 12:00 & \ok & \ok & \ok & \cod{Active}: 1.1.0 & La 1.1.0 la ha sustituido. \\
2 oct, 12:00 & \ok & \ok & \ok & \cod{Active}: 1.1.0; 2.0.0 programada & La 2.0.0 se conoce, pero no vale aún. \\
3 oct, 09:59:55 & \ok & \warn & \ok & \cod{Transition}: 2.0.0 cabeza, 1.1.0 vigente & La 2.0.0 entra un margen de arrendamiento antes. \\
10 oct, 12:00 & \ok & \warn & \ok & \cod{Transition} & Ventana de gracia. \\
2 nov, 10:00:05 & \ok & \warn & \ok & \cod{Transition} & Aún dentro del margen tras \cod{enforced_from}. \\
5 nov, 12:00 & \ok & \viol & \ok & \cod{Active}: 2.0.0; 1.1.0 retenida & Solo la 2.0.0. \\
\bottomrule
\end{longtable}

Es el patrón \emph{expandir, migrar, contraer}. Primero se acepta lo nuevo. Lo viejo sigue funcionando con un
aviso. El rechazo empieza solo después de \cod{enforced_from}.

\subsection{El aviso}

Durante la ventana, \cod{m2-profiling} se reenvía al conector y la respuesta vuelve con una cabecera:

\begin{lstlisting}[style=out]
m2-profiling   HTTP/1.1 200 OK | VHub-Verdict: warning;gate=policy_intake;code="VERDICT_NOT_IN_HEAD";rule="contract-request";version="1.1.0";sunset=@1793613600
\end{lstlisting}

\noindent\begin{tabularx}{\linewidth}{@{}L{5.7cm}Y@{}}
\toprule
\textbf{Parámetro} & \textbf{Significado} \\
\midrule
\cod{warning} & La clasificación. \\
\cod{gate=policy_intake} & El gate que la produjo. \\
\texttt{code="VERDICT\_NOT\_IN\_HEAD"} & El mensaje no conforma con la versión cabeza, pero sí con una anterior aún vigente. \\
\texttt{version="1.1.0"} & La versión que lo aceptó. \\
\cod{sunset=@1793613600} & Instante en que dejará de aceptarse: 2 de noviembre de 2026 a las 10:00 UTC. \\
\bottomrule
\end{tabularx}

\medskip
La comercializadora no tiene que leer ningún boletín. Su propio tráfico le dice que lo que envía tiene fecha de
caducidad, y cuál.

\cod{:explain} muestra el razonamiento completo:

\begin{lstlisting}[style=sh]
ex m2-profiling | short
\end{lstlisting}
\begin{lstlisting}[style=out]
rule=contract-request state=Transition action=forward-notice code=None
  gate=policy_intake  class=warning   codes=['VERDICT_NOT_IN_HEAD'] matched=1.1.0
     - Violation json-schema /permission/0/action validation failed
     - Info loader  POC_EVAL_UNSUPPORTED_ROLE: shacl artefact profiles/odrl/metering/2.0.0/shapes.ttl is not evaluated
  verdict: warning;gate=policy_intake;code="VERDICT_NOT_IN_HEAD";rule="contract-request";version="1.1.0";sunset=@1793613600
\end{lstlisting}

Los hallazgos son los de la cabeza, la 2.0.0, que rechaza la acción. \cod{matched=1.1.0} dice quién lo aceptó.

\begin{bajoElCapo}[El algoritmo de clasificación]
Para un gate, después de las filas de estado:
\begin{enumerate}[leftmargin=*, itemsep=1pt]
  \item Cuerpo ilegible, o falta el nodo de \cod{select}: violación, \cod{MESSAGE_UNPARSEABLE}.
  \item Ninguna versión vigente: violación, \cod{NO_VERSION_IN_FORCE}.
  \item Se evalúa contra la \textbf{cabeza}. Si conforma: ok.
  \item Si no, contra las anteriores aún vigentes, de la más nueva a la más vieja. Si una lo acepta: aviso,
        \cod{VERDICT_NOT_IN_HEAD}.
  \item Si ninguna: violación. \cod{CONTEXT_NOT_IN_PACKAGE} si el motivo es un contexto ajeno;
        \cod{VERDICT_VIOLATION} en otro caso.
\end{enumerate}
\end{bajoElCapo}

\subsection{El margen de arrendamiento}

Dos filas de la tabla tienen horas extrañas. El 3 de octubre a las 09:59:55 la 2.0.0 ya acepta mensajes, cinco
segundos antes de su \cod{new_content_from}. El 2 de noviembre a las 10:00:05, cinco segundos después de
\cod{enforced_from}, la 1.1.0 todavía se acepta.

Es el \cod{lease_margin}, de diez segundos: el cliente ensancha cada frontera por los dos lados. Los relojes de
dos participantes nunca coinciden del todo, y sin margen un mensaje enviado en el segundo exacto del cambio
podría ser válido para quien lo envía e inválido para quien lo recibe.

\section{Los acuerdos sobreviven a su versión}

El 10 de octubre, en plena ventana, la comercializadora y la distribuidora cierran un acuerdo que incluye el
perfilado. Es \cod{examples/msgs/a1-agreement.json}:

\begin{lstlisting}[style=json]
{
  "@context": [ … ],
  "@type": "ContractAgreementMessage",
  "providerPid": "urn:uuid:9c1f7a52-dso-neg-0042",
  "consumerPid": "urn:uuid:7d2a9f10-ret-neg-0007",
  "agreement": {
    "@id": "urn:uuid:agreement-retail-0007", "@type": "Agreement",
    "profile": "https://flexgrid.example/odrl/profile/metering",
    "target": "urn:uuid:dataset-qh-load-madrid-2026",
    "assigner": "did:web:dso.flexgrid.example",
    "assignee": "did:web:retailer.flexgrid.example",
    "permission": [ { "action": "profileConsumer", "constraint": [
        {"leftOperand": "purpose", "operator": "eq", "rightOperand": {"@id": "enr:billing"}},
        {"leftOperand": "retentionDays", "operator": "lteq", "rightOperand": 30} ] } ]
  }
}
\end{lstlisting}

\begin{lstlisting}[style=sh]
post /negotiations/urn:uuid:7d2a9f10-ret-neg-0007/agreement a1-agreement
curl -s "$ADMINAPI/agreements/urn:uuid:agreement-retail-0007"
\end{lstlisting}
\begin{lstlisting}[style=out]
a1-agreement   HTTP/1.1 200 OK | VHub-Verdict: warning;gate=policy_intake;code="VERDICT_NOT_IN_HEAD";rule="contract-agreement-in";version="1.1.0";sunset=@1793613600

{ "agreement_id": "urn:uuid:agreement-retail-0007",
  "provider_pid": "urn:uuid:9c1f7a52-dso-neg-0042",
  "consumer_pid": "urn:uuid:7d2a9f10-ret-neg-0007",
  "role": "consumer",
  "version": "1.1.0",
  "manifest": "ni:///sha-256;8kvT3m-30MLD6JA7VFtn-iNDqtFe2fYvyx_uG5ImLgI",
  "declared_profile": "https://flexgrid.example/odrl/profile/metering",
  "recorded_at": "2026-10-10T12:00:00Z",
  "source": "dsp-ingress",
  "usable": true }
\end{lstlisting}

El acuerdo pasa con aviso, y el cliente lo \textbf{captura}: guarda su identificador, los de las dos partes, su
papel y, sobre todo, la versión que lo validó, la 1.1.0. La captura ocurre cuando el conector responde \cod{2xx}
y antes de devolver la respuesta.

\begin{ideaClave}
Un acuerdo queda ligado a la versión del vocabulario bajo la que se firmó, no a la vigente. Un contrato no se
reinterpreta cada vez que cambian las reglas.
\end{ideaClave}

Con el acuerdo capturado, una transferencia bajo él se admite. \cod{examples/msgs/t1-transfer.json} es un
\cod{TransferRequestMessage} con \cod{agreementId}:

\begin{lstlisting}[style=sh]
post /transfers/request t1-transfer
post /transfers/request t2-transfer-no-format -b
\end{lstlisting}
\begin{lstlisting}[style=out]
t1-transfer             HTTP/1.1 200 OK | -
t2-transfer-no-format   HTTP/1.1 400 Bad Request | VHub-Verdict: violation;gate=dsp_message;code="VERDICT_VIOLATION";rule="transfer-request"
    {"@context":["https://w3id.org/dspace/2025/1/context.jsonld"],"@type":"TransferError",
     "code":"VERDICT_VIOLATION","consumerPid":"urn:uuid:7d2a9f10-ret-trf-0001","providerPid":null,
     "reason":["'allOf' failed","missing property 'format'"]}
\end{lstlisting}

La segunda transferencia es la primera vez que el slot \cod{dsp} vendorizado rechaza algo: al mensaje le falta
\cod{format}, que el esquema de \cod{vendor/dsp/2025-1/} exige. El error es ahora un \cod{TransferError}, el que
DSP define para ese intercambio.

\subsection{Después de la retirada}

El 5 de noviembre la 1.1.0 ya no está vigente. Una petición nueva con perfilado se rechaza. Pero:

\begin{lstlisting}[style=sh]
curl -s $ADMINAPI/versions | python3 -c "import json,sys;[print(v['version'], v['class'], v['state'], 'retained' if v['retained'] else '') for v in json.load(sys.stdin)]"
post /transfers/request t1-transfer
\end{lstlisting}
\begin{lstlisting}[style=out]
1.0.0 initial Superseded
1.1.0 compatible Retired retained
2.0.0 breaking Published retained

t1-transfer             HTTP/1.1 200 OK | -
\end{lstlisting}

La 1.1.0 está \cod{Retired}, pero su paquete se \textbf{retiene} porque el acuerdo la referencia, y la
transferencia se sigue admitiendo. La 1.0.0, a la que no apunta ningún acuerdo, ya no se conserva.

\noindent\begin{tabularx}{\linewidth}{@{}L{2.6cm}Y@{}}
\toprule
\textbf{Estado} & \textbf{Significado, en un instante de referencia} \\
\midrule
\cod{Scheduled} & Antes de \cod{new_content_from}. \\
\cod{Published} & Vigente y cabeza. \\
\cod{Deprecated} & Vigente, pero no cabeza. \\
\cod{Superseded}, \cod{Retired} & Ya no vigente. \\
\cod{Revoked} & Revocada. \\
\bottomrule
\end{tabularx}

\section{Revocar}

El 10 de noviembre el regulador declara ilícito el perfilado individual. Retirar ya no basta: los acuerdos
firmados bajo la 1.1.0 no deben amparar más transferencias.

\begin{lstlisting}[style=sh]
syncat 2026-11-10T00:00:00Z
R=$HUB/v1/namespaces/flexgrid/versions
curl -sk -X POST "$R/1.1.0:revoke" -H "Authorization: Bearer $(cat steward.token)" -d '{}'
curl -sk -X POST "$R/2.0.0:revoke" -H "Authorization: Bearer $(cat steward.token)" -d '{"justification":"test"}'
\end{lstlisting}
\begin{lstlisting}[style=out]
{"…","title":"DATES_JUSTIFICATION_REQUIRED","status":422,"detail":"a revocation requires a justification", … }
{"…","title":"DATES_REVOKE_LEAVES_NO_HEAD","status":422,"detail":"no version would be in force at 2026-11-12T00:00:00Z", … }
\end{lstlisting}

Los dos intentos fallan, y con razón: una revocación exige justificación, y no se puede revocar la única
versión vigente, porque el espacio se quedaría sin reglas.

\noindent\begin{tabularx}{\linewidth}{@{}C{1.3cm}Y L{6.1cm}@{}}
\toprule
\textbf{Regla} & \textbf{Comprobación} & \textbf{Código} \\
\midrule
REV-1 & La versión no está ya revocada. & \cod{DATES_ALREADY_REVOKED} \\
REV-2 & Hay justificación. & \cod{DATES_JUSTIFICATION_REQUIRED} \\
REV-3 & \cod{revoked_from} es al menos el instante de la petición más la ventana de propagación. & \cod{DATES_REVOCATION_TOO_EARLY} \\
REV-4 & Queda alguna versión vigente. & \cod{DATES_REVOKE_LEAVES_NO_HEAD} \\
\bottomrule
\end{tabularx}

\begin{lstlisting}[style=sh]
curl -sk -X POST "$R/1.1.0:revoke" -H "Authorization: Bearer $(cat steward.token)" \
  -d '{"justification":"Individual consumer profiling was ruled unlawful by the regulator"}'
\end{lstlisting}
\begin{lstlisting}[style=out]
{ … "credentialSubject": { "version": "1.1.0", "class": "compatible",
    "enforced_from": "2026-09-22T10:00:00Z", "revoked_from": "2026-11-12T00:00:00Z", … } }
\end{lstlisting}

La revocación tampoco es instantánea: \cod{revoked_from} es el 12 de noviembre, dos días después. Las fechas ya
emitidas no cambian; \cod{revoked_from} se añade una sola vez.

El listado siguiente lleva la revocación, y el cliente, al verla, avisa de cada acuerdo afectado:

\begin{lstlisting}[style=sh]
grep AGREEMENT var/client/events.jsonl
\end{lstlisting}
\begin{lstlisting}[style=json]
{"type":"vhub.agreement","ts":"…","code":"AGREEMENT_REVOCATION_SCHEDULED","severity":"info",
 "detail":"revoked_from 2026-11-12T00:00:00Z","version":"1.1.0",
 "agreement":"urn:uuid:agreement-retail-0007"}
\end{lstlisting}

Ese evento es lo que el operador de la comercializadora debe conectar a su sistema de alertas: tiene dos días
para renegociar.

\begin{lstlisting}[style=sh]
clock 2026-11-12T00:00:05Z; sleep 2; post /transfers/request t1-transfer
clock 2026-11-12T00:00:10Z; sleep 2; post /transfers/request t1-transfer -b
\end{lstlisting}
\begin{lstlisting}[style=out]
t1-transfer   HTTP/1.1 200 OK | -

t1-transfer   HTTP/1.1 400 Bad Request | VHub-Verdict: violation;gate=agreement_admission;code="AGREEMENT_VERSION_REVOKED";rule="transfer-request"
    {"@context":["https://w3id.org/dspace/2025/1/context.jsonld"],"@type":"TransferError",
     "code":"AGREEMENT_VERSION_REVOKED","consumerPid":"urn:uuid:7d2a9f10-ret-trf-0001","providerPid":null,
     "reason":["AGREEMENT_VERSION_REVOKED"]}
\end{lstlisting}

A las 00:00:05 la transferencia aún pasa, dentro del margen de arrendamiento. A las 00:00:10 se rechaza, y el
acuerdo figura como \texttt{"usable": false}.

\begin{ideaClave}
\textbf{Retirar} no rompe los acuerdos existentes. \textbf{Revocar} sí, y se anuncia con antelación. Son dos
herramientas distintas: la primera es para evolucionar el vocabulario, la segunda para cuando algo que se
permitió no debió permitirse.
\end{ideaClave}

\section{El rastro en el hub}

Los eventos del \emph{namespace} cuentan la misma historia desde el otro lado. Resumidos:

\begin{lstlisting}[style=out]
10 2026-09-20T10:00:00Z draft.created        lucia
11 2026-09-20T10:00:00Z candidate.proposed   lucia
12 2026-09-20T10:00:00Z release.aborted      lucia     <- el intento con grace_window
13 2026-09-20T10:00:00Z candidate.approved   lucia
14 2026-09-20T10:00:00Z release.published    lucia     versions/1.1.0
15 2026-09-20T10:00:00Z listing.issued       vhub
…
20 2026-10-01T10:00:00Z release.aborted      lucia     <- el intento con P3D
21 2026-10-01T10:00:00Z candidate.approved   lucia
22 2026-10-01T10:00:00Z release.published    lucia     versions/2.0.0
…
30 2026-11-10T00:00:00Z version.revoked      lucia     versions/1.1.0
31 2026-11-10T00:00:00Z listing.issued       vhub
\end{lstlisting}

Los dos intentos de aprobación fallidos también quedaron registrados, como \cod{release.aborted}. Una aprobación
que falla no deja nada visible para los clientes, pero sí deja constancia de que se intentó.

\begin{retos}
\begin{enumerate}[leftmargin=*, itemsep=2pt]
  \item Repite la migración aprobando la 2.0.0 con \texttt{\{"grace\_window":"P10D"\}} y calcula antes a mano el
        nuevo \cod{enforced_from} y el \cod{sunset} de la cabecera.
  \item Pon \texttt{unknown\_agreement: reject} en el cliente y envía una transferencia con un \cod{agreementId} que
        nunca se capturó.
  \item ¿Qué pasaría con el acuerdo si en lugar de revocar la 1.1.0 se publicara una 3.0.0? Compruébalo.
\end{enumerate}
\end{retos}

\chapter{Cuando algo falla, y dónde se despliega}
\label{ch:resiliencia}

\section{El hub desaparece}

El cliente decide con estado local, así que un hub caído no detiene el tráfico. Lo que cambia con el tiempo es
cuánto se fía el cliente de lo último que supo.

Se para el hub, y solo el hub:

\begin{lstlisting}[style=sh]
kill $(lsof -nP -iTCP:18443 -sTCP:LISTEN -t)
\end{lstlisting}

\subsection{Primera fase: no pasa nada}

El último listado aceptado vale hasta el 13 de noviembre a las 00:00:05. Doce horas antes:

\begin{lstlisting}[style=sh]
clock 2026-11-12T12:00:00Z; sleep 2
curl -s $ADMINAPI/status | python3 -c "import json,sys;d=json.load(sys.stdin);print(d['state'], d['last_accepted_listing'])"
row m1-forecast
\end{lstlisting}
\begin{lstlisting}[style=out]
Active {'sequence': 15, 'valid_from': '2026-11-12T00:00:05Z', 'valid_until': '2026-11-13T00:00:05Z'}
m1-forecast   HTTP/1.1 200 OK | -
\end{lstlisting}

El estado sigue siendo \cod{Active}. El fallo de sincronización queda registrado, pero no cambia nada:

\begin{lstlisting}[style=json]
{"type":"vhub.sync","ts":"2026-11-12T12:00:00Z","code":"SYNC_UNREACHABLE","severity":"error",
 "detail":"Get \"https://127.0.0.1:18443/v1/namespaces/flexgrid/versions\": dial tcp 127.0.0.1:18443: connect: connection refused"}
\end{lstlisting}

\subsection{Segunda fase: \texttt{Degraded}}

Pasada la caducidad del listado:

\begin{lstlisting}[style=sh]
clock 2026-11-13T06:00:00Z; sleep 2
row m1-forecast -b
\end{lstlisting}
\begin{lstlisting}[style=out]
m1-forecast   HTTP/1.1 400 Bad Request | VHub-Verdict: violation;gate=policy_intake;code="CLIENT_DEGRADED_REJECT";rule="contract-request"
    {"@context":["https://w3id.org/dspace/2025/1/context.jsonld"],"@type":"ContractNegotiationError",
     "code":"CLIENT_DEGRADED_REJECT", … ,"reason":["CLIENT_DEGRADED_REJECT"]}
\end{lstlisting}

El estado es \cod{Degraded} y \cod{/readyz} sigue respondiendo \cod{200}. Un mensaje que ayer era válido hoy se
rechaza, y el código explica que no es culpa del mensaje.

Es exactamente lo que el \cod{SpaceProfile} pidió en el capítulo~\ref{ch:spaceprofile}:

\begin{lstlisting}[style=yaml]
  degraded:
    policy_intake: reject
    dsp_message: continue
\end{lstlisting}

Sin contacto con el hub, Flexgrid deja de admitir \textbf{políticas nuevas}, porque podría haberse revocado
algo sin que el cliente lo sepa. Pero un gate con \cod{continue} sigue clasificando con normalidad. Una
transferencia mal formada lo demuestra, porque se sigue evaluando contra el esquema de DSP:

\begin{lstlisting}[style=out]
t2-transfer-no-format   HTTP/1.1 400 Bad Request | VHub-Verdict: violation;gate=dsp_message;code="VERDICT_VIOLATION";rule="transfer-request", violation;gate=agreement_admission;code="AGREEMENT_VERSION_REVOKED";rule="transfer-request"
\end{lstlisting}

La cabecera lleva aquí dos elementos, uno por gate. Cuando una regla tiene varios gates se toma la acción más
restrictiva, y cada gate que no es \ok\ aporta su veredicto.

\begin{ideaClave}
\cod{degraded} es una decisión de negocio que toma la autoridad, no cada participante: viaja en el manifiesto
firmado. Responde a la pregunta «si no sé si las reglas han cambiado, ¿qué es peor: parar o seguir?», y la
respuesta puede ser distinta para cada gate.
\end{ideaClave}

\subsection{Tercera fase: \texttt{FailClosed}}

Pasado \cod{max_staleness}, siete días, desde la caducidad:

\begin{lstlisting}[style=sh]
clock 2026-11-21T06:00:00Z; sleep 2
row m1-forecast
curl -s -o /dev/null -w '%{http_code}\n' $ADMINAPI/readyz
\end{lstlisting}
\begin{lstlisting}[style=out]
m1-forecast   HTTP/1.1 400 Bad Request | VHub-Verdict: violation;gate=policy_intake;code="CLIENT_FAIL_CLOSED";rule="contract-request"
503
\end{lstlisting}

Ahora todos los gates rechazan, digan lo que digan sus ajustes, y \cod{/readyz} pasa a \cod{503}. Los cambios de
estado quedan como eventos:

\begin{lstlisting}[style=json]
{"type":"vhub.state","ts":"2026-11-13T06:00:00Z","severity":"info","detail":"state changed","from":"Active","to":"Degraded"}
{"type":"vhub.state","ts":"2026-11-21T06:00:00Z","severity":"info","detail":"state changed","from":"Degraded","to":"FailClosed"}
\end{lstlisting}

\subsection{La recuperación}

\begin{lstlisting}[style=sh]
$BIN/alexandria-server -config server.yaml -clock-file clock >> server.log 2>&1 &
clock 2026-11-21T06:10:00Z; sleep 3
curl -s $ADMINAPI/status | python3 -c "import json,sys;d=json.load(sys.stdin);print(d['state'], d['last_accepted_listing'])"
row m1-forecast
\end{lstlisting}
\begin{lstlisting}[style=out]
Active {'sequence': 16, 'valid_from': '2026-11-21T06:00:00Z', 'valid_until': '2026-11-22T06:00:00Z'}
m1-forecast   HTTP/1.1 200 OK | -
\end{lstlisting}

Al volver, el hub emite un listado nuevo, de secuencia 16. El cliente reintenta un minuto después del último
fallo, lo acepta y vuelve a \cod{Active} sin intervención. No hay que reiniciar nada.

\begin{trampaHabitual}[Con reloj de fichero, el reintento también espera al reloj]
El reintento se cuenta en tiempo del reloj, no de pared. Si arrancas el hub y no mueves el fichero, el cliente
no reintenta nunca: para él no ha pasado ni un segundo. Por eso el comando anterior avanza diez minutos.
\end{trampaHabitual}

\begin{center}
\begin{tikzpicture}[font=\scriptsize, x=1.45cm, y=1cm]
  \draw[alxBorder, -{Stealth[length=2mm]}] (0,0) -- (9.6,0);
  \fill[alxTeal!35] (0,0.15) rectangle (3,0.75);    \node at (1.5,0.45) {\texttt{Active}};
  \fill[alxAmber!40] (3,0.15) rectangle (7,0.75);   \node at (5,0.45) {\texttt{Degraded}};
  \fill[alxRed!30] (7,0.15) rectangle (8.6,0.75);   \node at (7.8,0.45) {\texttt{FailClosed}};
  \fill[alxTeal!35] (8.6,0.15) rectangle (9.5,0.75);
  \draw (1,0.1) -- (1,-0.1) node[below] {hub cae};
  \draw (3,0.1) -- (3,-0.1) node[below] {\texttt{validUntil}};
  \draw (7,0.1) -- (7,-0.1) node[below] {+ \texttt{max\_staleness}};
  \draw (8.6,0.1) -- (8.6,-0.1) node[below] {listado nuevo};
  \node[above] at (5,0.78) {por gate: \texttt{continue} o \texttt{reject}};
  \node[above] at (7.8,0.78) {todo se rechaza};
\end{tikzpicture}
\end{center}

\section{Cuando el hub responde, pero mal}

Un hub que no contesta es un problema de disponibilidad. Uno que contesta algo incorrecto es un problema de
seguridad, y el cliente lo trata con más dureza: descarta el listado entero y sigue con la instantánea anterior.

\noindent\begin{tabularx}{\linewidth}{@{}L{6.7cm}Y@{}}
\toprule
\textbf{Código} & \textbf{Qué ha visto el cliente} \\
\midrule
\cod{SYNC_SIGNATURE_INVALID} & La firma del listado no verifica contra el DID del hub. \\
\cod{SYNC_ROLLBACK} & Una secuencia menor que la última aceptada: alguien sirve un listado antiguo. \\
\cod{SYNC_SEQUENCE_CONFLICT} & La misma secuencia con otro contenido. \\
\cod{SYNC_REVOCATION_DROPPED} & Una revocación que estaba ha desaparecido. \\
\cod{SYNC_DATES_CHANGED} & Una fecha ya emitida ha cambiado. \\
\cod{SYNC_MANIFEST_DIGEST_MISMATCH}, \cod{SYNC_FILE_HASH_MISMATCH} & Un paquete o un fichero no coincide con su resumen. \\
\bottomrule
\end{tabularx}

\medskip
Estos casos no se han provocado en este tutorial, porque exigen un hub que mienta. La batería de conformidad
del repositorio tiene un caso negativo mínimo para cada uno:

\begin{lstlisting}[style=sh]
go test ./conformance -run 'TestN_' -v
\end{lstlisting}

\section{Mensajes problemáticos}

\noindent\begin{tabularx}{\linewidth}{@{}L{5.5cm}L{5.9cm}C{2.2cm}@{}}
\toprule
\textbf{Mensaje} & \textbf{Código, con la regla en \texttt{enforce}} & \textbf{¿Ejecutado?} \\
\midrule
Política fuera del esquema & \cod{VERDICT_VIOLATION} & \si \\
Contexto JSON-LD que no está en el paquete & \cod{CONTEXT_NOT_IN_PACKAGE} & \si \\
Envoltura DSP incompleta & \cod{VERDICT_VIOLATION} en \cod{dsp_message} & \si \\
Ninguna versión vigente todavía & \cod{NO_VERSION_IN_FORCE} & \si \\
Acuerdo de una versión revocada & \cod{AGREEMENT_VERSION_REVOKED} & \si \\
Listado caducado, gate en \cod{reject} & \cod{CLIENT_DEGRADED_REJECT} & \si \\
Listado caducado más de \cod{max_staleness} & \cod{CLIENT_FAIL_CLOSED} & \si \\
Cuerpo que no es JSON, o sin el nodo de \cod{select} & \cod{MESSAGE_UNPARSEABLE} & no \\
Petición sin regla, con \texttt{default\_action: reject} & \cod{RULE_DEFAULT_REJECT} & no \\
\bottomrule
\end{tabularx}

\section{Dos formas de desplegar el cliente}

El cliente decide sobre cada mensaje. \emph{Quién reenvía} el mensaje es una decisión de despliegue, que se
elige con \cod{mode}.

\begin{center}
\begin{tikzpicture}[font=\small, node distance=7mm and 13mm,
  box/.style={draw=alxNavy, fill=alxLightBlue, rounded corners=2pt, minimum height=9mm, minimum width=29mm, align=center},
  ext/.style={draw=alxBorder, fill=alxSlate, rounded corners=2pt, minimum height=9mm, minimum width=24mm, align=center},
  arr/.style={-{Stealth[length=2.2mm]}, alxNavy, thick}]
  \node[ext] (p1) {Contraparte};
  \node[box, right=of p1] (c1) {\texttt{alexandria-client}};
  \node[ext, right=of c1] (k1) {Conector};
  \draw[arr] (p1) -- (c1); \draw[arr] (c1) -- (k1);
  \node[left=3mm of p1, anchor=east] {\texttt{proxy}};

  \node[ext, below=17mm of p1] (p2) {Contraparte};
  \node[box, right=of p2, fill=alxSlate, draw=alxBorder] (e2) {Envoy};
  \node[ext, right=of e2] (k2) {Conector};
  \node[box, below=8mm of e2] (c2) {\texttt{alexandria-client}};
  \draw[arr] (p2) -- (e2); \draw[arr] (e2) -- (k2);
  \draw[arr] (e2) to[bend right=18] node[left]{\scriptsize\texttt{ext\_authz}} (c2);
  \draw[arr] (c2) to[bend right=18] node[right]{\scriptsize 200 o rechazo} (e2);
  \node[left=3mm of p2, anchor=east] {\texttt{validator}};
\end{tikzpicture}
\end{center}

\noindent\begin{tabularx}{\linewidth}{@{}L{3.3cm}Y Y@{}}
\toprule
 & \textbf{\texttt{mode: proxy}}, por defecto & \textbf{\texttt{mode: validator}} \\
\midrule
Quién reenvía & El cliente. & Un proxy externo, como Envoy. \\
Papel del cliente & Decide y reenvía. & Solo decide; nunca toca el conector. \\
\cod{upstreams} & Obligatorio. & No se usa. \\
\emph{Listeners} de tráfico & Proxies inverso y directo. & Puntos de consulta a los que llama el proxy. \\
\bottomrule
\end{tabularx}

\medskip
El camino de decisión es el mismo código en los dos modos: reglas, gates, clasificación, cabecera, eventos,
métricas y API de administración. En modo \cod{validator}, el cliente responde \cod{200} para permitir y el
rechazo habitual, con cuerpo y cabecera, para denegar. Es el contrato del filtro \cod{ext_authz} de Envoy.

\begin{lstlisting}[style=yaml]
mode: validator
trust:
  hub: "https://127.0.0.1:8443"
  namespace: "demo"
  hub_did: "did:key:z6Mk…"
listeners:
  dsp-ingress: { bind: "127.0.0.1:9085" }   # solo debe alcanzarlo el proxy
  admin: { bind: "127.0.0.1:9090" }
rules:
  seed: "file:deploy/rules.seed.json"
events:
  sink: "file:var/client/events.jsonl"
\end{lstlisting}

Lo que cambia en modo \cod{validator}, porque el cliente ya no ve la respuesta del conector:

\begin{itemize}[leftmargin=*, itemsep=2pt]
  \item La \textbf{captura de acuerdos} ocurre cuando el mensaje se permite, no cuando el conector responde
        \cod{2xx}.
  \item Las reglas con \cod{validate_response} no se evalúan.
  \item El \textbf{reenvío byte a byte} pasa a ser responsabilidad del proxy.
  \item La \textbf{disponibilidad} es un ajuste del proxy: el ejemplo de Envoy cierra el paso si el validador
        no responde.
\end{itemize}

El repositorio trae un ejemplo ejecutable de cada topología, que levanta su propio laboratorio y repite el
escenario F6 comprobando la misma tabla de 29 veredictos:

\begin{lstlisting}[style=sh]
go run ./examples/proxy -bin-dir bin
\end{lstlisting}
\begin{lstlisting}[style=out]
…
2026-11-12T00:00:10Z F6.4 A1 transfer    400 AGREEMENT_VERSION_REVOKED    violation;gate=agreement_admission;code="AGREEMENT_VERSION_REVOKED";rule="transfer-request"
F6 scenario completed: 29/29 rows verified.
\end{lstlisting}

\begin{notaPractica}[Qué no se ha ejecutado]
\texttt{go run ./\allowbreak{}examples/\allowbreak{}envoy} necesita Envoy en el \cod{PATH} o descargar su imagen de Docker, y no se ha
ejecutado para este documento. Lo que se dice aquí del modo \cod{validator} procede de \cod{examples/README.md},
de \cod{examples/envoy/README.md} y del código de configuración.
\end{notaPractica}

\section{Cerrar el laboratorio}

\begin{lstlisting}[style=sh]
for port in 18443 19400 18282 19090; do kill $(lsof -nP -iTCP:$port -sTCP:LISTEN -t) 2>/dev/null; done
rm -rf var clock *.token *.log last-candidate.* client.local.yaml idp.key
\end{lstlisting}

Borrar \cod{var/} elimina la base de datos, las claves y el estado del cliente. El siguiente arranque genera un
DID nuevo.

\begin{retos}
\begin{enumerate}[leftmargin=*, itemsep=2pt]
  \item Cambia \cod{degraded.policy_intake} a \cod{continue}, publícalo como una versión nueva y repite la caída
        del hub. ¿Qué cambia en la segunda fase? ¿Y en la tercera?
  \item Con el hub caído, intenta cambiar una regla del cliente. ¿Depende del hub?
  \item Ejecuta \texttt{go run ./\allowbreak{}examples/\allowbreak{}proxy -keep} y envía \cod{m1-forecast} a su \emph{ingress}. ¿Por qué se
        rechaza con ese laboratorio y no con el tuyo?
\end{enumerate}
\end{retos}


\part{Revisión: la parte de SHACL}
\chapter{Revisión de la parte de SHACL}
\label{ch:shacl}

\section{Qué está hecho y qué no}

\noindent\begin{tabularx}{\linewidth}{@{}L{5.6cm}C{1.5cm}Y@{}}
\toprule
\textbf{Pieza} & \textbf{Estado} & \textbf{Dónde} \\
\midrule
Emisión de \cod{shapes.ttl} por perfil & \si & \cod{internal/dsl/emit_shacl.go} \\
Emisión de \cod{space/shacl-union.ttl} & \si & \cod{internal/dsl/compile.go:336} \\
Comprobación de las \emph{shapes} al compilar & parcial & Solo estructura: \cod{internal/dsl/turtle_check.go} \\
Ejecución de SHACL en el cliente & \no & \cod{evaluator_poc/evaluator.go:88}: expande el JSON-LD y anota un aviso \\
Avisos de severidad \cod{sh:Warning} en el veredicto & \no & \cod{VERDICT_SHACL_WARNING} es inalcanzable; \cod{ShaclWarnings} (\cod{classify.go:43}) nunca se activa \\
\bottomrule
\end{tabularx}

\medskip
En resumen: las \emph{shapes} se generan y se publican firmadas, pero \textbf{nadie las ha ejecutado nunca}, ni
el compilador ni el cliente. Esta revisión es la primera vez que pasan por un motor SHACL.

\section{Método}

Se han tomado 27 ofertas ODRL compactas y cada una se ha validado dos veces: contra \cod{schema.json} con
\cod{jsonschema}, y contra \cod{shapes.ttl} con \cod{pyshacl}~0.40.2 tras expandir el JSON-LD con el contexto
generado. El banco de pruebas es \cod{examples/shacl/harness.py}.

\begin{lstlisting}[style=sh]
python3 -m venv /tmp/venv && /tmp/venv/bin/pip install pyshacl jsonschema
/tmp/venv/bin/python shacl/harness.py /tmp/metering/profiles/odrl/metering/1.0.0 -v
\end{lstlisting}

El contexto de ODRL que usa el banco es una copia mínima escrita a mano, para no depender de la red.

\section{Resultados}

Primera conclusión, y es buena: las \emph{shapes} generadas son SHACL válido, el motor las carga sin errores y
\textbf{coinciden con el esquema en 21 de 27 casos}. El diseño de \cod{emit_shacl.go} es correcto en lo esencial.

\begin{longtable}{@{}C{0.7cm}L{7.6cm}C{1.6cm}C{2.5cm}C{1.4cm}@{}}
\toprule
\textbf{N.º} & \textbf{Oferta} & \textbf{Esquema} & \textbf{SHACL} & \textbf{¿Igual?} \\
\midrule
\endhead
01 & Oferta válida & \pasa & \pasa & \si \\
02 & Acción ajena (\cod{sell}) & \falla & \falla, más aviso & \si \\
03 & Operador fuera de la firma & \falla & \falla & \si \\
04 & Literal de tipo equivocado & \falla & \falla & \si \\
05 & Finalidad fuera de la lista & \falla & \falla & \si \\
06 & \cod{isAnyOf} con dos regiones & \pasa & \pasa & \si \\
07 & Acción del núcleo (\cod{use}) & \pasa & \pasa, con aviso & \si \\
08 & \textbf{Perfil no declarado} & \falla & \pasa & \no \\
09 & Restricción lógica \cod{and} & \pasa & \pasa & \si \\
10 & Restricción lógica con hoja inválida & \falla & \falla & \si \\
11 & Fecha como cadena sin tipar & \falla & \falla & \si \\
12 & Fecha tipada & \pasa & \pasa & \si \\
13 & Entero como valor tipado & \pasa & \pasa & \si \\
14 & Decimal donde va un entero & \falla & \falla & \si \\
15 & Referencia no permitida & \falla & \falla & \si \\
16 & Operando izquierdo ajeno & \falla & \falla & \si \\
17 & Dos valores sin operador de conjunto & \falla & \falla & \si \\
18 & \cod{eq} con dos finalidades & \pasa & \pasa & \si \\
19 & Finalidad como cadena, no \cod{@id} & \falla & \falla & \si \\
20 & Deber con acción ajena & \falla & \falla, más aviso & \si \\
21 & Acción como objeto con \cod{@id} & \pasa & \pasa & \si \\
22 & \textbf{Fecha tipada con forma léxica inválida} & \pasa & \falla & \no \\
23 & Prohibición de un término obsoleto (1.1.0) & \pasa & \pasa, con aviso & \si \\
24 & \textbf{Regla sin acción} & \pasa & \falla & \no \\
25 & \textbf{Entero tipado con forma léxica inválida} & \pasa & \falla & \no \\
26 & Permiso con término obsoleto (1.1.0) & \pasa & \pasa, con aviso & \si \\
27 & Finalidad \cod{research} bajo la 1.0.0 & \falla & \falla & \si \\
\bottomrule
\end{longtable}

A estas se suma un caso del perfil de salud: una restricción sobre \cod{dataPermit} \textbf{sin operando derecho
ni referencia} pasa en SHACL y falla en el esquema.

\section{Hallazgos}

Ordenados por gravedad. H1 y H2 son la carencia; H3 a H9 son defectos que solo se ven al ejecutar.

\begin{hallazgo}{H1 (alta). Un slot obligatorio validado solo con SHACL no valida nada}
En el cliente, un artefacto de rol \cod{shacl} produce una expansión JSON-LD y un hallazgo \cod{Info}
(\cod{evaluator.go:88-90}). Si un \emph{binding} \cod{mandatory} enlaza \emph{únicamente}
\cod{shacl} a un gate, ese gate acepta cualquier política que use contextos del paquete. Lo confirma el propio
test del repositorio \cod{TestShaclIsReportedAsUnsupportedInfo}, que espera \cod{Pass} verdadero.

Es exactamente lo que hace \cod{internal/dsl/testdata/demo.space.yaml}: \cod{shacl} en \cod{policy_intake}.
Compila, se publica como obligatorio y no controla nada. \emph{Propuesta inmediata}, sin motor SHACL: que el
compilador avise cuando un \emph{binding} obligatorio no tenga ningún \cod{json-schema} en alguno de sus gates.
\end{hallazgo}

\begin{hallazgo}{H2 (alta). N5 no comprueba que las shapes sean SHACL}
\cod{checkTurtle} cuenta corchetes, comillas e IRI. No analiza Turtle (falla con cadenas de triple comilla, que
\cod{extra_turtle} puede traer) y no carga las \emph{shapes}. Un error en un emisor se publicaría firmado.
\end{hallazgo}

\begin{hallazgo}{H3 (media). Los dos artefactos no dicen lo mismo en cinco casos}
\begin{itemize}[leftmargin=*, itemsep=1pt]
  \item \textbf{Caso 08.} Ninguna \emph{shape} exige \cod{odrl:profile}. El esquema sí. Una política que no
        declara el perfil pasaría el gate SHACL.
  \item \textbf{Caso 24.} \cod{gen:ActionShape} exige \texttt{sh:minCount 1} en \cod{odrl:action}; el esquema no
        marca \cod{action} como obligatorio en \cod{Rule} (\cod{emit_json.go}).
  \item \textbf{Casos 22 y 25.} El esquema acepta cualquier cadena en \cod{@value}; SHACL comprueba la forma
        léxica del tipo. Aquí SHACL es el estricto.
  \item \textbf{Referencias.} Con \texttt{allow\_reference: true}, la \emph{shape} no exige «al menos uno de
        \cod{rightOperand} o \cod{rightOperandReference}»; el esquema sí.
\end{itemize}
Mientras solo se ejecute el esquema no tiene efecto. El día que se ejecuten los dos, el mismo mensaje tendrá dos
veredictos.
\end{hallazgo}

\begin{hallazgo}{H4 (media). \texttt{sh:conforms} es falso con solo avisos, y hay avisos espurios}
En SHACL, \cod{sh:conforms} es falso si hay \emph{cualquier} resultado, también de severidad \cod{sh:Warning}
(casos 07, 23 y 26). Quien implemente el gate tiene que clasificar por \cod{sh:resultSeverity}, no por
\cod{sh:conforms}.

Además, \cod{gen:ActionCoreShape} (\cod{emit_shacl.go:26}) avisa «las acciones del núcleo se ignoran» también
para una acción que no es del núcleo ni del perfil (casos 02 y 20): la misma acción produce una violación y un
aviso con un mensaje que no le corresponde.
\end{hallazgo}

\begin{hallazgo}{H5 (baja). Los mensajes de tres shapes se pierden}
En \cod{gen:ActionShape}, \cod{gen:ActionCoreShape} y \cod{gen:LeftOperandShape}, \cod{sh:message} está en el
\cod{NodeShape}, pero la restricción está en una \emph{property shape} anónima anidada. El resultado sale con el
mensaje genérico del motor («Value odrl:sell not in list…») y con un nodo anónimo como \cod{sh:sourceShape}: no
se puede saber qué \emph{shape} con nombre falló. Las \cod{gen:Sig_*} y la de obsolescencia sí conservan su
mensaje. \emph{Arreglo}: mover \cod{sh:message} dentro de la \emph{property shape}.
\end{hallazgo}

\begin{hallazgo}{H6 (media). \texttt{shacl-union.ttl} no se puede consumir por gate}
El gate solo figura en comentarios (\texttt{\# gate policy\_intake · slot odrl · …}), que un analizador RDF
descarta. Un artefacto enlazado a dos gates se concatena dos veces. Y \cod{SHACL_UNION_CONFLICT}
(\cod{compile.go:355}) es inalcanzable: solo mira el slot \cod{odrl}, y \cod{SLOT_DUP} garantiza un único
\emph{binding} por slot. Hoy el fichero equivale a \cod{shapes.ttl} (223 tripletas ambos) y ningún componente lo
lee.
\end{hallazgo}

\begin{hallazgo}{H7 (media). \texttt{scheme} genera una shape que nunca puede cumplirse}
\cod{emit_shacl.go:98} emite \texttt{sh:node [ sh:property [ sh:path skos:inScheme ; sh:hasValue … ] ]}. Pero
ningún campo del lenguaje declara la pertenencia a un esquema, \cod{vocab.ttl} no emite \cod{skos:inScheme}, y
el grafo de datos de una política no contiene el vocabulario. El compilador ya avisa
(\cod{SIG_SCHEME_UNSUPPORTED}); convendría que fuera un error hasta que exista el campo.
\end{hallazgo}

\begin{hallazgo}{H8 (baja). El aviso del PoC sale del sistema}
\cod{POC_EVAL_UNSUPPORTED_ROLE} aparece en el \cod{reason} del error DSP que recibe la contraparte
(\cod{decide.go:148}, que copia todos los hallazgos sin filtrar por severidad), y el evento \cod{vhub.sync} que
lo anuncia lleva \texttt{"severity":"error"} aunque el hallazgo es \cod{Info}.
\end{hallazgo}

\begin{hallazgo}{H9 (baja). Detalles}
\begin{itemize}[leftmargin=*, itemsep=1pt]
  \item Una firma con un operador de conjunto relaja la cardinalidad para \emph{todos} sus operadores: \cod{eq}
        con dos valores pasa (caso 18), en los dos artefactos.
  \item \cod{profile/slots.json} repite el gate cuando dos roles lo comparten:
        \texttt{["policy\_intake","policy\_intake"]} (\cod{compiler_dsl/compiler.go:97}).
  \item Los objetivos de las \emph{shapes} no cubren \cod{odrl:consequence} ni \cod{odrl:remedy}.
\end{itemize}
\end{hallazgo}

\section{Qué haría falta para cerrarlo}

\begin{enumerate}[leftmargin=*, itemsep=3pt]
  \item \textbf{Sin motor, ya.} Corregir H5 y H4 en el emisor, igualar esquema y \emph{shapes} (H3), filtrar
        H8, y añadir el aviso de H1. Son cambios pequeños en \cod{emit_shacl.go}, \cod{emit_json.go},
        \cod{decide.go} y \cod{space.go}.
  \item \textbf{Una prueba de equivalencia en CI.} El banco de este capítulo, convertido en test: cada fila
        declara el resultado esperado en ambos artefactos. Puede vivir fuera del módulo Go, con pySHACL, sin
        tocar las dependencias permitidas.
  \item \textbf{Elegir motor.} Las opciones que recoge \cod{PROFILES_DSL.md} siguen siendo las reales:
        goRDFlib en Go, reciente y de un solo autor; o un proceso auxiliar con pySHACL o Jena detrás del puerto
        \cod{app.Evaluator}. La arquitectura hexagonal ya lo permite: es un adaptador nuevo junto a
        \cod{evaluator_poc}.
  \item \textbf{Cablear severidades.} Violación de un slot obligatorio: violación. \cod{sh:Warning}:
        \cod{ShaclWarnings} y \cod{VERDICT_SHACL_WARNING}. Con eso el aviso de obsolescencia del
        capítulo~\ref{ch:evolucion} llegaría por fin a quien usa el término.
  \item \textbf{N5 semántico.} Cargar las \emph{shapes} generadas con el mismo motor al compilar.
\end{enumerate}


\appendix
\chapter{Chuleta}
\label{ch:chuleta}

\section{El laboratorio completo, de un tirón}

Desde \cod{tutorial/manual/examples/lab}, con los binarios compilados. Es la secuencia exacta que se ha
ejecutado para este documento.

\begin{lstlisting}[style=sh]
source env.sh
clock 2026-09-01T09:00:00Z
$BIN/fake-idp serve --addr 127.0.0.1:19400 --issuer $IDP --key idp.key --clock-file clock > idp.log 2>&1 &
$BIN/fake-connector -listen 127.0.0.1:18282 > connector.log 2>&1 &
$BIN/alexandria-server -config server.yaml -clock-file clock > server.log 2>&1 &
sleep 3

# Gobierno
tok --sub ana --admin > admin.token
tok --sub lucia > steward.token
curl -sk -X POST $HUB/v1/namespaces -H "Authorization: Bearer $(cat admin.token)" \
  -d '{"slug":"flexgrid","title":"Flexgrid Energy Data Space"}'
curl -sk -X POST $HUB/v1/namespaces/flexgrid/stewards -H "Authorization: Bearer $(cat admin.token)" \
  -d '{"iss":"http://127.0.0.1:19400","sub":"lucia","label":"Lucia (DSO)"}'

# 1.0.0
clock 2026-09-01T10:00:00Z
publish ../dsl/flexgrid-1.0.0.yaml && approve '{}'

# Cliente
sed "s/HUB_DID_HERE/$(hubdid)/" client.yaml > client.local.yaml
$BIN/alexandria-client -config client.local.yaml -state-dir var/client/state -clock-file clock > client.log 2>&1 &
sleep 3
setrule contract-request; setrule contract-agreement-in; setrule transfer-request
syncat 2026-09-03T12:00:00Z && rows

# 1.1.0, compatible
syncat 2026-09-20T10:00:00Z
publish ../dsl/flexgrid-1.1.0.yaml && approve '{}'
syncat 2026-09-23T12:00:00Z && rows

# 2.0.0, ruptura
syncat 2026-10-01T10:00:00Z
publish ../dsl/flexgrid-2.0.0.yaml && approve '{}'
syncat 2026-10-10T12:00:00Z && rows
post /negotiations/urn:uuid:7d2a9f10-ret-neg-0007/agreement a1-agreement
syncat 2026-11-05T12:00:00Z && rows
post /transfers/request t1-transfer

# Revocación
syncat 2026-11-10T00:00:00Z
curl -sk -X POST "$HUB/v1/namespaces/flexgrid/versions/1.1.0:revoke" \
  -H "Authorization: Bearer $(cat steward.token)" \
  -d '{"justification":"Individual consumer profiling was ruled unlawful by the regulator"}'
sleep 2; clock 2026-11-12T00:00:10Z; sleep 2
post /transfers/request t1-transfer
\end{lstlisting}

\section{\texttt{alexandriac}}

\noindent\begin{tabularx}{\linewidth}{@{}L{9cm}Y@{}}
\toprule
\textbf{Comando} & \textbf{Para qué} \\
\midrule
\texttt{alexandriac validate f.yaml} & Validar, sin escribir nada. \\
\texttt{alexandriac compile -o out f.yaml} & Validar y generar los artefactos. \\
\texttt{… -vendor vendor/\allowbreak{} bundle.yaml} & Resolver \emph{bindings} vendorizados. \\
\texttt{… -previous v1.yaml v2.yaml} & Comprobar el versionado contra la versión anterior. \\
\bottomrule
\end{tabularx}

\section{API del hub}

\begin{longtable}{@{}L{8.8cm}L{6.6cm}@{}}
\toprule
\textbf{Operación} & \textbf{Quién} \\
\midrule
\endhead
\texttt{POST /\allowbreak{}v1/\allowbreak{}namespaces} & Admin \\
\texttt{POST /\allowbreak{}v1/\allowbreak{}namespaces/\allowbreak{}\{ns\}/\allowbreak{}stewards} & Admin \\
\texttt{DELETE /\allowbreak{}v1/\allowbreak{}namespaces/\allowbreak{}\{ns\}/\allowbreak{}stewards/\allowbreak{}\{id\}} & Admin \\
\texttt{GET /\allowbreak{}v1/\allowbreak{}events} & Admin \\
\texttt{POST /\allowbreak{}v1/\allowbreak{}namespaces/\allowbreak{}\{ns\}/\allowbreak{}vocabularies} & Steward \\
\texttt{PUT /\allowbreak{}v1/\allowbreak{}namespaces/\allowbreak{}\{ns\}/\allowbreak{}vocabularies/\allowbreak{}\{nombre\}/\allowbreak{}content} & Steward \\
\texttt{POST /\allowbreak{}v1/\allowbreak{}namespaces/\allowbreak{}\{ns\}/\allowbreak{}vocabularies/\allowbreak{}\{nombre\}:refresh} & Steward \\
\texttt{POST /\allowbreak{}v1/\allowbreak{}namespaces/\allowbreak{}\{ns\}/\allowbreak{}drafts} & Steward \\
\texttt{PUT /\allowbreak{}v1/\allowbreak{}namespaces/\allowbreak{}\{ns\}/\allowbreak{}drafts/\allowbreak{}\{id\}}, con \cod{If-Match} & Steward \\
\texttt{POST /\allowbreak{}v1/\allowbreak{}namespaces/\allowbreak{}\{ns\}/\allowbreak{}drafts/\allowbreak{}\{id\}:validate} & Steward \\
\texttt{POST /\allowbreak{}v1/\allowbreak{}namespaces/\allowbreak{}\{ns\}/\allowbreak{}drafts/\allowbreak{}\{id\}:propose} & Steward \\
\texttt{POST /\allowbreak{}v1/\allowbreak{}namespaces/\allowbreak{}\{ns\}/\allowbreak{}candidates/\allowbreak{}\{id\}:approve} & Steward \\
\texttt{POST /\allowbreak{}v1/\allowbreak{}namespaces/\allowbreak{}\{ns\}/\allowbreak{}candidates/\allowbreak{}\{id\}:reject} & Steward \\
\texttt{POST /\allowbreak{}v1/\allowbreak{}namespaces/\allowbreak{}\{ns\}/\allowbreak{}versions/\allowbreak{}\{v\}:revoke} & Steward \\
\texttt{GET /\allowbreak{}v1/\allowbreak{}namespaces/\allowbreak{}\{ns\}/\allowbreak{}events} & Steward \\
\texttt{GET /\allowbreak{}v1/\allowbreak{}namespaces/\allowbreak{}\{ns\}/\allowbreak{}versions} & Cualquiera: el listado firmado \\
\texttt{GET /\allowbreak{}v1/\allowbreak{}namespaces/\allowbreak{}\{ns\}/\allowbreak{}versions/\allowbreak{}\{v\}} & Cualquiera: el registro firmado \\
\texttt{GET /\allowbreak{}pub/\allowbreak{}…} & Cualquiera \\
\bottomrule
\end{longtable}

\section{API de administración del cliente}

\noindent\begin{tabularx}{\linewidth}{@{}L{5.6cm}Y@{}}
\toprule
\textbf{Ruta} & \textbf{Para qué} \\
\midrule
\texttt{GET /\allowbreak{}status} & Estado, versiones locales con su papel, último listado, DID del hub y revisión de reglas. \\
\texttt{GET /\allowbreak{}versions} & Registros conocidos con su estado y si se conserva el paquete. \\
\texttt{GET /\allowbreak{}rules} & El conjunto de reglas; su revisión es el \cod{ETag}. \\
\texttt{PUT /\allowbreak{}rules}, \texttt{POST /\allowbreak{}rules} & Reemplazar el conjunto o añadir una regla. Con \cod{If-Match}. \\
\cod{GET}, \cod{PUT}, \texttt{DELETE /\allowbreak{}rules/\allowbreak{}\{id\}} & Leer, reemplazar o borrar una regla. \\
\texttt{POST /\allowbreak{}rules:explain} & Qué haría el proxy con una petición, sin efectos. \\
\texttt{GET /\allowbreak{}agreements/\allowbreak{}\{id\}} & Un acuerdo capturado y si es utilizable. \\
\texttt{GET /\allowbreak{}audit/\allowbreak{}reports/\allowbreak{}latest} & Último informe de auditoría pasiva. \\
\texttt{GET /\allowbreak{}healthz}, \texttt{GET /\allowbreak{}readyz} & Vida y disponibilidad. \\
\texttt{GET /\allowbreak{}metrics} & OpenMetrics. \\
\bottomrule
\end{tabularx}

\medskip
La API de administración no tiene autenticación: su \emph{listener} no debe ser alcanzable desde fuera.

\section{Códigos vistos en este tutorial}

\begin{longtable}{@{}L{6.9cm}L{2.5cm}L{6cm}@{}}
\toprule
\textbf{Código} & \textbf{Quién} & \textbf{Cuándo} \\
\midrule
\endhead
\cod{API_UNAUTHENTICATED} & Hub, 401 & Operación protegida sin token. \\
\cod{API_FORBIDDEN_ROLE} & Hub, 403 & Token válido sin el papel necesario. \\
\cod{VOCAB_UNPARSEABLE} & Hub, 422 & Contenido que no es del tipo declarado. \\
\cod{VERSION_BUMP_TOO_LOW} & Compilador & El salto de versión no cubre los cambios. \\
\cod{TERM_RETIRED_WITHOUT_DEPRECATION} & Compilador & Retirar sin haber marcado obsoleto antes. \\
\cod{TERM_REMOVED} & Compilador & Un término ha desaparecido. \\
\cod{PROFILE_VERSION_CHANGED} & Compilador & Misma versión, otro contenido. \\
\cod{GATE_SHACL_ONLY} & Compilador, aviso & Gate obligatorio cubierto solo por \cod{shacl}. \\
\cod{BINDING_UNRESOLVED} & Compilador & Rango sin versión, o falta la copia vendorizada. \\
\cod{DATES_GRACE_ON_COMPATIBLE} & Hub, 422 & Gracia en una versión compatible. \\
\cod{DATES_GRACE_TOO_SHORT} & Hub, 422 & Gracia por debajo de la mínima. \\
\cod{DATES_JUSTIFICATION_REQUIRED} & Hub, 422 & Revocación sin justificación. \\
\cod{DATES_REVOKE_LEAVES_NO_HEAD} & Hub, 422 & Revocar la única versión vigente. \\
\cod{NO_VERSION_IN_FORCE} & Cliente & Aún no hay versión vigente. \\
\cod{VERDICT_VIOLATION} & Cliente & El mensaje no conforma con ninguna versión vigente. \\
\cod{VERDICT_NOT_IN_HEAD} & Cliente & Conforma con una versión anterior aún vigente. \\
\cod{CONTEXT_NOT_IN_PACKAGE} & Cliente & Contexto JSON-LD ajeno al paquete. \\
\cod{AGREEMENT_REVOCATION_SCHEDULED} & Cliente, evento & La versión de un acuerdo va a revocarse. \\
\cod{AGREEMENT_VERSION_REVOKED} & Cliente & Transferencia bajo un acuerdo de versión revocada. \\
\cod{SYNC_UNREACHABLE} & Cliente, evento & El hub no responde. \\
\cod{CLIENT_DEGRADED_REJECT} & Cliente & Listado caducado y gate en \cod{reject}. \\
\cod{CLIENT_FAIL_CLOSED} & Cliente & Listado caducado más de \cod{max_staleness}. \\
\cod{POC_EVAL_UNSUPPORTED_ROLE} & Cliente, aviso & Artefacto \cod{shacl} que no se ejecuta. \\
\bottomrule
\end{longtable}

\section{Lo que conviene recordar}

\begin{enumerate}[leftmargin=*, itemsep=3pt]
  \item \textbf{Un \emph{binding} obligatorio necesita un \cod{json-schema} en cada gate.} Con solo \cod{shacl}
        compila, se publica y no valida nada; el compilador avisa con \cod{GATE_SHACL_ONLY}.
  \item \textbf{Una regla valida un nodo.} Para la política de una petición de contrato, \texttt{select: /\allowbreak{}offer}; y
        entonces esa regla no puede validar la envoltura.
  \item \textbf{El artefacto \cod{shacl} activa la comprobación de contextos.} Si lo enlazas, todo contexto que
        citen los mensajes tiene que estar en el paquete con su \cod{url}, incluido el de DSP.
  \item \textbf{Nada se borra y nada publicado cambia.} Obsoleto, luego retirado; y una errata es una versión
        de parche.
  \item \textbf{Retirar no rompe acuerdos; revocar sí}, con dos días de aviso.
  \item \textbf{Con reloj de fichero, todo espera al reloj}: los tokens, la reemisión del listado y los
        reintentos.
\end{enumerate}

\section{Recompilar este documento}

Las fuentes están en \cod{tutorial/manual/}: \cod{main.tex} y un fichero por capítulo en \cod{chapters/}.

\begin{lstlisting}[style=sh]
cd tutorial/manual
make            # con latexmk, tectonic o pdflatex en el PATH
make docker     # con la imagen texlive/texlive:latest-small
\end{lstlisting}

La imagen \cod{latest-small} no trae \cod{tcolorbox}, \cod{titlesec} ni \cod{xurl}; el objetivo \cod{docker}
los instala con \cod{tlmgr} antes de compilar.


\end{document}
